\documentclass[11pt,openany]{book}

\usepackage[T1]{fontenc}
\usepackage[utf8]{inputenc}
\usepackage{lmodern}
\usepackage{microtype}
\usepackage[margin=1in,headheight=15pt]{geometry}
\usepackage{amsmath,amssymb,amsthm,mathtools,mathrsfs}
\usepackage{booktabs,longtable,tabularx,array,multirow}
\usepackage{enumitem}
\usepackage{graphicx}
\usepackage{xcolor}
\usepackage{tikz}
\usetikzlibrary{arrows.meta,positioning,calc,decorations.pathreplacing,fit,backgrounds}
\usepackage[most]{tcolorbox}
\usepackage{fancyhdr}
\usepackage{titlesec}
\usepackage{hyperref}
\usepackage[capitalise,noabbrev]{cleveref}
\usepackage{bookmark}
\usepackage{url}
\usepackage{needspace}

\definecolor{DeepBlue}{HTML}{17365D}
\definecolor{SkyBlue}{HTML}{DCEAF7}
\definecolor{DeepGreen}{HTML}{245B45}
\definecolor{PaleGreen}{HTML}{E4F2EA}
\definecolor{DeepRed}{HTML}{8B2E2E}
\definecolor{PaleRed}{HTML}{F8E6E4}
\definecolor{DeepGold}{HTML}{7A5A00}
\definecolor{PaleGold}{HTML}{FFF4CC}
\definecolor{SoftGray}{HTML}{F3F4F6}

\hypersetup{
  colorlinks=true,
  linkcolor=DeepBlue,
  citecolor=DeepGreen,
  urlcolor=DeepRed,
  pdftitle={Transseries: A Step-by-Step Tutorial},
  pdfsubject={Asymptotic expansions, generalized series, transseries, resurgence, and applications},
  pdfkeywords={transseries, asymptotics, Hahn series, Borel summation, resurgence, Stokes phenomenon, H-fields, surreal numbers}
}

\setlist[itemize]{leftmargin=2em,itemsep=0.35em,topsep=0.5em}
\setlist[enumerate]{leftmargin=2.2em,itemsep=0.45em,topsep=0.5em}
\setlength{\parindent}{1.25em}
\setlength{\parskip}{0.2em}
\allowdisplaybreaks
\setlength{\emergencystretch}{2em}

\pagestyle{fancy}
\fancyhf{}
\fancyhead[LE]{\small\nouppercase{\leftmark}}
\fancyhead[RO]{\small\nouppercase{\rightmark}}
\fancyfoot[C]{\thepage}
\renewcommand{\headrulewidth}{0.4pt}

\titleformat{\chapter}[display]
  {\normalfont\bfseries\color{DeepBlue}}
  {\filleft\Large\chaptername\ \thechapter}
  {0.7ex}
  {\titlerule\vspace{1.2ex}\huge}
  [\vspace{0.8ex}\titlerule]
\titleformat{\section}{\Large\bfseries\color{DeepBlue}}{\thesection}{0.7em}{}
\titleformat{\subsection}{\large\bfseries\color{DeepGreen}}{\thesubsection}{0.7em}{}

\theoremstyle{plain}
\newtheorem{theorem}{Theorem}[chapter]
\newtheorem{proposition}[theorem]{Proposition}
\newtheorem{lemma}[theorem]{Lemma}
\newtheorem{corollary}[theorem]{Corollary}
\theoremstyle{definition}
\newtheorem{definition}[theorem]{Definition}
\newtheorem{example}[theorem]{Example}
\newtheorem{exercise}{Exercise}[chapter]
\theoremstyle{remark}
\newtheorem{remark}[theorem]{Remark}

\tcbset{enhanced,breakable,boxrule=0.8pt,arc=1.5mm,left=2mm,right=2mm,top=1mm,bottom=1mm}
\newtcolorbox{ideabox}[1][]{colback=SkyBlue,colframe=DeepBlue,title=Core idea,#1}
\newtcolorbox{intuitionbox}[1][]{colback=PaleGreen,colframe=DeepGreen,title=Intuition,#1}
\newtcolorbox{warningbox}[1][]{colback=PaleRed,colframe=DeepRed,title=Common pitfall,#1}
\newtcolorbox{workedbox}[1]{colback=SoftGray,colframe=DeepBlue,title={Worked example: #1}}
\newtcolorbox{advancedbox}[1][]{colback=PaleGold,colframe=DeepGold,title=Optional advanced viewpoint,#1}
\newtcolorbox{checkpointbox}[1][]{colback=white,colframe=DeepGreen,title=Checkpoint,#1}
\newtcolorbox{sourcebox}[1][]{colback=SoftGray,colframe=black!55,title=Sources and scope,#1}

\newcommand{\R}{\mathbb{R}}
\newcommand{\C}{\mathbb{C}}
\newcommand{\Q}{\mathbb{Q}}
\newcommand{\Z}{\mathbb{Z}}
\newcommand{\N}{\mathbb{N}}
\newcommand{\Tfield}{\mathbb{T}}
\newcommand{\G}{\mathfrak{G}}
\newcommand{\m}{\mathfrak{m}}
\newcommand{\n}{\mathfrak{n}}
\newcommand{\g}{\mathfrak{g}}
\newcommand{\supp}{\operatorname{supp}}
\newcommand{\magop}{\operatorname{mag}}
\newcommand{\domop}{\operatorname{dom}}
\newcommand{\ord}{\operatorname{ord}}
\newcommand{\Ei}{\operatorname{Ei}}
\newcommand{\Borel}{\mathcal{B}}
\newcommand{\Laplace}{\mathcal{L}}
\newcommand{\Stokes}{\mathfrak{S}}
\newcommand{\smallo}{\mathrm{o}}
\newcommand{\bigO}{\mathrm{O}}
\newcommand{\dd}{\,\mathrm{d}}
\newcommand{\e}{\mathrm{e}}
\newcommand{\logit}{\operatorname{logit}}
\newcommand{\abs}[1]{\left|#1\right|}
\newcommand{\set}[1]{\left\{#1\right\}}
\newcommand{\brak}[1]{\left(#1\right)}
\newcommand{\angles}[1]{\left\langle#1\right\rangle}
\newcommand{\asympseries}{\sim_{\mathrm{as}}}
\newcommand{\defterm}[1]{\textbf{#1}}
\newcommand{\llbracket}{\mathopen{[\![}}
\newcommand{\rrbracket}{\mathclose{]\!]}}

\begin{document}

\frontmatter

\begin{titlepage}
\centering
\vspace*{1.4cm}
{\Huge\bfseries\color{DeepBlue} Transseries\par}
\vspace{0.5cm}
{\LARGE A Step-by-Step Tutorial for Mere Mortals\par}
\vspace{0.8cm}
{\Large From asymptotic scales and Hahn series to Borel summation, resurgence, and surreal numbers\par}
\vfill
\begin{tikzpicture}[scale=1.05]
  \node[font=\Large] (a) at (0,0) {$1$};
  \node[font=\Large] (b) at (1.6,0) {$\log x$};
  \node[font=\Large] (c) at (3.5,0) {$x^a$};
  \node[font=\Large] (d) at (5.5,0) {$\e^{\sqrt{x}}$};
  \node[font=\Large] (e1) at (7.8,0) {$\e^x$};
  \node[font=\Large] (f) at (10.0,0) {$\e^{\e^x}$};
  \draw[-{Stealth[length=3mm]},very thick,DeepBlue] (-0.7,-0.8) -- (10.8,-0.8)
      node[midway,below=2mm,font=\large] {increasing orders of growth};
  \foreach \u/\v in {a/b,b/c,c/d,d/e1,e1/f}{\draw[DeepGreen,thick] (\u) -- (\v);}
  \node[draw=DeepRed,rounded corners,fill=PaleRed,align=center] at (5.1,-2.2)
    {$x^{-1}+2x^{-2}+6x^{-3}+\cdots$\\meets\\$C\e^{-x}x^2(1-4/x+\cdots)$};
\end{tikzpicture}
\vfill
{\large Prepared as a self-contained mathematical tutorial\par}
\vspace{0.25cm}
{\large August 2026\par}
\end{titlepage}

\chapter*{Preface}
\addcontentsline{toc}{chapter}{Preface}

Transseries are what happens when ordinary asymptotic series learn to speak the languages of exponentials, logarithms, infinitesimals, infinite quantities, and ``beyond-all-orders'' effects at the same time. A typical transseries may contain powers of $x$, iterated logarithms, exponentials such as $\e^{-x}$, and whole asymptotic series multiplying those exponentials:
\[
  F(x)=\sum_{n\ge 0}\frac{a_n}{x^n}
  +C\e^{-x}x^\beta\sum_{n\ge 0}\frac{b_n}{x^n}
  +C^2\e^{-2x}x^{2\beta}\sum_{n\ge 0}\frac{c_n}{x^n}+\cdots.
\]
This one expression can record ordinary perturbation theory, exponentially small corrections, and the way those corrections interact.

The subject has several overlapping traditions. In algebra and model theory, one studies formal logarithmic-exponential transseries as a large ordered differential field. In asymptotic analysis and mathematical physics, one often studies resurgent transseries whose sectors are connected by Borel-plane singularities and Stokes jumps. In the theory of surreal numbers, transseries appear as canonical infinite expressions evaluated at positive infinite numbers. These traditions share a common grammar but emphasize different questions.

This tutorial starts at the level of a strong second-year undergraduate who has seen calculus, elementary differential equations, and power series. Every major construction is first motivated informally, then stated carefully. Proofs are included when they illuminate the mechanism; technically deep closure theorems are clearly labeled and explained rather than buried under specialist machinery.

\begin{sourcebox}
The attached archive supplied with the request contains, among other materials, G.~A. Edgar's expository source \emph{Transseries for Beginners}, his papers on composition and fractional iteration, Dahn--G\"oring's notes on exponential-logarithmic terms, work of Kuhlmann--Tressl on LE and EL series, papers of Berarducci--Mantova on surreal numbers and transseries, and an exposition by Aschenbrenner--van den Dries--van der Hoeven on numbers, germs, and transseries. It also contains primers and lecture notes on resurgent transseries and Stokes phenomena. Those sources were supplemented by the original arXiv records, the linked Wikipedia overview as an initial map of terminology, and standard monographs listed in the bibliography \cite{WikipediaTransseries}. The prose, organization, examples, and pedagogical development here are new; this is not a reproduction of any one source.
\end{sourcebox}

\section*{How to read this book}
The shortest useful path is:
\[
\text{Chapters 1--4}\longrightarrow
\text{Chapters 5--10}\longrightarrow
\text{Chapters 11--16}\longrightarrow
\text{any application chapters you like}.
\]
Chapters marked by an ``Optional advanced viewpoint'' box may be skipped on a first reading. Part~V, on Borel summation and resurgence, uses complex variables but begins with concrete real integrals. Part~VI explains the wider algebraic landscape without assuming model theory or surreal-number theory.

\section*{A promise about notation}
The variable $x$ tends to $+\infty$ unless stated otherwise. The relation $f\prec g$ means that $f/g\to0$ in an analytic setting, or that the dominant monomial of $f$ is strictly smaller than that of $g$ in a formal setting. We write $f\sim g$ when $f/g\to1$, or formally when $f-g\prec g$. The symbol $\asymp$ means ``same order of magnitude'' and is weaker than $\sim$.

\tableofcontents

\mainmatter

\part{Why ordinary asymptotics asks for transseries}

\chapter{Asymptotics is about scale, not convergence}

\section{The question asked by asymptotic analysis}
Suppose a function is difficult to evaluate exactly, but we care about it for very large $x$. A convergent power series is one possible answer, but it is often the wrong kind of answer. Asymptotic analysis asks a more local and more flexible question:

\begin{quote}
Which simpler terms describe the function in decreasing order of importance as $x\to+\infty$?
\end{quote}

The phrase ``decreasing order of importance'' is the heart of the subject. The sequence
\[
  1,\quad x^{-1},\quad x^{-2},\quad x^{-3},\ldots
\]
is an \defterm{asymptotic scale} because each term is negligible compared with the preceding one:
\[
  x^{-(n+1)}=\smallo(x^{-n}).
\]
More generally, a sequence $(\phi_n)$ is an asymptotic scale if
\[
  \phi_{n+1}(x)=\smallo(\phi_n(x))
  \qquad(x\to+\infty).
\]

\begin{definition}[Poincar\'e asymptotic expansion]
Let $(\phi_n)_{n\ge0}$ be an asymptotic scale. We write
\[
  f(x)\asympseries\sum_{n=0}^{\infty}a_n\phi_n(x)
\]
if, for every fixed $N\ge1$,
\[
  f(x)-\sum_{n=0}^{N-1}a_n\phi_n(x)=\smallo(\bigl(\phi_{N-1}(x)\bigr)).
\]
A slightly stronger and common version requires the remainder to be $\bigO(\phi_N)$.
\end{definition}

This definition deliberately says nothing about convergence of the infinite sum. It only says that each fixed truncation becomes more accurate in a precise relative sense.

\begin{workedbox}{the geometric series viewed from infinity}
For $x>1$,
\[
  \frac{x}{x-1}=\frac{1}{1-x^{-1}}
  =1+x^{-1}+x^{-2}+\cdots.
\]
Here the series actually converges. Its $N$-term remainder is
\[
  \frac{x^{-N}}{1-x^{-1}}=\bigO(x^{-N}),
\]
so it is also an asymptotic expansion. Convergence is a bonus, not part of the definition.
\end{workedbox}

\section{A divergent series can be an excellent approximation}
Consider Euler's factorial series
\[
  \widetilde F(x)=\sum_{n=0}^{\infty}(-1)^n n!\,x^{-n-1}.
\]
For every fixed $x$, the terms eventually grow because
\[
  \frac{(n+1)!x^{-n-2}}{n!x^{-n-1}}=\frac{n+1}{x}.
\]
Thus the series diverges for every finite $x$. Nevertheless, it is the asymptotic expansion of the perfectly ordinary function
\[
  F(x)=\int_0^\infty \frac{\e^{-xp}}{1+p}\,\dd p
      =-\e^x\Ei(-x),
  \qquad x>0.
\]
Repeated integration by parts gives, for every $N$,
\[
  F(x)=\sum_{n=0}^{N-1}(-1)^n n!\,x^{-n-1}
       +(-1)^N N!\int_0^\infty\frac{\e^{-xp}}{(1+p)^{N+1}}\,\dd p.
\]
The final integral is $\bigO(x^{-N-1})$, so the divergent series is asymptotic to $F$.

\begin{intuitionbox}
A convergent series is designed to be summed all the way to infinity. A divergent asymptotic series is designed to be stopped. Its early terms improve the approximation; its late terms eventually make it worse. The location of the smallest term is usually the best stopping point.
\end{intuitionbox}

If the $n$th term behaves like $n!/x^{n+1}$, then it decreases until roughly $n\approx x$ and increases afterward. Truncating near $n=x$ leaves an error on the order of $\e^{-x}$, which is smaller than every power $x^{-N}$. This observation points directly toward transseries.

\section{Beyond all algebraic orders}
The elementary limit
\[
  \e^{-x}=\smallo(x^{-N})\qquad\text{for every fixed }N
\]
means that the exponential $\e^{-x}$ is invisible to the entire power scale
\[
  1,x^{-1},x^{-2},\ldots.
\]
Consequently, the functions $0$ and $\e^{-x}$ have the same asymptotic power series: all coefficients are zero. Likewise,
\[
  \frac{1}{x}+\e^{-x}
  \qquad\text{and}\qquad
  \frac{1}{x}+17\e^{-x}
\]
have the same expansion in inverse powers of $x$.

A power series therefore cannot distinguish functions that differ by an exponentially small amount. A transseries repairs this loss of information by extending the scale:
\[
  1\succ x^{-1}\succ x^{-2}\succ\cdots
  \succ \e^{-x}\succ x^{-1}\e^{-x}\succ\cdots
  \succ \e^{-2x}\succ\cdots.
\]
The notation $A\succ B$ means that $A$ is asymptotically much larger than $B$.

\begin{workedbox}{a boundary layer that ordinary perturbation misses}
Let $0<\varepsilon\ll1$ and solve
\[
  \varepsilon y'(t)+y(t)=1,\qquad y(0)=0.
\]
The exact solution is
\[
  y(t)=1-\e^{-t/\varepsilon}.
\]
For every fixed $t>0$, the exponential is smaller than every power of $\varepsilon$:
\[
  \e^{-t/\varepsilon}=\smallo(\varepsilon^N)
  \quad\text{for all }N.
\]
A power-series ansatz $y\sim y_0+\varepsilon y_1+\cdots$ produces only $y\sim1$. The missing term is not a minor correction near $t=0$; it is the entire mechanism that enforces the initial condition. In the variable $x=1/\varepsilon$, the complete expression is the two-term transseries
\[
  y(t)=1-\e^{-tx}.
\]
\end{workedbox}

\section{Optimal truncation and the first glimpse of resurgence}
Suppose
\[
  a_n\sim C\,A^{-n}\Gamma(n+\beta).
\]
Then the terms $a_n x^{-n}$ are smallest near $n\approx A x$. Stirling's formula suggests that the least term is on the scale
\[
  x^{\beta-1/2}\e^{-Ax}.
\]
Thus factorial divergence in the perturbative coefficients predicts an exponentially small scale. In resurgent problems this is not an accident: the large-order growth of one sector encodes the existence and structure of another sector. We will make this statement concrete in Part~V.

\begin{checkpointbox}
At this stage you should be comfortable with four distinct ideas:
\begin{enumerate}
\item an asymptotic expansion is a hierarchy of remainders, not necessarily a convergent sum;
\item a divergent series may approximate a function extremely well when truncated;
\item exponentials such as $\e^{-x}$ are smaller than every inverse power;
\item a richer expansion must include several asymptotic scales at once.
\end{enumerate}
\end{checkpointbox}

\chapter{The growth ladder}

\section{Comparing familiar functions}
For large positive $x$ we have the chain
\[
  \log\log x\prec\log x\prec x^a\prec \e^{x^b}\prec\e^x\prec\e^{\e^x}
\]
whenever $a,b>0$ and $0<b<1$. The comparison is multiplicative. For example,
\[
  \frac{\log x}{x^a}\to0,
  \qquad
  \frac{x^a}{\e^{x^b}}\to0,
  \qquad
  \frac{\e^x}{\e^{\e^x}}\to0.
\]
This is not merely a list of limits; it behaves like an ordered multiplicative group. Products, reciprocals, and real powers remain meaningful:
\[
  \frac{x^{\sqrt2}\e^{-x}}{(\log x)^3},
  \qquad
  x^\pi\e^{x^{1/2}-2x^{1/3}},
  \qquad
  \exp\!\left(-\frac{x}{\log x}\right).
\]
Such building blocks are called \defterm{transmonomials}.

\section{Three comparison relations}
For positive quantities $f$ and $g$ near infinity:
\begin{align*}
  f\prec g &\quad\Longleftrightarrow\quad f/g\to0,\\
  f\asymp g &\quad\Longleftrightarrow\quad f/g\to c\in(0,\infty),\\
  f\sim g &\quad\Longleftrightarrow\quad f/g\to1.
\end{align*}
The first relation compares orders of magnitude, the second says ``same monomial scale,'' and the third says ``same leading term.'' For instance,
\[
  3x^2+\log x\asymp x^2,
  \qquad
  3x^2+\log x\sim3x^2,
  \qquad
  \log x\prec x^2.
\]

\section{Why the resulting field is non-Archimedean}
In the real numbers, no positive number is larger than every integer. In a formal field containing an infinite symbol $x$, we impose
\[
  x>n\qquad\text{for every }n\in\N.
\]
Then $x^{-1}$ is a positive infinitesimal:
\[
  0<x^{-1}<\frac1n\qquad\text{for every }n\in\N.
\]
An ordered field with such elements is \defterm{non-Archimedean}. This feature is not pathology; it is exactly what allows the field to store asymptotic size algebraically.

\begin{ideabox}
A transseries does not first receive a numerical value and then get compared. Its ordering is read from its largest monomial. The asymptotic hierarchy is built into the syntax.
\end{ideabox}

\section{A visual growth map}
\begin{center}
\begin{tikzpicture}[node distance=9mm and 6mm, every node/.style={font=\small}]
\node[draw,rounded corners,fill=PaleGreen] (tiny) {$\e^{-\e^x}$};
\node[draw,rounded corners,fill=PaleGreen,right=of tiny] (ex) {$\e^{-x}$};
\node[draw,rounded corners,fill=PaleGreen,right=of ex] (powneg) {$x^{-N}$};
\node[draw,rounded corners,fill=SkyBlue,right=of powneg] (one) {$1$};
\node[draw,rounded corners,fill=SkyBlue,right=of one] (log) {$\log x$};
\node[draw,rounded corners,fill=SkyBlue,right=of log] (pow) {$x^N$};
\node[draw,rounded corners,fill=PaleGold,right=of pow] (exp) {$\e^x$};
\node[draw,rounded corners,fill=PaleGold,right=of exp] (ee) {$\e^{\e^x}$};
\draw[-{Stealth[length=2.5mm]},thick] (tiny) -- (ex) -- (powneg) -- (one) -- (log) -- (pow) -- (exp) -- (ee);
\node[below=5mm of one,align=center] {larger monomials to the right\\reciprocals reflect the picture about $1$};
\end{tikzpicture}
\end{center}

The actual transmonomial group is vastly denser than this picture. Between $x$ and $\e^x$ lie $\e^{\sqrt{x}}$, $\e^{x/\log x}$, $\e^{x^{0.999}}$, and infinitely many other scales. Between $\e^{-x}$ and $\e^{-2x}$ lie $\e^{-x-x^{1/2}}$, $x^{-100}\e^{-x}$, and so on.

\section{Comparability and oscillation}
The classical logarithmic-exponential world is designed for eventually non-oscillatory growth. The functions $x$, $\log x$, and $\e^x$ can be consistently ranked. The function $\sin x$ cannot: it changes sign infinitely often and is neither eventually larger nor eventually smaller than zero.

This explains a basic limitation. Standard real LE-transseries do not directly contain a monomial behaving like $\sin x$. Complex exponentials $\e^{ix}$ can be adjoined in other frameworks, but then the simple total ordering is lost. Transseries are exceptionally powerful for regular, non-oscillatory asymptotics; they are not intended to replace Fourier analysis.

\chapter{Your first transseries}

\section{An informal definition}
A transseries is a formal sum
\[
  T=\sum_{\m\in S} a_{\m}\m,
  \qquad a_{\m}\in\R,
\]
where the symbols $\m$ are transmonomials and the support $S$ is arranged so that there is always a largest remaining monomial. Typical examples are
\begin{align*}
  T_1&=x^3-2x+7+x^{-1}+2x^{-2}+6x^{-3}+\cdots,\\
  T_2&=\log x-\log\log x+\frac{\log\log x}{\log x}+\cdots,\\
  T_3&=1+\sum_{n\ge0}a_nx^{-n-1}
      +C\e^{-x}x^2\sum_{n\ge0}b_nx^{-n},\\
  T_4&=\e^{x+\sqrt{x}+x^{1/3}+\cdots}+x^\pi-\log x+\e^{-x}.
\end{align*}
The infinite dots are allowed only when the set of monomials is well organized. Formal generalized series are not arbitrary bags of terms.

\section{Dominant term, magnitude, and sign}
If
\[
  T=3\e^x-x^{100}+7+\e^{-x}+\cdots,
\]
then the dominant monomial is $\e^x$, the dominant term is $3\e^x$, and $T$ is positive because the leading coefficient is positive. We write
\[
  \magop(T)=\e^x,
  \qquad
  \domop(T)=3\e^x.
\]
The lower terms cannot overturn the sign of the leading term, just as $3x^5-10^{100}x^4$ is positive for sufficiently large $x$.

Every nonzero transseries has a multiplicative normal form
\[
  T=a\m(1+\varepsilon),
  \qquad a\in\R^\times,
  \quad \m=\magop(T),
  \quad \varepsilon\prec1.
\]
This one decomposition drives inversion, logarithms, real powers, and many recursive calculations.

\section{Examples of legal and illegal supports}
At $x\to+\infty$, the formal series
\[
  1+x^{-1}+x^{-2}+\cdots
\]
is legal because its monomials decrease. The series
\[
  1+x+x^2+x^3+\cdots
\]
is not a classical transseries support: there is no largest monomial, so the expression has no dominant term.

Likewise,
\[
  \log x+\log\log x+\log\log\log x+\cdots
\]
uses unbounded logarithmic depth. It does not belong to the classical finite-depth LE field, although larger systems such as omega-series and hyperseries are designed to include related expressions.

A third non-example is
\[
  x+\e^x+\e^{\e^x}+\e^{\e^{\e^x}}+\cdots,
\]
which has unbounded exponential height. Again, larger transserial fields can go beyond the classical boundary.

\begin{warningbox}
``Every term gets smaller'' is not by itself sufficient. One also needs a support condition that makes products and infinite sums well defined. The precise condition is reverse well-ordering, introduced in \cref{chap:hahn}.
\end{warningbox}

\section{Four elementary computations}

\begin{workedbox}{a geometric inverse}
Since $x^{-1}\prec1$,
\[
  (1-x^{-1})^{-1}=1+x^{-1}+x^{-2}+\cdots.
\]
Therefore
\[
  (x-1)^{-1}=x^{-1}+x^{-2}+x^{-3}+\cdots.
\]
Multiplication is formal, and every coefficient receives only finitely many contributions.
\end{workedbox}

\begin{workedbox}{the inverse of $\e^x+x$}
Factor out the dominant monomial:
\[
  \e^x+x=\e^x(1+x\e^{-x}).
\]
The correction $x\e^{-x}$ is small, so
\[
  (\e^x+x)^{-1}
  =\e^{-x}\sum_{n\ge0}(-1)^n x^n\e^{-nx}
  =\sum_{n\ge0}(-1)^n x^n\e^{-(n+1)x}.
\]
The ordinary geometric series has been transplanted to an exponential scale.
\end{workedbox}

\begin{workedbox}{the logarithm of $\sinh x$}
For large positive $x$,
\[
  \sinh x=\frac{\e^x}{2}(1-\e^{-2x}).
\]
Hence
\[
  \log(\sinh x)
  =x-\log2+\log(1-\e^{-2x})
  =x-\log2-\sum_{n\ge1}\frac{\e^{-2nx}}{n}.
\]
An elementary function becomes an infinite exponential transseries.
\end{workedbox}

\begin{workedbox}{an algebraic equation with exponential sectors}
Solve
\[
  y^2-y=\e^{-x}.
\]
The two roots are
\[
  y_{\pm}=\frac{1\pm\sqrt{1+4\e^{-x}}}{2}.
\]
Expanding the square root gives
\begin{align*}
  y_+&=1+\e^{-x}-\e^{-2x}+2\e^{-3x}-5\e^{-4x}+\cdots,\\
  y_-&=-\e^{-x}+\e^{-2x}-2\e^{-3x}+5\e^{-4x}-\cdots.
\end{align*}
The coefficients are signed Catalan numbers. Even an algebraic equation can naturally produce a purely exponential transseries.
\end{workedbox}

\section{\texorpdfstring{The decomposition $T=L+c+S$}{The decomposition T=L+c+S}}
Every real transseries can be split uniquely as
\[
  T=L+c+S,
\]
where $L$ is \defterm{purely large} (all its monomials are $\succ1$), $c\in\R$, and $S$ is \defterm{small} (all its monomials are $\prec1$). For example,
\[
  x+\sqrt{x}-3+\frac{\log x}{x}+\e^{-x}
  =\underbrace{x+\sqrt{x}}_{L}
   \underbrace{-3}_{c}
   +\underbrace{\frac{\log x}{x}+\e^{-x}}_{S}.
\]
Then exponentiation is defined by
\[
  \e^T=\e^L\e^c\sum_{n\ge0}\frac{S^n}{n!}.
\]
This formula illustrates the recursive nature of the subject: the new monomial $\e^L$ is multiplied by a generalized power series in the small quantity $S$.

\chapter{What the word ``transseries'' can mean}

\section{One word, several closely related objects}
The literature uses ``transseries'' for a family of constructions rather than one universally fixed set. The differences matter, but they should not obscure the common idea.

\subsection{Formal logarithmic-exponential transseries}
These are generalized series built recursively from real constants, an infinite variable $x$, real powers, exponentials, logarithms, and admissible infinite sums. They form an ordered differential field usually denoted $\Tfield$. This is the main algebraic object in Parts~II--IV. Edgar's beginner exposition \cite{Edgar2009}, the constructions of van den Dries--Macintyre--Marker \cite{DMM1997,DMM2001}, and van der Hoeven's monograph \cite{Hoeven2006} are central references.

\subsection{Grid-based and well-based transseries}
Grid-based transseries have supports contained in finitely generated multiplicative grids. They are convenient for algorithms and explicit calculations. Well-based transseries allow arbitrary reverse well-ordered supports and are more flexible. Most formulas look the same, but some closure and convergence arguments require separate proofs. We compare them in \cref{chap:grids}.

\subsection{\'{E}calle's analyzable and resurgent transseries}
In resurgence theory a transseries often has the sector form
\[
  \Phi(x;\sigma)=\sum_{k\ge0}\sigma^k\e^{-kAx}x^{k\beta}\Phi_k(x),
  \qquad
  \Phi_k(x)=\sum_{n\ge0}a_{k,n}x^{-n}.
\]
The focus is not only formal algebra. One studies Borel transforms, analytic continuation, Stokes automorphisms, and the way different sectors determine one another. This tradition descends from \'{E}calle's work \cite{Ecalle1992} and is developed in modern expositions such as Costin \cite{Costin2009}, Dorigoni \cite{Dorigoni2019}, and Aniceto--Ba\c{s}ar--Schiappa \cite{ABS2019}.

\subsection{Transseries as germs or surreal functions}
A \defterm{germ at infinity} identifies two functions if they agree for all sufficiently large $x$. Hardy fields organize non-oscillating germs. In another direction, one can evaluate suitable formal series at positive infinite surreal numbers. The $H$-field viewpoint creates a bridge among germs, transseries, and surreal numbers \cite{ADH2017,ADH2018,BM2018,BM2019}.

\section{Formal object versus represented function}
It is essential to separate three statements:
\begin{enumerate}
\item $T$ is a legal formal transseries.
\item A function $f$ has $T$ as an asymptotic expansion.
\item A summation procedure assigns a specific analytic function to $T$.
\end{enumerate}
The first is algebraic. The second is an estimate about remainders. The third requires extra analytic information and may depend on a direction in the complex plane.

\begin{warningbox}
A formal transseries is not automatically a convergent numerical expression. Manipulating it can be perfectly rigorous even when no ordinary infinite sum exists. This is no stranger than manipulating a formal power series in algebra.
\end{warningbox}

\section{A useful working convention}
Throughout the formal chapters we use real, well-based logarithmic-exponential transseries at $x\to+\infty$. When a proof is easiest in the grid-based setting, we say so. In the analytic chapters, $x$ may be complex and the word ``sector'' refers to an angular region where a Borel-Laplace sum is defined.

\section{Historical orientation}
The intellectual roots include du Bois-Reymond's orders of infinity, Hardy's logarithmico-exponential functions, Hahn's generalized series, and Borel's theory of divergent series. Modern transseries were developed independently in the 1980s by Dahn and G\"oring in connection with exponential fields and by \'{E}calle in his work on analyzable functions and the Dulac problem. Later developments connected the subject with computer algebra, o-minimality, differential algebra, model theory, resurgence, and surreal numbers; see \cite{DahnGoring1987,Ecalle1992,DMM2001,Hoeven2006,ADH2017}.

\begin{checkpointbox}
You are now ready for the formal core. The key ingredients will be:
\begin{enumerate}
\item an ordered multiplicative group of monomials;
\item generalized series whose supports are reverse well ordered;
\item recursive closure under exponentials and logarithms;
\item operations controlled by dominant monomials and small remainders.
\end{enumerate}
\end{checkpointbox}

\part{Generalized series and the formal construction}

\chapter{Ordered monomial groups}
\label{chap:monomials}

\section{Why monomials should form a group}
Ordinary Laurent series use monomials $x^n$ with $n\in\Z$. They multiply according to
\[
  x^m x^n=x^{m+n},
  \qquad
  (x^n)^{-1}=x^{-n}.
\]
The exponents form the additive group $\Z$, while the monomials form a multiplicative group. Generalized series keep this pattern but replace $\Z$ by a much richer ordered group.

\begin{definition}[ordered abelian group]
An \defterm{ordered abelian group} $(\G,\cdot,1,\prec)$ is an abelian group together with a total order preserved by multiplication:
\[
  \m\prec\n\quad\Longrightarrow\quad \m\g\prec\n\g
  \qquad(\g\in\G).
\]
We think of elements of $\G$ as positive monomials. The notation $\m\prec\n$ means that $\m$ is asymptotically smaller than $\n$.
\end{definition}

The simplest example is
\[
  \G_0=x^\R:=\set{x^a:a\in\R},
\]
ordered by $x^a\prec x^b$ exactly when $a<b$. This already allows generalized power series such as
\[
  3x^{\sqrt2}-7x^\pi+x^{-1/3}+\cdots.
\]

\section{Multiplicative notation and logarithmic coordinates}
Ordered abelian groups are often written additively. If $\Gamma$ is additive, the associated monomial group may be written
\[
  t^\Gamma=\set{t^\gamma:\gamma\in\Gamma},
  \qquad t^\alpha t^\beta=t^{\alpha+\beta}.
\]
Transseries are more naturally written multiplicatively because their monomials look like products and exponentials. Formally, however, taking a logarithm turns the monomial group into an additive ordered group:
\[
  \log(\m\n)=\log\m+\log\n.
\]
For instance,
\[
  \m=x^a\e^L
  \qquad\Longrightarrow\qquad
  \log\m=a\log x+L.
\]
This observation is what makes exponentiation fit recursively into the construction.

\section{Coefficients are not part of the monomial}
A term has the form $a\m$ with $a\in\R^\times$ and $\m\in\G$. The coefficient $a$ affects sign and leading constant but not order of magnitude. Thus
\[
  7x^3\asymp -10^{100}x^3,
  \qquad
  x^3\succ x^2,
\]
although $7x^3$ and $-10^{100}x^3$ have opposite signs.

The separation of coefficient and monomial gives a unique normal form for each nonzero term. This uniqueness is essential when like terms are collected.

\section{A lexicographic example}
Consider monomials
\[
  \m=x^a\e^L,
\]
where $L$ is purely large. Compare two of them by first comparing the exponential parts:
\[
  x^{a_1}\e^{L_1}\prec x^{a_2}\e^{L_2}
\]
when either $L_1<L_2$, or $L_1=L_2$ and $a_1<a_2$. This is a lexicographic order. It expresses the elementary fact that an exponential difference dominates every power:
\[
  L_2-L_1\succ1
  \quad\Longrightarrow\quad
  \frac{x^{a_1}\e^{L_1}}{x^{a_2}\e^{L_2}}
  =x^{a_1-a_2}\e^{-(L_2-L_1)}\prec1.
\]

\begin{workedbox}{sorting a mixed collection of monomials}
Arrange the following from largest to smallest:
\[
  \e^{x+\sqrt{x}},\quad x^{100}\e^x,\quad x^{-3}\e^x,
  \quad \e^{\sqrt{x}},\quad x^\pi,\quad \e^{-x}.
\]
First compare exponential exponents:
\[
  x+\sqrt{x}>x>\sqrt{x}>0>-x.
\]
Within the equal exponent $x$, compare powers of $x$. The result is
\[
  \e^{x+\sqrt{x}}
  \succ x^{100}\e^x
  \succ x^{-3}\e^x
  \succ \e^{\sqrt{x}}
  \succ x^\pi
  \succ \e^{-x}.
\]
\end{workedbox}

\section{Monomial logarithmic derivatives}
For a nonzero differentiable monomial $\m$, define
\[
  \m^\dagger:=\frac{\m'}{\m}.
\]
This is the \defterm{logarithmic derivative}. If $\m=x^a\e^L$, then
\[
  \m^\dagger=\frac{a}{x}+L'.
\]
The logarithmic derivative is often simpler than the derivative itself and is a central tool for asymptotic integration and differential equations.

\begin{example}
For
\[
  \m=x^3\exp\!\left(-x^2+\sqrt{x}\right),
\]
we have
\[
  \m^\dagger=\frac3x-2x+\frac{1}{2\sqrt{x}},
  \qquad
  \m'=\left(-2x+\frac{1}{2\sqrt{x}}+\frac3x\right)\m.
\]
The dominant factor in the derivative is $-2x\m$.
\end{example}

\section{The role of finite height}
A classical transmonomial is built using only finitely many nested exponentials and logarithms. Each individual monomial therefore has finite complexity. A transseries may have infinitely many terms, but there is a finite bound on the exponential height and logarithmic depth appearing in a given classical LE-transseries.

This restriction makes the recursive construction manageable. It also explains why an infinite exponential tower is not an element of the classical field. Later systems of hyperseries relax this boundary.

\chapter{Hahn series and well-based support}
\label{chap:hahn}

\section{Generalized series}
Fix an ordered abelian monomial group $\G$. A formal generalized series is an expression
\[
  F=\sum_{\m\in\G}a_{\m}\m,
\]
where only a selected set of coefficients is nonzero. The support is
\[
  \supp F=\set{\m\in\G:a_{\m}\ne0}.
\]
If arbitrary supports were allowed, multiplication could require infinite, undefined coefficient sums. Hahn's insight was to impose an order condition on the support.

\begin{definition}[reverse well-ordered support]
A subset $S\subseteq\G$ is \defterm{reverse well ordered} if every nonempty subset of $S$ has a largest element. Equivalently, there is no infinite strictly increasing chain
\[
  \m_1\prec\m_2\prec\m_3\prec\cdots
\]
inside $S$.
\end{definition}

Many authors reverse the group order and simply say ``well ordered.'' In transseries it is natural to write terms from largest to smallest, so ``reverse well ordered'' is the transparent phrase.

\begin{example}
The set
\[
  \set{1,x^{-1},x^{-2},\ldots}
\]
is reverse well ordered. Its largest member is $1$, and every tail has a largest member. By contrast,
\[
  \set{1,x,x^2,\ldots}
\]
has no largest member and is not reverse well ordered.
\end{example}

\section{Hahn series}
\begin{definition}[Hahn field notation]
The set of all real generalized series with reverse well-ordered support in $\G$ is denoted
\[
  \R((\G))
\]
or, in transseries typography, $\R\llbracket\G\rrbracket$.
\end{definition}

A nonzero Hahn series automatically has a largest monomial, hence a dominant term and a sign. This turns the formal series field into an ordered field.

\section{Addition}
Addition is coefficientwise:
\[
  \left(\sum_{\m}a_{\m}\m\right)
  +\left(\sum_{\m}b_{\m}\m\right)
  =\sum_{\m}(a_{\m}+b_{\m})\m.
\]
The support of the sum is contained in the union of the two supports. A finite union of reverse well-ordered sets is reverse well ordered, so addition is safe.

\section{Multiplication and the Hahn--Neumann mechanism}
Formally distribute:
\[
  \left(\sum_{\m}a_{\m}\m\right)
  \left(\sum_{\n}b_{\n}\n\right)
  =\sum_{\g\in\G}
    \left(\sum_{\m\n=\g}a_{\m}b_{\n}\right)\g.
\]
Two facts are required:
\begin{enumerate}
\item the product support is reverse well ordered;
\item for each fixed $\g$, only finitely many pairs $(\m,\n)$ contribute.
\end{enumerate}
A classical lemma of Hahn and Neumann supplies these facts under the support hypotheses. This is the hidden combinatorial engine behind generalized-series multiplication.

\begin{workedbox}{a support of order type $\omega^2$}
Let
\[
  U=\sum_{j\ge1}j\e^{-jx},
  \qquad
  V=\sum_{k\ge0}x^{-k}.
\]
Then
\[
  UV=\sum_{j\ge1}\sum_{k\ge0}j\,x^{-k}\e^{-jx}.
\]
For fixed $j$, the powers $x^{-k}\e^{-jx}$ descend as $k$ increases. After all terms with $j=1$ come those with $j=2$, and so on. The support has the nested ordering pattern
\[
  \omega+\omega+\omega+\cdots=\omega^2.
\]
Transfinite order types appear naturally even though every displayed calculation is elementary.
\end{workedbox}

\section{Summable families}
To define an infinite sum of Hahn series, one needs a stronger notion than termwise convergence.

\begin{definition}[summable family]
A family $(F_i)_{i\in I}$ of Hahn series is \defterm{summable} if
\begin{enumerate}
\item the union $\bigcup_i\supp F_i$ is reverse well ordered;
\item each monomial belongs to the supports of only finitely many $F_i$.
\end{enumerate}
Then $\sum_iF_i$ is defined by finite coefficientwise summation at every monomial.
\end{definition}

This notion explains why the geometric series
\[
  \sum_{n\ge0}\varepsilon^n
\]
is legal whenever $\varepsilon\prec1$: the supports move downward, and each target monomial receives finitely many contributions under the standard hypotheses.

\begin{warningbox}
Formal summability is not numerical convergence. A summable family may define a valid Hahn series even when substituting a finite real number for $x$ makes the numerical series diverge.
\end{warningbox}

\section{The field theorem}
\begin{theorem}[Hahn]
If $K$ is a field and $\G$ is an ordered abelian group, then the generalized series $K((\G))$ with reverse well-ordered support form a field.
\end{theorem}

The nontrivial step is inversion. Let
\[
  F=a\m(1+\varepsilon),
  \qquad \varepsilon\prec1.
\]
Then
\[
  F^{-1}=a^{-1}\m^{-1}
  \sum_{n\ge0}(-\varepsilon)^n.
\]
The generalized geometric family is summable, so the inverse exists formally.

\begin{proof}[Proof idea]
The leading term $a\m$ can be inverted in the coefficient field and monomial group. The remainder $\varepsilon$ is small. Every power $\varepsilon^n$ has monomials strictly below those in $\varepsilon^{n-1}$ in the relevant multiplicative sense. Neumann's lemma shows that the union of supports remains well based and that each coefficient is a finite sum. Formal multiplication then gives
\[
  (1+\varepsilon)\sum_{n\ge0}(-\varepsilon)^n=1.
\]
\end{proof}

\section{Truncations}
For a cutoff monomial $\g$, define the truncation
\[
  F_{\succeq\g}:=
  \sum_{\m\succeq\g}a_{\m}\m.
\]
A field of generalized series is \defterm{truncation closed} if these initial segments remain in the field. Classical transseries fields are constructed to have strong truncation properties. Truncation is the formal analogue of computing ``all terms down to a requested size.''

\begin{ideabox}
A transseries computation is exact above a chosen cutoff. One need not compute an infinite expression all at once. The support order guarantees that only finitely much algebra influences any fixed leading segment.
\end{ideabox}

\chapter{Grid-based versus well-based series}
\label{chap:grids}

\section{Ratio sets}
A grid-based series begins with finitely many small monomials
\[
  \boldsymbol\mu=(\mu_1,\ldots,\mu_r),
  \qquad \mu_i\prec1.
\]
For an integer vector $\mathbf{k}=(k_1,\ldots,k_r)$ write
\[
  \boldsymbol\mu^{\mathbf{k}}
  =\mu_1^{k_1}\cdots\mu_r^{k_r}.
\]
Given a lower-bound vector $\mathbf{m}\in\Z^r$, the associated grid is
\[
  \mathfrak{J}^{\boldsymbol\mu,\mathbf{m}}
  =\set{\boldsymbol\mu^{\mathbf{k}}:\mathbf{k}\ge\mathbf{m}}
\]
with coordinatewise inequality.

\begin{example}
Take $\mu_1=x^{-1}$ and $\mu_2=\e^{-x}$. The grid contains
\[
  x^{-k}\e^{-jx}
  \qquad(k,j\ge0).
\]
A two-scale transseries such as
\[
  \sum_{j\ge0}\e^{-jx}\sum_{k\ge0}a_{j,k}x^{-k}
\]
is grid based.
\end{example}

\section{Why finite grids help computation}
A finite ratio set turns support bookkeeping into integer-vector bookkeeping. Multiplication adds exponent vectors. To compute down to a cutoff, algorithms can enumerate a finite region of $\Z^r$. Dickson's lemma---every subset of $\N^r$ has finitely many minimal elements---underlies many finiteness arguments.

Grid-based transseries are therefore natural in computer algebra. Van der Hoeven's approach develops effective operations and Newton methods in this setting \cite{Hoeven2006}.

\section{Why well-based series are more general}
Not every reverse well-ordered support lies in a finitely generated grid. For example, a family of monomials may involve infinitely many independent scales. Well-based series permit such supports as long as the order condition holds.

The tradeoff is familiar:
\begin{center}
\begin{tabularx}{0.92\textwidth}{@{}>{\bfseries}lXX@{}}
\toprule
 & Grid based & Well based\\
\midrule
Support & contained in a finite-rank grid & arbitrary reverse well-ordered set\\
Algorithms & especially concrete & more abstract\\
Closure proofs & often combinatorial in $\Z^r$ & rely on general support lemmas\\
Generality & smaller class & larger class\\
\bottomrule
\end{tabularx}
\end{center}

\section{Witnesses and ratio sets}
In grid-based work, a finite ratio set can \defterm{witness} a comparison such as $A\prec B$ by expressing every monomial in $A/B$ as a positive product of the chosen small ratios. This gives an effective certificate that a remainder is small.

For example, with $\boldsymbol\mu=(x^{-1},\e^{-x})$, the transseries
\[
  S=3x^{-2}+x\e^{-x}+\e^{-2x}
\]
is small because every monomial is generated by products of $x^{-1}$ and $\e^{-x}$ with exponents in a suitable grid. The positive power $x$ in $x\e^{-x}$ does not make the term large; the exponential factor wins.

\section{Which version will we use?}
We state the basic algebra in well-based language because it reveals the general idea. Most explicit examples live in a small grid and may be computed as if one were manipulating multivariate formal power series. When composition, convergence, or fractional iteration has version-dependent subtleties, we mention them.

\chapter[Constructing logarithmic-exponential transseries]{Constructing logarithmic-exponential\\transseries}
\label{chap:construction}

\section{Stage zero: generalized powers}
Begin with the monomial group
\[
  \G_0=x^\R
\]
and the Hahn field
\[
  \Tfield_0=\R((x^\R)).
\]
Elements of $\Tfield_0$ are generalized power series
\[
  \sum_{a\in S}c_a x^a
\]
with reverse well-ordered exponent support. This stage already contains Laurent series, Puiseux series, and real powers.

Every $F\in\Tfield_0$ decomposes as
\[
  F=L+c+S,
\]
where $L$ contains powers $x^a$ with $a>0$, $c$ is constant, and $S$ contains powers with $a<0$.

\section{Adjoining exponentials}
Assume a field $\Tfield_n$ has been built. Let $\Tfield_n^{\mathrm{pl}}$ denote its purely large elements. Form new monomials
\[
  \G_{n+1}
  =\set{x^b\e^L:b\in\R,\ L\in\Tfield_n^{\mathrm{pl}}}.
\]
Then form the next Hahn field
\[
  \Tfield_{n+1}=\R((\G_{n+1})).
\]
The \defterm{exponential height} counts how many such stages are needed.

\begin{example}
The monomial $\e^{x^2+\sqrt{x}}$ has height one over generalized powers. The monomial
\[
  \exp\!\left(\e^x+x^{1/2}\right)
\]
has height two because its exponent already contains $\e^x$.
\end{example}

The log-free field is the union
\[
  \Tfield_{\bullet}=\bigcup_{n\ge0}\Tfield_n,
\]
with compatible embeddings between stages. The full construction takes care that a series represented at one stage keeps the same meaning at later stages.

\section{Defining the exponential of an arbitrary transseries}
Let
\[
  T=L+c+S
\]
with $L$ purely large, $c\in\R$, and $S\prec1$. Define
\[
  \exp(T)=\e^L\e^c\sum_{n\ge0}\frac{S^n}{n!}.
\]
The factor $\e^L$ is a monomial introduced at the next height, while the Taylor series in $S$ is a summable Hahn family.

\begin{workedbox}{expanding an exponential}
Let
\[
  T=x+\log x+x^{-1}+2\e^{-x}.
\]
Treat $x$ as purely large, $\log x$ as part of the monomial exponent after logarithms are admitted, and
\[
  S=x^{-1}+2\e^{-x}
\]
as small. Then
\[
  \e^T=x\e^x\exp(x^{-1}+2\e^{-x}).
\]
To a few orders,
\begin{align*}
  \exp(x^{-1}+2\e^{-x})
  &=1+x^{-1}+\frac{1}{2}x^{-2}
    +2\e^{-x}+2x^{-1}\e^{-x}+2\e^{-2x}+\cdots.
\end{align*}
The ordering decides which of these terms should be displayed first.
\end{workedbox}

\section{Adjoining logarithms}
A positive transseries has normal form
\[
  T=a\m(1+\varepsilon),
  \qquad a>0,
  \quad \varepsilon\prec1.
\]
Define
\[
  \log T
  =\log a+\log\m+
   \sum_{n\ge1}\frac{(-1)^{n+1}}{n}\varepsilon^n.
\]
For a monomial $\m=x^b\e^L$,
\[
  \log\m=b\log x+L.
\]
To make $\log x,\log\log x,$ and further iterates available, one may substitute $\log_M x$ for $x$ in the log-free field, where
\[
  \log_0x=x,
  \qquad
  \log_{M+1}x=\log(\log_Mx).
\]
The classical LE field is a union over finite height and finite logarithmic depth.

\section{Why the recursive definition is not circular}
At first sight, ``a monomial is the exponential of a transseries, and a transseries is a sum of monomials'' sounds circular. The height filtration resolves the issue:
\begin{enumerate}
\item stage $0$ monomials are known;
\item exponentials at stage $n+1$ use purely large exponents from earlier stages;
\item Hahn sums are then formed from the newly known monomials.
\end{enumerate}
Every individual object appears at a finite stage, so the recursion is well founded.

\section{A schematic of the construction}
\begin{center}
\begin{tikzpicture}[node distance=7mm, every node/.style={align=center,font=\small}]
\node[draw,rounded corners,fill=SkyBlue] (g0) {$\G_0=x^\R$};
\node[draw,rounded corners,fill=PaleGreen,below=of g0] (t0) {$\Tfield_0=\R((\G_0))$};
\node[draw,rounded corners,fill=PaleGold,below=of t0] (g1) {$\G_1=\{x^b\e^L:L\in\Tfield_0^{\rm pl}\}$};
\node[draw,rounded corners,fill=PaleGreen,below=of g1] (t1) {$\Tfield_1=\R((\G_1))$};
\node[below=of t1] (dots) {$\vdots$};
\node[draw,rounded corners,fill=SkyBlue,below=of dots] (union) {$\Tfield_{\bullet}=\bigcup_n\Tfield_n$};
\node[draw,rounded corners,fill=PaleRed,right=18mm of union] (logs) {substitute $\log_Mx$\\and take the union over $M$};
\node[draw,very thick,rounded corners,fill=white,right=18mm of logs] (T) {classical LE field $\Tfield$};
\draw[-{Stealth},thick] (g0)--(t0);
\draw[-{Stealth},thick] (t0)--(g1);
\draw[-{Stealth},thick] (g1)--(t1);
\draw[-{Stealth},thick] (t1)--(dots);
\draw[-{Stealth},thick] (dots)--(union);
\draw[-{Stealth},thick] (union)--(logs);
\draw[-{Stealth},thick] (logs)--(T);
\end{tikzpicture}
\end{center}

\begin{advancedbox}
There are several equivalent or closely related constructions: logarithmic-exponential series, exponential-logarithmic series, grid-based transseries, well-based transseries, and fields built inside the surreal numbers. Their exact closure procedures differ. For a first course, the height-and-depth picture above captures the operational core, while the references supply the set-theoretic details \cite{DMM2001,Hoeven2006,KuhlmannTressl2011,Schmeling2001}.
\end{advancedbox}

\chapter{Dominance and normal forms}
\label{chap:valuation}

\section{The dominant monomial}
Let
\[
  T=\sum_{\m\in S}a_{\m}\m\ne0.
\]
Because $S$ is reverse well ordered, it has a largest element $\m_0$. Define
\[
  \magop(T)=\m_0,
  \qquad
  \domop(T)=a_{\m_0}\m_0.
\]
Then
\[
  T\sim\domop(T).
\]
The sign of $T$ is the sign of $a_{\m_0}$.

\section{Formal asymptotic relations}
For nonzero transseries $A,B$ define
\begin{align*}
  A\prec B
  &\quad\Longleftrightarrow\quad
    \magop(A)\prec\magop(B),\\
  A\asymp B
  &\quad\Longleftrightarrow\quad
    \magop(A)=\magop(B),\\
  A\sim B
  &\quad\Longleftrightarrow\quad
    A-B\prec B.
\end{align*}
The last condition means that $A$ and $B$ have the same dominant term, not merely the same dominant monomial.

\begin{example}
Let
\[
  A=3x^2+x,
  \quad B=5x^2-7,
  \quad C=3x^2+\e^{-x}.
\]
Then $A\asymp B$ but $A\not\sim B$, since their leading coefficients differ. On the other hand, $A\sim C$: indeed, $A-C=x-\e^{-x}\asymp x\prec x^2\asymp A$. Thus $A$ and $C$ share the dominant term $3x^2$, even though their lower-order parts look very different.
\end{example}

\section{Large, bounded, and infinitesimal elements}
A transseries $T$ is
\begin{itemize}
\item \defterm{large} if $|T|>r$ for every real $r$;
\item \defterm{bounded} if $|T|<r$ for some positive real $r$;
\item \defterm{infinitesimal} if $|T|<r$ for every positive real $r$.
\end{itemize}
In monomial language,
\[
  T\text{ large}\iff \magop(T)\succ1,
  \qquad
  T\text{ infinitesimal}\iff \magop(T)\prec1.
\]
A bounded transseries has the form
\[
  c+\varepsilon,
  \qquad c\in\R,
  \quad\varepsilon\prec1.
\]

\section{The natural valuation ring}
The bounded elements form a valuation ring
\[
  \mathcal O=\set{T:|T|\le r\text{ for some }r\in\R_{>0}}.
\]
Its maximal ideal is the set of infinitesimals
\[
  \mathfrak o=\set{T:T\prec1}.
\]
The map
\[
  \operatorname{st}:\mathcal O\to\R,
  \qquad c+\varepsilon\mapsto c,
\]
is the formal \defterm{standard part} or residue map.

\begin{intuitionbox}
The valuation forgets coefficients and remembers order of magnitude. The residue map forgets infinitesimal corrections and remembers the constant part. Together, ordering and valuation turn asymptotic estimates into algebra.
\end{intuitionbox}

\section{Multiplicative normal form}
Every nonzero $T$ can be written uniquely as
\[
  T=a\m(1+\varepsilon),
  \qquad a\in\R^\times,
  \quad \varepsilon\prec1.
\]
From this one obtains immediately:
\begin{align*}
  T^{-1}
  &=a^{-1}\m^{-1}\sum_{n\ge0}(-\varepsilon)^n,\\
  \log T
  &=\log a+\log\m+
    \sum_{n\ge1}\frac{(-1)^{n+1}}{n}\varepsilon^n
    \quad (T>0),\\
  T^\alpha
  &=a^\alpha\m^\alpha
    \sum_{n\ge0}\binom{\alpha}{n}\varepsilon^n
    \quad (T>0,\ \alpha\in\R).
\end{align*}

\section{Leading-term arithmetic}
If $A,B\ne0$, then
\[
  \magop(AB)=\magop(A)\magop(B),
  \qquad
  \domop(AB)=\domop(A)\domop(B).
\]
For sums, the larger magnitude wins unless leading terms cancel. If $A\succ B$, then
\[
  A+B\sim A.
\]
If $A\asymp B$, one must add the leading coefficients; cancellation may reveal a lower scale.

\begin{workedbox}{cancellation exposes a hidden scale}
Let
\[
  A=x+1+\e^{-x},
  \qquad
  B=-x-1+x^{-3}.
\]
Individually, $A\asymp B\asymp x$. But
\[
  A+B=x^{-3}+\e^{-x}\sim x^{-3},
\]
because $x^{-3}\succ\e^{-x}$. Two cancellations move the answer from the scale $x$ to the scale $x^{-3}$.
\end{workedbox}

\section{Formal limits}
The order language lets us write analogues of limits:
\[
  T\to c
  \quad\text{means}\quad
  T-c\prec1,
\]
and
\[
  T\to+\infty
  \quad\text{means}\quad
  T>\R.
\]
These are algebraic statements inside the transseries field. When a transseries represents a germ, they agree with the usual analytic interpretation.

\begin{checkpointbox}
The entire formal calculus rests on two normal forms:
\[
  T=L+c+S
  \quad\text{and}\quad
  T=a\m(1+\varepsilon).
\]
The first is adapted to exponentiation; the second is adapted to inversion, logarithms, and powers.
\end{checkpointbox}

\part{Algebra, calculus, and composition}

\chapter{Formal arithmetic: familiar rules on an unfamiliar scale}
\label{chap:arithmetic}

\section{Why the ordinary algebraic rules still work}
A transseries is more elaborate than a polynomial, but its arithmetic is built from the same two operations:
\[
  \left(\sum_{\m}a_{\m}\m\right)
  +\left(\sum_{\m}b_{\m}\m\right)
  =\sum_{\m}(a_{\m}+b_{\m})\m,
\]
and
\[
  \left(\sum_{\m}a_{\m}\m\right)
  \left(\sum_{\n}b_{\n}\n\right)
  =\sum_{\g}
   \left(\sum_{\m\n=\g}a_{\m}b_{\n}\right)\g.
\]
The support conditions from \cref{chap:hahn} are what make the inner coefficient sum finite. Once this finiteness theorem has been proved, addition and multiplication may be performed exactly as one would multiply ordinary formal power series.

\begin{workedbox}{multiplying two sectors}
Let
\[
  A=1+\frac{2}{x}+3\e^{-x},
  \qquad
  B=1-\frac{1}{x}+5\e^{-x}.
\]
Then
\begin{align*}
  AB={}&1+\frac1x-\frac{2}{x^2}
       +8\e^{-x}+\frac{7}{x}\e^{-x}
       +15\e^{-2x}.
\end{align*}
There is no mystery: collect equal monomials, but remember that
$x^{-1}\e^{-x}$ and $\e^{-x}$ are distinct monomials.
\end{workedbox}

\section{Inversion by separating the leading monomial}
Suppose
\[
  T=a\m(1+\varepsilon),
  \qquad \varepsilon\prec1.
\]
The formal geometric series gives
\[
  \frac1T=a^{-1}\m^{-1}
  \left(1-\varepsilon+\varepsilon^2-\varepsilon^3+\cdots\right).
\]
This is legitimate because the family $(\varepsilon^n)_{n\ge0}$ is summable: for every fixed monomial, only finitely many powers can contribute to its coefficient.

\begin{workedbox}{an inverse containing two small scales}
Write
\[
  T=x\left(1+x^{-1}+x\e^{-x}\right).
\]
Set $u=x^{-1}+x\e^{-x}$. Then
\[
  T^{-1}=x^{-1}(1-u+u^2-u^3+\cdots).
\]
Keeping terms down to the scale $x^{-3}$ and the first exponential sector gives
\begin{align*}
  T^{-1}
  ={}&x^{-1}-x^{-2}+x^{-3}
      -\e^{-x}+2x^{-1}\e^{-x}
      -3x^{-2}\e^{-x}
      +x\e^{-2x}+\cdots.
\end{align*}
The order in which one displays the terms is dictated by magnitude, not by the order in which the algebra produced them.
\end{workedbox}

\section{Infinite sums are allowed only when they are summable}
The symbol
\[
  \sum_{i\in I}T_i
\]
means a transseries only when the family is summable in the Hahn sense. Two conditions are useful to remember:
\begin{enumerate}
\item the union of all supports must be well based;
\item each fixed monomial must occur in only finitely many $T_i$.
\end{enumerate}
For example,
\[
  \sum_{n\ge0}x^{-n}
\]
is summable, while
\[
  \sum_{n\ge0}x^n
\]
is not a legal expansion at $x=+\infty$. The second family keeps introducing larger monomials and has no largest term.

\begin{warningbox}
The phrase ``formal'' does not mean ``anything goes.'' Formal transseries have strict support rules. Those rules replace numerical convergence and are what prevent ambiguous rearrangements.
\end{warningbox}

\section{A transseries parameter is a coefficient, not a monomial}
Many differential equations produce a family
\[
  T(x;C)=\sum_{k\ge0}C^k\e^{-kAx}\Phi_k(x),
\]
where $C$ is an arbitrary constant. One should distinguish two viewpoints:
\begin{itemize}
\item for a fixed real or complex $C$, this is an ordinary transseries in $x$;
\item algebraically, one may work over the coefficient field $\R(C)$ or with a formal parameter $C$.
\end{itemize}
The exponential scale $\e^{-Ax}$ belongs to the monomial group. The integration constant $C$ belongs to the coefficient field. Confusing these roles leads to errors when comparing magnitudes.

\section{Truncation is an algebraic operation}
If $\n$ is a monomial, write
\[
  T\big|_{\succeq\n}
\]
for the sum of all terms of $T$ whose monomials are at least $\n$. Because supports are well based, this is a legitimate transseries. Truncation is not generally a field homomorphism: multiplying two discarded tails can sometimes contribute at a retained scale. Nevertheless, if one keeps a sufficiently generous ``guard band,'' finite truncations can be computed reliably.

\begin{ideabox}
Computer algebra with transseries is possible because every requested finite output depends on only finitely many input terms. The global object may be infinite, but a well-chosen cutoff turns each calculation into finite bookkeeping.
\end{ideabox}

\chapter{Exponentials, logarithms, and real powers}
\label{chap:explog}

\section{The three-part additive decomposition}
Recall that every transseries has a unique decomposition
\[
  T=L+c+S,
\]
where $L$ is purely large, $c$ is constant, and $S$ is infinitesimal. Exponentiation separates these three scales:
\[
  \e^T=\underbrace{\e^L}_{\text{new monomial}}
        \underbrace{\e^c}_{\text{coefficient}}
        \underbrace{\left(1+S+\frac{S^2}{2!}+\cdots\right)}_{\text{unit}}.
\]
This formula is the conceptual center of the LE construction.

\begin{workedbox}{exponentiating a mixed transseries}
Let
\[
 T=x^2-3x+2\log x+\log 7+x^{-1}+\e^{-x}.
\]
Then
\[
 \e^T=7x^2\e^{x^2-3x}
 \exp(x^{-1}+\e^{-x}).
\]
To a few orders,
\begin{align*}
 \e^T
 =7x^2\e^{x^2-3x}\biggl(&1+x^{-1}+\frac{x^{-2}}2
 +\e^{-x}+x^{-1}\e^{-x}
 +\frac{\e^{-2x}}2+\cdots\biggr).
\end{align*}
The very large part determines one monomial; only the small part is Taylor-expanded.
\end{workedbox}

\section{Why we do not Taylor-expand a large exponent}
It would be meaningless at infinity to write
\[
  \e^x=1+x+\frac{x^2}{2!}+\cdots
\]
as a transseries ordered from large to small: its support runs upward forever and has no dominant term. Instead, $\e^x$ is treated as a single monomial. Taylor series are used only for infinitesimal arguments.

This principle is worth memorizing:
\begin{center}
\emph{large pieces create monomials; small pieces create Taylor expansions.}
\end{center}

\section{Logarithms by multiplicative normal form}
For $T>0$, write
\[
  T=a\m(1+\varepsilon),
  \qquad a>0,
  \quad \varepsilon\prec1.
\]
Then
\[
  \log T=\log a+\log\m+
  \log(1+\varepsilon),
\]
where
\[
  \log(1+\varepsilon)
  =\varepsilon-\frac{\varepsilon^2}{2}
  +\frac{\varepsilon^3}{3}-\cdots.
\]

\begin{workedbox}{taking a logarithm}
Let
\[
  T=3x^5\e^{\sqrt{x}}
  \left(1+\frac2x-\e^{-x}\right).
\]
Set $u=2x^{-1}-\e^{-x}$. Then
\begin{align*}
  \log T
  ={}&\log3+5\log x+\sqrt{x}
       +u-\frac{u^2}{2}+\frac{u^3}{3}-\cdots\\
  ={}&\sqrt{x}+5\log x+\log3
       +2x^{-1}-2x^{-2}+\frac{8}{3}x^{-3}\\
     &{}-\e^{-x}+2x^{-1}\e^{-x}
       -4x^{-2}\e^{-x}-\frac12\e^{-2x}+\cdots.
\end{align*}
\end{workedbox}

\section{Real powers and binomial series}
For $T>0$ and $\alpha\in\R$, define
\[
  T^\alpha=\exp(\alpha\log T).
\]
Equivalently,
\[
  \bigl(a\m(1+\varepsilon)\bigr)^\alpha
  =a^\alpha\m^\alpha
   \sum_{n\ge0}\binom{\alpha}{n}\varepsilon^n.
\]
This includes square roots, reciprocals, and arbitrary real powers.

\begin{workedbox}{a square root across two scales}
For
\[
  T=x^2\left(1+\frac1x+\e^{-x}\right),
\]
we obtain
\begin{align*}
  \sqrt{T}
  =x\biggl(&1+\frac{1}{2x}-\frac{1}{8x^2}
  +\frac{1}{16x^3}
  +\frac12\e^{-x}-\frac{1}{4x}\e^{-x}\\
  &{}+\frac{3}{16x^2}\e^{-x}
  -\frac18\e^{-2x}+\cdots\biggr).
\end{align*}
\end{workedbox}

\section{Identities and their domains}
The familiar identities
\[
  \log(AB)=\log A+\log B,
  \qquad
  \e^{A+B}=\e^A\e^B
\]
hold in the positive real transseries field. Complex logarithms require branch choices, exactly as in complex analysis. Even over the reals, $\log A$ is defined only for $A>0$.

\begin{warningbox}
The formal Taylor series for $\log(1+S)$ requires $S\prec1$, not merely $S<1$ in some informal sense. If $S$ contains a large monomial, factor the dominant monomial first.
\end{warningbox}

\chapter{Differentiation}
\label{chap:differentiation}

\section{Differentiate term by term}
The derivative on ordinary monomials extends recursively:
\[
  (x^a)'=ax^{a-1},
  \qquad
  (\e^L)'=L'\e^L,
  \qquad
  (\log x)'=x^{-1}.
\]
For a transseries
\[
  T=\sum_{\m}a_{\m}\m,
\]
define
\[
  T'=\sum_{\m}a_{\m}\m'.
\]
A nontrivial support theorem shows that this derivative family is summable. The resulting operation obeys the usual rules:
\[
  (A+B)'=A'+B',
  \quad
  (AB)'=A'B+AB',
  \quad
  (\e^A)'=A'\e^A,
  \quad
  (\log A)'=\frac{A'}A.
\]

\section{Logarithmic derivatives}
For a nonzero monomial $\m$, its logarithmic derivative is
\[
  \m^\dagger:=\frac{\m'}{\m}=(\log\m)'.
\]
This often contains much less information than $\m$ itself and is easier to compare. Examples are
\begin{align*}
  (x^a)^\dagger&=\frac{a}{x},\\
  (\e^{x^2})^\dagger&=2x,\\
  \left(x^a(\log x)^b\e^{-x}\right)^\dagger
  &=-1+\frac{a}{x}+\frac{b}{x\log x}.
\end{align*}
Thus
\[
  \left(x^a(\log x)^b\e^{-x}\right)'
  =x^a(\log x)^b\e^{-x}
   \left(-1+\frac{a}{x}+\frac{b}{x\log x}\right).
\]

\section{Differentiation can change the comparison scale}
If $A\prec B$, it is tempting to assume $A'\prec B'$. This is often true in tame settings but is not a purely formal consequence of $A\prec B$. For instance,
\[
  A=1,
  \qquad B=x
\]
satisfies $A\prec B$, yet $A'=0$ and $B'=1$. A more dramatic issue in arbitrary function spaces is rapid oscillation, although the standard ordered transseries field excludes such behavior.

The H-field axiom gives the right replacement: if $F$ is larger than every constant, then $F'>0$. In particular, large positive transseries are eventually increasing in their germ interpretation.

\section{Constants of differentiation}
In the standard real transseries field,
\[
  T'=0\quad\Longrightarrow\quad T\in\R.
\]
This is not automatic for every differential extension, but it is a defining strength of the standard construction. Thus integration constants behave exactly as expected.

\begin{workedbox}{differentiating a two-sector transseries}
Let
\[
  T=x^3+\log x+\e^{-x}x^2
  \left(1-\frac4x-\frac{8}{x^3}\right).
\]
For the exponential sector write $T_1=\e^{-x}x^2H$, where
$H=1-4x^{-1}-8x^{-3}$. Then
\[
  T_1'=\e^{-x}x^2
  \left(H'+\left(-1+\frac2x\right)H\right).
\]
Hence
\begin{align*}
  T'={}&3x^2+x^{-1}
  +\e^{-x}x^2\left(
    -1+\frac6x-\frac4{x^2}
    +\frac8{x^3}-\frac8{x^4}+\frac{24}{x^5}
  \right).
\end{align*}
The logarithmic-derivative method avoids repeatedly applying the product rule.
\end{workedbox}

\section{A proof of the chain rule in the small-perturbation case}
Let $F$ be represented by a formal series around $G$ and let $U\prec G$ be small enough that
\[
  F(G+U)=\sum_{n\ge0}\frac{F^{(n)}(G)}{n!}U^n
\]
is summable. Differentiating term by term gives
\begin{align*}
  \frac{\dd}{\dd x}F(G+U)
  &=(G'+U')
    \sum_{n\ge0}\frac{F^{(n+1)}(G)}{n!}U^n\\
  &=(G'+U')F'(G+U).
\end{align*}
The general transseries chain rule follows by the recursive construction of composition.

\chapter{Integration and asymptotic antidifferentiation}
\label{chap:integration}

\section{The simplest cases}
For $a\ne-1$,
\[
  \int x^a\dd x=\frac{x^{a+1}}{a+1},
\]
while
\[
  \int x^{-1}\dd x=\log x.
\]
For a rapidly varying exponential, integration by parts naturally produces a transseries. If $L'\ne0$, seek
\[
  \int \e^L A\dd x=\e^L B.
\]
Differentiation yields the first-order equation
\[
  B'+L'B=A.
\]
If $B'/L'B$ is small, solve recursively:
\[
  B=\frac{A}{L'}-\frac{1}{L'}
      \left(\frac{A}{L'}\right)'
    +\frac{1}{L'}\left[
       \frac{1}{L'}\left(\frac{A}{L'}\right)'
     \right]'-\cdots.
\]
This is repeated integration by parts written as an operator expansion.

\begin{workedbox}{the large-$x$ expansion of an exponential integral}
Take $L=x^2$ and $A=1$. Then $L'=2x$, so
\begin{align*}
 \int^x \e^{t^2}\dd t
 &\asympseries \frac{\e^{x^2}}{2x}
 \left(1+\frac{1}{2x^2}
 +\frac{3}{(2x^2)^2}
 +\frac{15}{(2x^2)^3}+\cdots\right)\\
 &=\frac{\e^{x^2}}{2x}
 \left(1+\frac{1}{2x^2}
 +\frac{3}{4x^4}
 +\frac{15}{8x^6}+\cdots\right).
\end{align*}
A derivative check determines the recurrence. If
$B=\sum_{n\ge0}b_nx^{-2n-1}$, the equation $B'+2xB=1$ gives
\[
  b_0=\frac12,
  \qquad b_{n+1}=\frac{2n+1}{2}b_n.
\]
\end{workedbox}

\section{Asymptotic integration of a monomial}
Let $\m=\e^L$. We want a transseries $F$ with $F'=\m$. The leading guess is
\[
  F_0=\frac{\m}{L'}.
\]
Indeed,
\[
  F_0'=\m\left(1-\frac{L''}{(L')^2}\right).
\]
When $L''/(L')^2\prec1$, the relative error is small and may be corrected iteratively. This covers many familiar cases, such as $L=x^p$ with $p>0$.

There is also a ``slow'' regime. For example, integrating $x^a$ is not best viewed as dividing by $(a\log x)'=a/x$ when $a=-1$, because a resonance occurs and the answer moves to the new logarithmic scale $\log x$.

\section{Exact closure is a deep theorem}
The preceding algorithms produce formal asymptotic antiderivatives. A major theorem says more: the standard field of logarithmic-exponential transseries is closed under integration. Thus for every $A\in\Tfield$ there is $B\in\Tfield$ with $B'=A$.

This statement is not proved by naively integrating every term independently. Derivatives can mix neighboring monomials, and resonant logarithms may be required. The proof uses the structure of asymptotic couples and is one of the gateways from elementary transseries calculation to differential algebra \cite{Hoeven2006,ADH2017}.

\section{Formal definite integrals require endpoints}
An antiderivative is a transseries near infinity. A numerical definite integral such as
\[
  \int_0^\infty f(x)\dd x
\]
involves behavior at both $0$ and $\infty$ and generally cannot be recovered from a transseries at infinity alone. Two functions can have the same complete asymptotic expansion at infinity but different integrals because they differ on a compact interval.

\begin{warningbox}
A transseries is local asymptotic data. It may determine a germ near infinity without determining a global function or a global integral.
\end{warningbox}

\chapter{Composition and Taylor expansion}
\label{chap:composition}

\section{When composition makes asymptotic sense}
If $F$ is a transseries at $+\infty$, then $F\circ G$ should be defined when the inner transseries $G$ is positive and infinitely large:
\[
  G>r\qquad\text{for every }r\in\R.
\]
This mirrors ordinary asymptotics: a large-$x$ expansion for $F(x)$ can be substituted into only when $G(x)\to+\infty$.

For monomials, composition begins with
\[
  x^a\circ G=G^a,
  \qquad
  \e^L\circ G=\exp(L\circ G),
  \qquad
  \log x\circ G=\log G.
\]
It then extends to summable series.

\section{Taylor's formula is the workhorse}
Suppose
\[
  G=G_0+U,
  \qquad U\prec G_0,
\]
and the derivatives of $F$ form a summable Taylor family. Then
\[
  F(G_0+U)=
  \sum_{n\ge0}\frac{F^{(n)}(G_0)}{n!}U^n.
\]
Unlike numerical Taylor series, this identity is controlled by support rather than by a radius of convergence.

\begin{workedbox}{composing an exponential with a small correction}
Let
\[
  F(x)=\e^{-x},
  \qquad G=x+\log x+x^{-1}.
\]
Then directly,
\[
  F\circ G
  =\e^{-x}x^{-1}\e^{-x^{-1}}
  =\e^{-x}x^{-1}
   \left(1-x^{-1}+\frac{x^{-2}}2-\frac{x^{-3}}6+\cdots\right).
\]
The inner logarithm changes the algebraic prefactor; the infinitesimal $x^{-1}$ is Taylor-expanded.
\end{workedbox}

\begin{workedbox}{Taylor composition around a large base}
Let $F(x)=\log x$ and
\[
  G=x^2+x+\e^{-x}=x^2(1+u),
  \qquad u=x^{-1}+x^{-2}\e^{-x}.
\]
Then
\begin{align*}
  \log G
  &=2\log x+\log(1+u)\\
  &=2\log x+x^{-1}-\frac{x^{-2}}2
    +\frac{x^{-3}}3+ x^{-2}\e^{-x}
    -x^{-3}\e^{-x}+\cdots.
\end{align*}
\end{workedbox}

\section{Associativity and the chain rule}
On its natural domain of positive infinite inner arguments, transseries composition satisfies
\[
  (F\circ G)\circ H=F\circ(G\circ H)
\]
and
\[
  (F\circ G)'=(F'\circ G)G'.
\]
The proof is not just a symbolic slogan: one must show that every family created by substitution is summable and that the result is independent of the chosen construction stage. Edgar's work gives a particularly approachable treatment in the well-based setting \cite{EdgarComposition}.

\section{Composition is not substitution into an arbitrary small quantity}
The expression $F(1/x)$ cannot in general be obtained from a transseries expansion of $F$ at $+\infty$, because $1/x\to0$. One needs a separate expansion of $F$ near zero. Similarly, substituting a negative infinite transseries may leave the real logarithmic-exponential domain.

\begin{checkpointbox}
For composition, always ask two questions:
\begin{enumerate}
\item Toward which endpoint is the outer expansion valid?
\item Does the inner transseries approach that endpoint in the required sector or direction?
\end{enumerate}
Formal syntax does not erase analytic domain information.
\end{checkpointbox}

\chapter{Compositional inversion and implicit equations}
\label{chap:inverse}

\section{The inverse problem}
Given a positive infinite transseries $F$, we seek $G$ such that
\[
  F\circ G=x,
  \qquad
  G\circ F=x.
\]
At a practical level, the process resembles series reversion:
\begin{enumerate}
\item invert the dominant behavior;
\item substitute the approximation;
\item measure the error;
\item add a smaller correction;
\item repeat.
\end{enumerate}

\section{Newton correction for composition}
Suppose $G$ is an approximate inverse and
\[
  E=F\circ G-x
\]
is small compared with $x$. Replace $G$ by $G+\Delta$. Taylor expansion gives
\[
  F(G+\Delta)-x
  =E+(F'\circ G)\Delta
   +\frac12(F''\circ G)\Delta^2+\cdots.
\]
The first Newton correction is therefore
\[
  \Delta=-\frac{E}{F'\circ G}.
\]
Under the transseries ordering, the new error is usually quadratically smaller.

\begin{workedbox}{inverting $x+\log x$}
Let
\[
  F(x)=x+\log x.
\]
The dominant inverse is $G_0=x$. Since
\[
  F(G_0)-x=\log x,
\]
we try $G_1=x-\log x$. Then
\begin{align*}
 F(G_1)-x
 &= -\log x+\log(x-\log x)\\
 &=\log\left(1-\frac{\log x}{x}\right)\\
 &=-\frac{\log x}{x}
   -\frac{(\log x)^2}{2x^2}-\cdots.
\end{align*}
The next correction begins with $(\log x)/x$. Continuing gives
\[
  F^{[-1]}(x)
  =x-\log x+\frac{\log x}{x}
   +\frac{(\log x)^2-2\log x}{2x^2}
   +\cdots.
\]
Direct substitution verifies the displayed orders.
\end{workedbox}

\section{A structural theorem}
Let
\[
  \mathcal P=\set{F\in\Tfield:F>\R}
\]
be the set of positive infinite transseries. With composition, $\mathcal P$ forms a group: every $F\in\mathcal P$ has a compositional inverse in $\mathcal P$. This is a deep closure property, not a consequence of ordinary field inversion.

\section{Implicit equations by dominant balance}
Consider
\[
  P(x,y)=0
\]
with transseries coefficients. A robust strategy is:
\begin{enumerate}
\item guess the dominant monomial of $y$ by balancing the largest terms;
\item write $y=y_0(1+u)$ or $y=y_0+u$ with $u$ smaller;
\item expand the equation in $u$;
\item solve recursively for the next term.
\end{enumerate}

\begin{workedbox}{an algebraic equation with exponentially small roots}
Solve
\[
  y^2-y=\e^{-x}.
\]
The exact roots are
\[
  y_{\pm}=\frac{1\pm\sqrt{1+4\e^{-x}}}{2}.
\]
Using the binomial series,
\begin{align*}
  y_+&=1+\e^{-x}-\e^{-2x}
      +2\e^{-3x}-5\e^{-4x}+14\e^{-5x}-\cdots,\\
  y_-&=-\e^{-x}+\e^{-2x}
      -2\e^{-3x}+5\e^{-4x}-14\e^{-5x}+\cdots.
\end{align*}
The coefficients are Catalan numbers with alternating signs. The two branches arise from two different dominant balances, $y\sim1$ and $y\sim0$.
\end{workedbox}

\section{Implicit differentiation and transseries branches}
If $P_y(x,y_0)$ is nonzero and dominates the nonlinear corrections, a transseries version of the implicit-function method constructs a branch near $y_0$. The formal Newton step is
\[
  y_{n+1}=y_n-\frac{P(x,y_n)}{P_y(x,y_n)}.
\]
Each step removes the current dominant residual. This ``cancel the largest error'' principle will recur in every differential-equation example.

\chapter{How to compute a finite transseries in practice}
\label{chap:algorithm}

\section{The cancellation algorithm}
Most hand calculations follow the same loop.
\begin{enumerate}
\item Choose a scale or ratio set, such as
\[
  \boldsymbol\mu=(x^{-1},\e^{-x}).
\]
\item Guess the dominant term from the equation.
\item Substitute the partial transseries.
\item Find the largest uncancelled residual monomial.
\item Add one term whose image under the linearized operator cancels it.
\item Stop when the residual is below the requested cutoff.
\end{enumerate}
This is long division, Newton's method, and asymptotic matching rolled into one.

\section{A coefficient-matching template}
Suppose an equation suggests
\[
  y=\sum_{n=0}^{\infty}a_nx^{-n}.
\]
Substitute and rewrite every expression in the same monomial basis. If the result is
\[
  \sum_{n\ge0}R_n(a_0,\ldots,a_n)x^{-n}=0,
\]
then solve
\[
  R_n=0
\]
sequentially. Triangularity is the reason this works: the coefficient at a new scale usually depends linearly on the newest unknown and polynomially on earlier ones.

\section{Several sectors}
If the homogeneous linearization contains an exponentially small solution $\e^{-Ax}x^\beta$, enlarge the ansatz to
\[
  y=\sum_{k\ge0}C^k\e^{-kAx}x^{k\beta}
    \sum_{n\ge0}a_{k,n}x^{-n}.
\]
The integer $k$ labels the \defterm{nonperturbative sector}. Substitution often yields a two-dimensional triangular recurrence:
\begin{itemize}
\item solve sector $k=0$ first;
\item solve $k=1$ using the linearized homogeneous equation;
\item nonlinear products of lower sectors force $k=2,3,\ldots$.
\end{itemize}

\section{Bookkeeping with weights}
Assign a weight to each generator. For example,
\[
  w(x^{-1})=1,
  \qquad
  w(\e^{-x})=M,
\]
where $M$ is chosen larger than any algebraic order being retained. Then sort monomials by the pair
\[
  (k,n)	ext{ for }\e^{-kx}x^{-n},
\]
first by sector $k$, then by power $n$, or according to the exact asymptotic order needed. This is a practical computational device, not a replacement for the true monomial ordering.

\section{A pseudocode skeleton}
\begin{verbatim}
input: equation E[y] = 0, ansatz generators, cutoff
T := dominant approximation
repeat:
    R := expand E[T] down to cutoff
    if R has no retained term: return T
    r*m := dominant term of R
    L := linearization of E about T
    choose correction c*n so dominant(L[c*n]) = -r*m
    T := T + c*n
end
\end{verbatim}
For nonlinear equations, recomputing the linearization after each correction is safest. For a fixed power-series ansatz, direct coefficient recurrences are faster.

\section{Verification by residual}
A finite transseries approximation should never be trusted merely because its coefficients look plausible. Substitute it back into the original equation and display the leading residual. If
\[
  E[T_N]=\bigO(\n),
\]
then the approximation is certified through all monomials strictly larger than $\n$, provided the linearization is nondegenerate on the omitted scale.

\begin{checkpointbox}
A good transseries calculation has three visible parts:
\begin{enumerate}
\item the proposed support or sectors;
\item the recurrence or cancellation rule;
\item a residual check at the first omitted order.
\end{enumerate}
The applications in the next part follow this pattern repeatedly.
\end{checkpointbox}

\part{Worked examples and applications}

\chapter{A gallery of elementary transseries}
\label{chap:gallery}

The best way to make the notation natural is to see many short examples. In each case, the main task is to factor out the dominant scale and expand only an infinitesimal remainder.

\section{Hyperbolic functions}
Since
\[
  \sinh x=\frac{\e^x}{2}(1-\e^{-2x}),
\]
we have the exact transseries identity
\[
  \log(\sinh x)
  =x-\log2+\log(1-\e^{-2x})
  =x-\log2-\sum_{n\ge1}\frac{\e^{-2nx}}{n}.
\]
Thus
\[
  \log(\sinh x)
  =x-\log2-\e^{-2x}-\frac12\e^{-4x}
   -\frac13\e^{-6x}-\cdots.
\]
Here the series converges for $x>0$, but it is also a perfectly good transseries at infinity.

Similarly,
\begin{align*}
  \tanh x
  &=\frac{1-\e^{-2x}}{1+\e^{-2x}}\\
  &=1-2\e^{-2x}+2\e^{-4x}
    -2\e^{-6x}+\cdots.
\end{align*}

\section{A reciprocal with exponential nesting}
Factor the dominant term:
\[
  \frac{1}{\e^x+x}
  =\frac{\e^{-x}}{1+x\e^{-x}}.
\]
Because $x\e^{-x}\prec1$,
\[
  \frac{1}{\e^x+x}
  =\e^{-x}-x\e^{-2x}+x^2\e^{-3x}
   -x^3\e^{-4x}+\cdots.
\]
The powers of $x$ grow within successive exponential sectors, but the whole sequence still decreases because every additional factor $\e^{-x}$ overwhelms any fixed power of $x$.

\section{Nested logarithms}
For
\[
  F(x)=\log(x+\log x),
\]
factor $x$:
\begin{align*}
  F(x)
  &=\log x+\log\left(1+\frac{\log x}{x}\right)\\
  &=\log x+\frac{\log x}{x}
    -\frac{(\log x)^2}{2x^2}
    +\frac{(\log x)^3}{3x^3}-\cdots.
\end{align*}
This example already lies beyond ordinary Laurent series because the coefficients contain powers of $\log x$.

\section{An exponential of an asymptotic series}
Let
\[
  A=x+\frac1x+\frac{2}{x^2}.
\]
Then
\begin{align*}
  \e^A
  &=\e^x\exp\left(x^{-1}+2x^{-2}\right)\\
  &=\e^x\left(
      1+x^{-1}+\frac{5}{2}x^{-2}
      +\frac{13}{6}x^{-3}
      +\frac{73}{24}x^{-4}+\cdots
    \right).
\end{align*}
A large exponential factor and a small power-series factor are kept conceptually separate.

\section{A square root near a large branch}
For fixed $a$,
\begin{align*}
 \sqrt{x^2+a}
 &=x\sqrt{1+ax^{-2}}\\
 &=x+\frac{a}{2x}-\frac{a^2}{8x^3}
   +\frac{a^3}{16x^5}
   -\frac{5a^4}{128x^7}+\cdots.
\end{align*}
If instead the radicand contains a beyond-all-orders term,
\[
  \sqrt{x^2+\e^{-x}}
  =x+\frac{\e^{-x}}{2x}
   -\frac{\e^{-2x}}{8x^3}
   +\frac{\e^{-3x}}{16x^5}-\cdots.
\]

\section{An inverse function close to a logarithm}
Let $F(t)=\e^t+t$ and write its inverse as $t=G(x)$. Since $\e^t$ dominates $t$, the first approximation is $G\sim\log x$. Put
\[
  G=\log x+u,
  \qquad u\prec1.
\]
The equation $\e^G+G=x$ becomes
\[
  x\e^u+\log x+u=x,
\]
or
\[
  \e^u=1-\frac{\log x+u}{x}.
\]
Taking logarithms and solving recursively gives
\begin{align*}
  G(x)=\log x
  &-\frac{\log x}{x}
  +\frac{\log x-\tfrac12(\log x)^2}{x^2}\\
  &+\frac{-\log x+\tfrac32(\log x)^2
    -\tfrac13(\log x)^3}{x^3}+\cdots.
\end{align*}
A residual substitution is an efficient way to verify each displayed coefficient.

\section{What these examples teach}
All six computations use the same small number of moves:
\begin{enumerate}
\item identify and factor the dominant monomial;
\item turn the remaining factor into $1+\varepsilon$ with $\varepsilon\prec1$;
\item apply a familiar Taylor series formally;
\item reorder the answer by asymptotic magnitude;
\item check the first omitted term by substitution.
\end{enumerate}

\chapter{Lambert's \texorpdfstring{$W$}{W} function, step by step}
\label{chap:lambertw}

\section{\texorpdfstring{The inverse of $w\mapsto w\e^w$}{The inverse of w mapped to w exp(w)}}
Lambert's function $W$ is defined by
\[
  W(x)\e^{W(x)}=x.
\]
For $x\to+\infty$, taking logarithms gives
\[
  W+\log W=\log x.
\]
The first term must therefore be close to $\log x$, but substituting $W\sim\log x$ leaves an error $\log\log x$. This forces a second scale.

Set
\[
  L_1=\log x,
  \qquad
  L_2=\log\log x.
\]
The first two terms are
\[
  W\sim L_1-L_2.
\]
This is the simplest famous example in which iterated logarithms arise naturally from inversion.

\section{Deriving the correction equation}
Write
\[
  W=L_1-L_2+u,
  \qquad u\prec L_2.
\]
Then
\begin{align*}
  \log W
  &=\log L_1+
    \log\left(1-\frac{L_2-u}{L_1}\right)\\
  &=L_2+
    \log\left(1-\frac{L_2-u}{L_1}\right).
\end{align*}
The equation $W+\log W=L_1$ reduces exactly to
\[
  u+\log\left(1-\frac{L_2-u}{L_1}\right)=0.
\]
Now $L_2/L_1\prec1$, so the logarithm may be expanded.

\section{Coefficient matching}
Seek
\[
  u=\frac{P_1(L_2)}{L_1}
    +\frac{P_2(L_2)}{L_1^2}
    +\frac{P_3(L_2)}{L_1^3}+\cdots,
\]
where $P_n$ is a polynomial. Substitution gives successively
\begin{align*}
 P_1(z)&=z,\\
 P_2(z)&=\frac{z^2}{2}-z,\\
 P_3(z)&=\frac{z^3}{3}-\frac32z^2+z,\\
 P_4(z)&=\frac{z^4}{4}-\frac{11}{6}z^3+3z^2-z.
\end{align*}
Therefore
\begin{align}
 W(x)={}&L_1-L_2+\frac{L_2}{L_1}
 +\frac{L_2^2/2-L_2}{L_1^2}\notag\\
 &+\frac{L_2^3/3-(3/2)L_2^2+L_2}{L_1^3}\notag\\
 &+\frac{L_2^4/4-(11/6)L_2^3+3L_2^2-L_2}{L_1^4}
 +\cdots.
 \label{eq:Wexpansion}
\end{align}

\begin{workedbox}{a numerical sanity check}
At $x=10^6$,
\[
  L_1\approx13.8155,
  \qquad L_2\approx2.6258.
\]
The successive approximations are roughly
\[
  11.1897,
  \qquad 11.3798,
  \qquad 11.3833,
  \qquad 11.38336,
\]
while the true principal value is about $11.38336$. The point is not the decimal itself but the rapid gain obtained by respecting the correct logarithmic scale.
\end{workedbox}

\section{A compact all-orders formula}
The polynomials can be expressed using unsigned Stirling numbers of the first kind. One standard form is
\[
 W=L_1-L_2+
 \sum_{\ell\ge0}\sum_{m\ge1}
 \frac{(-1)^\ell}{m!}
 \genfrac{[}{]}{0pt}{}{\ell+m}{\ell+1}
 \frac{L_2^m}{L_1^{\ell+m}}.
\]
Here $\genfrac{[}{]}{0pt}{}{n}{k}$ counts permutations of $n$ objects with $k$ cycles. The double sum is ordered by total inverse-logarithmic degree. A classical analytic account of the function and its branches is \cite{Corless1996}.

\section{The lower real branch}
For $-1/\e<x<0$, the equation has two real branches. As $x\to0^-$, the branch $W_{-1}(x)$ tends to $-\infty$. The same logarithmic mechanism applies, but branch choices must be tracked:
\[
  L_1=\log(-x)<0,
  \qquad
  L_2=\log(-L_1).
\]
Then
\[
  W_{-1}(x)
  =L_1-L_2+\frac{L_2}{L_1}+\cdots.
\]
The algebraic pattern is the same as \cref{eq:Wexpansion}; the analytic branch is different.

\section{Why this example matters}
Lambert $W$ displays three central features of transseries:
\begin{enumerate}
\item compositional inversion creates new scales;
\item logarithms iterate because the correction to a logarithm is another logarithm;
\item coefficient functions are polynomials in slower variables such as $\log\log x$.
\end{enumerate}
It is a model for much more complicated inverse problems.

\chapter{A linear differential equation with a hidden exponential}
\label{chap:linearODE}

\section{The equation}
Consider
\[
  y''+xy'+y=0.
  \label{eq:linearhidden}
\]
The left side factors as
\[
  \frac{\dd}{\dd x}(y'+xy),
\]
so one integration gives
\[
  y'+xy=C_2.
\]
Using the integrating factor $\e^{x^2/2}$,
\[
  y(x)=C_1\e^{-x^2/2}
  +C_2\e^{-x^2/2}\int^x\e^{t^2/2}\dd t.
\]

\section{The algebraic sector}
Seek a formal inverse-power solution
\[
  y_0=\sum_{n\ge0}a_nx^{-2n-1}.
\]
Substitution into $y'+xy=1$ gives
\[
  a_0=1,
  \qquad
  a_{n+1}=(2n+1)a_n.
\]
Hence
\[
  y_0=\frac1x+\frac1{x^3}+\frac3{x^5}
      +\frac{15}{x^7}+\frac{105}{x^9}+\cdots,
\]
with
\[
  a_n=(2n-1)!!=\frac{(2n)!}{2^nn!}.
\]
This series diverges factorially but is asymptotic to the second exact solution.

\section{The invisible homogeneous solution}
Every function
\[
  y=y_0+C\e^{-x^2/2}
\]
has the same inverse-power asymptotic expansion, because
\[
  \e^{-x^2/2}\prec x^{-N}
  \qquad\text{for every }N.
\]
The full formal information is therefore the one-parameter transseries
\[
  y\sim
  \left(\frac1x+\frac1{x^3}+\frac3{x^5}+\cdots\right)
  +C\e^{-x^2/2}.
\]
In this linear example there are no higher powers of $C\e^{-x^2/2}$ because superposition is linear.

\begin{ideabox}
An ordinary asymptotic series describes a quotient space of solutions: it identifies functions that differ by something beyond all retained orders. A transseries restores the missing homogeneous coordinate $C$.
\end{ideabox}

\section{Divergence predicts the hidden scale}
The $n$th algebraic term is approximately
\[
  \frac{(2n)!}{2^nn!x^{2n+1}}.
\]
Stirling's formula shows that the smallest term occurs near $n\approx x^2/2$ and has size comparable, up to an algebraic factor, to $\e^{-x^2/2}$. Thus the late terms of the algebraic sector know the exponential scale of the missing homogeneous solution.

\section{Residual verification}
Truncate after $N$ terms:
\[
  y_N=\sum_{n=0}^{N-1}(2n-1)!!x^{-2n-1}.
\]
A direct calculation gives
\[
  y_N'+xy_N=1-(2N-1)!!x^{-2N}.
\]
Differentiating once more verifies \cref{eq:linearhidden} up to the corresponding first omitted order. This exact residual formula makes the asymptotic status transparent.

\chapter{A nonlinear ODE and its tower of sectors}
\label{chap:nonlinearODE}

\section{The model problem}
Consider the Riccati-type equation
\[
  y'+y=\frac1x+y^2.
  \label{eq:riccati-model}
\]
It is simple enough for hand computation but rich enough to generate a genuine nonlinear transseries.

\section{The perturbative sector}
Assume
\[
  y_0=\sum_{m\ge1}a_mx^{-m}.
\]
Matching $x^{-1}$ gives $a_1=1$. For $m\ge2$,
\[
  a_m=(m-1)a_{m-1}
  +\sum_{j=1}^{m-1}a_ja_{m-j}.
  \label{eq:riccati-recurrence}
\]
The first coefficients are
\[
  a_1,a_2,a_3,a_4,a_5,a_6,a_7
  =1,2,8,44,296,2312,20384.
\]
Thus
\[
  y_0=\frac1x+\frac2{x^2}+\frac8{x^3}
      +\frac{44}{x^4}+\frac{296}{x^5}+\cdots.
\]
The rapid coefficient growth warns us that this is an asymptotic, not convergent, series.

\section{Linearization reveals the first exponential}
Put
\[
  y=y_0+\delta.
\]
Using the equation for $y_0$ gives
\[
  \delta'+(1-2y_0)\delta=\delta^2.
  \label{eq:delta-equation}
\]
At first order, discard $\delta^2$:
\[
  \delta_1'+(1-2y_0)\delta_1=0.
\]
Hence
\[
  \delta_1=C\exp\left(-x+2\int y_0\dd x\right).
\]
Since $y_0=x^{-1}+2x^{-2}+\cdots$,
\[
  \delta_1=C\e^{-x}x^2h_1(x),
\]
where
\[
  h_1=1-\frac4x-\frac8{x^3}
      -\frac{52}{x^4}-\frac{384}{x^5}
      -\frac{3200}{x^6}-\cdots.
\]
The power $x^2$ comes from integrating the term $2/x$ in $2y_0$.

\section{Nonlinearity forces higher sectors}
Because the right side of \cref{eq:delta-equation} contains $\delta^2$, a first sector proportional to $C\e^{-x}$ produces a second sector proportional to $C^2\e^{-2x}$, which then interacts to produce all higher sectors. The natural ansatz is
\[
  y=y_0+
  \sum_{k\ge1}C^k\e^{-kx}x^{2k}h_k(x),
  \qquad
  h_k(x)=\sum_{n\ge0}c_{k,n}x^{-n}.
  \label{eq:riccati-transseries}
\]
For $k=1$, the equation is homogeneous. For $k\ge2$,
\[
 h_k'+\left(1-k+\frac{2k}{x}-2y_0\right)h_k
 =\sum_{i=1}^{k-1}h_i h_{k-i}.
 \label{eq:sector-recurrence}
\]
For example,
\[
  h_2=-1+\frac6x-\frac6{x^2}
      +\frac8{x^3}+\frac{48}{x^4}
      +\frac{384}{x^5}+\cdots.
\]

\section{Where the free parameter enters}
The perturbative recurrence \cref{eq:riccati-recurrence} fixes every $a_m$. The first exponential equation is homogeneous, so its overall coefficient $C$ is free. Every later sector is forced polynomially by the earlier ones. This is the typical one-parameter pattern for a first-order nonlinear ODE.

\begin{workedbox}{checking the first sector}
Let
\[
  q=\e^{-x}x^2,
  \qquad h_1=1-4x^{-1}+\bigO(x^{-2}).
\]
Then
\[
  \frac{q'}q=-1+\frac2x.
\]
Since $1-2y_0=1-2x^{-1}-4x^{-2}+\cdots$, the coefficient of $qh_1$ in the linearized left side is
\[
  \left(-1+\frac2x\right)+(1-2y_0)
  =-4x^{-2}+\cdots.
\]
The derivative of $h_1$ begins with $4x^{-2}$, so the leading residual cancels. This explains the coefficient $-4/x$ without solving the entire recurrence.
\end{workedbox}

\section{Formal solution versus actual solution}
The formal family \cref{eq:riccati-transseries} does not by itself say which value of $C$ corresponds to a particular boundary condition, nor whether the sectors can be summed along a chosen complex direction. Those are analytic questions. Borel summation and Stokes phenomena, developed in Part~V, explain how the formal parameter can jump when the direction of summation changes.

\chapter{WKB analysis and the Airy equation}
\label{chap:wkb}

\section{Why exponentials appear in second-order equations}
Consider
\[
  \varepsilon^2 y''(x)=Q(x)y(x),
  \qquad 0<\varepsilon\ll1.
  \label{eq:wkb-general}
\]
If we try $y=\exp(S/\varepsilon)$, then
\[
  \frac{y''}{y}=\frac{(S')^2}{\varepsilon^2}
  +\frac{S''}{\varepsilon}.
\]
Equation \cref{eq:wkb-general} becomes
\[
  (S')^2+\varepsilon S''=Q.
\]
The leading balance is
\[
  S_0'=\pm\sqrt{Q},
\]
so one expects exponential scales
\[
  \exp\left(\pm\frac1\varepsilon
  \int^x\sqrt{Q(t)}\dd t\right).
\]
This is the origin of the two WKB sectors.

\section{Logarithmic derivative and the Riccati equation}
It is often cleaner to set
\[
  u=\frac{y'}{y}.
\]
Then
\[
  y''/y=u'+u^2.
\]
For the Airy equation
\[
  y''=xy,
  \label{eq:airy}
\]
we obtain
\[
  u'+u^2=x.
  \label{eq:airy-riccati}
\]
The dominant balance $u^2\sim x$ gives two branches
\[
  u\sim\pm x^{1/2}.
\]

\section{Solving the Riccati expansion}
Write
\[
  u_\pm=\pm x^{1/2}+v.
\]
Substitution into \cref{eq:airy-riccati} and matching descending powers of $x^{3/2}$ yields
\begin{align*}
 u_+={}&x^{1/2}-\frac1{4x}
 -\frac5{32x^{5/2}}-\frac{15}{64x^4}
 -\frac{1105}{2048x^{11/2}}+\cdots,\\
 u_-={}&-x^{1/2}-\frac1{4x}
 +\frac5{32x^{5/2}}-\frac{15}{64x^4}
 +\frac{1105}{2048x^{11/2}}+\cdots.
\end{align*}
The exponents advance by $3/2$ because differentiating $x^{1/2}$ lowers the exponent by one, while multiplication by the leading $x^{1/2}$ shifts corrections by one half.

\section{Integrating back to the Airy solutions}
Since $u=(\log y)'$, integrate the decaying branch:
\begin{align*}
  \log y_-={}&-\frac23x^{3/2}-\frac14\log x
  -\frac5{48x^{3/2}}+\frac5{64x^3}+\cdots.
\end{align*}
Exponentiating gives
\[
 y_-(x)=C_-x^{-1/4}\e^{-2x^{3/2}/3}
 \left(1-\frac5{48x^{3/2}}
 +\frac{385}{4608x^3}
 -\frac{85085}{663552x^{9/2}}+\cdots\right).
 \label{eq:airy-decay}
\]
With the normalization $C_-=1/(2\sqrt\pi)$, this is the classical expansion of $\operatorname{Ai}(x)$ on the positive real axis. The growing branch has the scale
\[
  x^{-1/4}\e^{+2x^{3/2}/3}.
\]

\section{Turning points}
A point where $Q(x)=0$ is a \defterm{turning point}. The WKB monomials involve $Q^{-1/4}$ and therefore become singular there. Near a simple turning point, rescaling reduces the differential equation to the Airy equation. Thus Airy transseries are not an isolated curiosity; they are the universal local model for a wide class of turning-point problems.

\section{Dominance depends on complex direction}
For complex $x=re^{i\theta}$,
\[
  x^{3/2}=r^{3/2}e^{3i\theta/2}.
\]
The sign of
\[
  \Re\left(\frac23x^{3/2}\right)
\]
determines which WKB exponential grows and which decays. Along rays where this real part vanishes, the two exponentials have equal magnitude. These rays are tied to Stokes and anti-Stokes geometry. The formal sectors stay the same, but their analytic importance changes with direction. A classical rigorous treatment of these asymptotic constructions is given by Olver \cite{Olver1997}.

\begin{ideabox}
WKB analysis can be summarized as follows: turn a differential equation into an equation for a logarithmic derivative, find the dominant algebraic branches, integrate them to create exponential monomials, and then solve recursively for the prefactor series.
\end{ideabox}

\section{Newton polygons in one paragraph}
For more complicated differential equations, one may plot how candidate terms scale under an ansatz $y\asymp\e^{S}x^\alpha$. Balancing terms that lie on a common edge of the resulting Newton polygon identifies possible dominant exponents. The polygon does not compute the full transseries, but it organizes which balances should be tried and which exponential scales can occur.

\chapter{Singular perturbations and boundary layers}
\label{chap:singular}

\section{Regular versus singular perturbation}
A parameter $\varepsilon$ is \defterm{regular} if setting $\varepsilon=0$ leaves an equation of the same essential type. It is \defterm{singular} if the highest derivative disappears, the order drops, or boundary conditions can no longer all be imposed. Singular perturbations routinely produce scales such as
\[
  \e^{-A/\varepsilon},
\]
which are invisible to every power of $\varepsilon$.

\section{A first-order boundary layer revisited}
The equation
\[
  \varepsilon y'+y=1,
  \qquad y(0)=0,
\]
has solution
\[
  y=1-\e^{-x/\varepsilon}.
\]
The outer expansion at fixed $x>0$ is simply $y\sim1$. To resolve the initial condition, introduce the inner variable
\[
  X=\frac{x}{\varepsilon}.
\]
Then the equation becomes
\[
  \frac{\dd y}{\dd X}+y=1,
\]
whose solution $1-\e^{-X}$ is not small when $X=\bigO(1)$. The same exponential is therefore negligible in the outer region and essential in the boundary layer.

\section{A two-point example}
Consider
\[
  \varepsilon y''+y'=0,
  \qquad y(0)=0,
  \quad y(1)=1.
  \label{eq:boundary-two-point}
\]
The exact solution is
\[
  y(x)=\frac{1-\e^{-x/\varepsilon}}
  {1-\e^{-1/\varepsilon}}.
\]
Expanding the denominator geometrically gives the transseries
\begin{align*}
 y(x)={}&(1-\e^{-x/\varepsilon})
 \left(1+\e^{-1/\varepsilon}
 +\e^{-2/\varepsilon}+\cdots\right)\\
 ={}&1-\e^{-x/\varepsilon}
 +\e^{-1/\varepsilon}
 -\e^{-(x+1)/\varepsilon}
 +\e^{-2/\varepsilon}-\cdots.
\end{align*}
Several exponential actions occur: $x$, $1$, $x+1$, and so on. Which term is largest depends on the location $x$.

\section{Outer, inner, and composite descriptions}
A standard matched-asymptotic calculation has three levels:
\begin{enumerate}
\item the \defterm{outer expansion}, valid away from the thin layer;
\item the \defterm{inner expansion}, written in a stretched variable such as $X=x/\varepsilon$;
\item a \defterm{composite approximation}, obtained by adding outer and inner approximations and subtracting their common overlap.
\end{enumerate}
Transseries offer a fourth viewpoint: retain both the algebraic and exponential scales in one symbolic object. This does not eliminate the need to state regions of validity, but it makes the hidden scales explicit.

\section{Multiple boundary layers}
If an equation has layers near several endpoints, one may encounter
\[
  \e^{-x/\varepsilon},
  \qquad
  \e^{-(1-x)/\varepsilon},
  \qquad
  \e^{-1/\varepsilon},
\]
and products of these monomials. The product
\[
  \e^{-x/\varepsilon}\e^{-(1-x)/\varepsilon}
  =\e^{-1/\varepsilon}
\]
shows that local layer sectors can interact to create a global exponentially small scale.

\section{Nonuniformity is not a contradiction}
For fixed $x>0$,
\[
  \e^{-x/\varepsilon}\prec\varepsilon^N
  \quad\text{for all }N,
\]
but this estimate is not uniform down to $x=0$. When $x=\varepsilon X$, the term becomes $\e^{-X}$, which is order one. An asymptotic statement is incomplete unless its limiting process and domain are specified.

\begin{warningbox}
Never say merely ``$\e^{-x/\varepsilon}$ is negligible.'' Say where: it is beyond all algebraic orders for fixed $x>0$ as $\varepsilon\to0^+$, but not inside the $x=\bigO(\varepsilon)$ boundary layer.
\end{warningbox}

\chapter{Difference equations and Stirling's series}
\label{chap:difference}

\section{The gamma function as a difference equation}
The defining recurrence
\[
  \Gamma(x+1)=x\Gamma(x)
\]
becomes, after taking logarithms,
\[
  F(x+1)-F(x)=\log x,
  \qquad F=\log\Gamma.
  \label{eq:gamma-difference}
\]
This is the difference-equation analogue of antidifferentiation.

\section{The shift operator}
Let $D=\dd/\dd x$. Formally,
\[
  F(x+1)=\e^D F(x),
\]
so \cref{eq:gamma-difference} is
\[
  (\e^D-1)F=\log x.
\]
Multiplying by $D/(\e^D-1)$ gives an equation for $F'$:
\[
  F'=\frac{D}{\e^D-1}\log x.
\]
The generating function for Bernoulli numbers is
\[
  \frac{z}{\e^z-1}
  =\sum_{n\ge0}B_n\frac{z^n}{n!}.
\]
Therefore
\[
  F'=\sum_{n\ge0}\frac{B_n}{n!}D^n\log x.
\]

\section{Integrating the formal expansion}
Using $B_1=-1/2$ and $B_{2n+1}=0$ for $n\ge1$, one obtains
\begin{align*}
 F'(x)\asympseries{}&\log x-\frac{1}{2x}
 -\frac{1}{12x^2}+\frac{1}{120x^4}
 -\frac{1}{252x^6}+\cdots.
\end{align*}
Integration yields
\begin{align}
 \log\Gamma(x)\asympseries{}&
 \left(x-\frac12\right)\log x-x+C
 +\sum_{n\ge1}
 \frac{B_{2n}}{2n(2n-1)x^{2n-1}}.
 \label{eq:stirling}
\end{align}
The difference equation does not determine the constant $C$. The standard normalization of $\Gamma$ gives
\[
  C=\frac12\log(2\pi).
\]

\section{Why Stirling's series diverges}
The Bernoulli numbers satisfy
\[
  |B_{2n}|\sim
  \frac{2(2n)!}{(2\pi)^{2n}}.
\]
Consequently, the terms of \cref{eq:stirling} eventually grow for every fixed $x$. The least term occurs at roughly
\[
  n\approx\pi x,
\]
and its scale is approximately
\[
  \e^{-2\pi x}
\]
times an algebraic factor. Once again, factorial divergence predicts a beyond-all-orders exponential.

\section{The periodic ambiguity}
If $F$ solves \cref{eq:gamma-difference}, then so does
\[
  F+P
\]
for every period-one function $P(x+1)=P(x)$. In the ordered real LE-transseries field, nonconstant oscillatory periodic functions are absent, so only a constant ambiguity remains. In broader complex frameworks, Fourier modes and exponentially small Stokes data encode related phenomena.

\section{Euler--Maclaurin as a transseries machine}
The same operator identity underlies the Euler--Maclaurin formula. For a smooth $f$,
\begin{align*}
 \sum_{k=a}^{b}f(k)
 ={}&\int_a^b f(t)\dd t+\frac{f(a)+f(b)}2\\
 &+\sum_{n=1}^{N-1}\frac{B_{2n}}{(2n)!}
 \left(f^{(2n-1)}(b)-f^{(2n-1)}(a)\right)+R_N.
\end{align*}
When the derivative series diverges, optimal truncation and exponentially small corrections enter. Thus transseries are as natural for difference equations and sums as they are for differential equations.

\chapter{Fractional iteration and Abel's equation}
\label{chap:iteration}

\section{From integer iterates to a continuous parameter}
For a function or transseries $T$, write
\[
  T^{[0]}(x)=x,
  \qquad
  T^{[1]}=T,
  \qquad
  T^{[n+m]}=T^{[n]}\circ T^{[m]}.
\]
A \defterm{fractional iterate} is a family $T^{[s]}$ with real or complex $s$ satisfying
\[
  T^{[s+t]}=T^{[s]}\circ T^{[t]}.
\]
The central device is to convert composition into addition.

\section{Abel's equation}
Seek an \defterm{Abel coordinate} $A$ satisfying
\[
  A(T(x))=A(x)+1.
  \label{eq:abel}
\]
Then define
\[
  T^{[s]}(x)=A^{[-1]}(A(x)+s).
  \label{eq:fractional-abel}
\]
Indeed,
\begin{align*}
 T^{[s]}(T^{[t]}(x))
 &=A^{[-1]}\bigl(A(T^{[t]}(x))+s\bigr)\\
 &=A^{[-1]}(A(x)+t+s)
 =T^{[s+t]}(x).
\end{align*}

\section{The trivial example}
For $T(x)=x+1$, take $A(x)=x$. Then
\[
  T^{[s]}(x)=x+s.
\]
The importance of Abel's equation is that the same idea works when no elementary coordinate is available.

\section{\texorpdfstring{A nontrivial example: $T(x)=x+1/x$}{A nontrivial example: T(x)=x+1/x}}
The map
\[
  T(x)=x+\frac1x
\]
increases $x$ only slightly. Since
\[
  T(x)^2-x^2=2+x^{-2},
\]
the leading Abel coordinate should be $x^2/2$. A correction by $\log x$ removes the next error. Solving \cref{eq:abel} term by term gives
\begin{align}
 A(x)={}&\frac{x^2}{2}-\frac12\log x
 +\frac{1}{8x^2}+\frac{5}{96x^4}
 +\frac{7}{288x^6}-\frac{1}{1280x^8}+\cdots.
 \label{eq:abel-example}
\end{align}
Substitution of $x+x^{-1}$ into this expression verifies
\[
  A(T(x))-A(x)=1
\]
to the displayed order.

\section{The fractional iterate}
Using \cref{eq:fractional-abel} and reverting the series gives
\begin{align}
 T^{[s]}(x)={}&x+\frac{s}{x}
 -\frac{s(s-1)}{2x^3}
 +\frac{s(s-1)^2}{2x^5}\notag\\
 &-\frac{s(s-1)(3s-2)(5s-7)}{24x^7}
 +\bigO(x^{-9}).
 \label{eq:fractional-example}
\end{align}
Several checks are immediate:
\[
  T^{[0]}(x)=x,
  \qquad
  T^{[1]}(x)=x+x^{-1}.
\]
For $s=2$, \cref{eq:fractional-example} gives
\[
  x+\frac2x-\frac1{x^3}+\frac1{x^5}-\frac1{x^7}+\cdots,
\]
which agrees with directly expanding $T(T(x))$.

\section{The iterative logarithm}
Differentiating $T^{[s]}$ with respect to $s$ at $s=0$ defines the \defterm{iterative logarithm} or infinitesimal generator
\[
  V(x)=\left.\frac{\partial}{\partial s}T^{[s]}(x)\right|_{s=0}.
\]
The flow equation is
\[
  \frac{\partial}{\partial s}T^{[s]}(x)
  =V(T^{[s]}(x)).
\]
From \cref{eq:fractional-example},
\[
  V(x)=\frac1x+\frac1{2x^3}
  +\frac1{2x^5}+\frac7{12x^7}+\cdots.
\]
This is the composition-theoretic analogue of the logarithm of a matrix.

\section{Existence and uniqueness subtleties}
Abel coordinates are not automatically unique as functions: if $P$ is period one, then $A+P\circ A$ may produce another coordinate under suitable conditions. In formal transseries classes, normalization and support restrictions often restore uniqueness. Edgar proved substantial existence results for large positive transseries of exponentiality zero and developed their fractional iterates within the transseries framework \cite{EdgarFractional}.

\begin{advancedbox}
Iteration near a fixed point is usually linearized by Schr\"oder's equation rather than Abel's equation. Iteration near infinity, where $T(x)-x$ is small or of controlled growth, naturally leads to Abel coordinates. Tetration and continuous iteration of the exponential require still richer asymptotic structures and delicate choices of domain.
\end{advancedbox}

\part{From formal expressions to analytic functions}

\chapter{Formal transseries versus analytic functions}
\label{chap:formal-analytic}

\section{A formal object is not automatically a function}
The expression
\[
  \widetilde F(x)=\sum_{n\ge0}n!x^{-n-1}
\]
is a valid formal power-series sector, but it diverges for every finite $x$. It therefore cannot be interpreted by ordinary pointwise summation. A formal transseries records ordered asymptotic data; an analytic function is an object with numerical values on a domain. Connecting the two requires an additional theorem or summation procedure.

\section{Asymptotic representation}
Let
\[
  \widetilde F=\sum_{j\ge0}a_j\m_j,
  \qquad
  \m_0\succ\m_1\succ\cdots.
\]
An analytic function $F$ is represented by this transseries on a domain if, for each retained cutoff,
\[
  F-\sum_{j<N}a_j\m_j
  =\smallo(\m_{N-1})
\]
uniformly in the specified limit or sector. For multiple exponential sectors, one must state how the transseries parameter is chosen and in which complex direction the comparison holds.

\section{The asymptotic map is not injective}
At infinity,
\[
  0
  \qquad\text{and}\qquad
  \e^{-x}
\]
have the same expansion in inverse powers of $x$. More generally, adding any flat function smaller than all members of the selected scale leaves that restricted asymptotic expansion unchanged. Consequently, a Poincar\'e series alone usually cannot identify a unique function.

A fuller transseries distinguishes many flat corrections, but even a formal transseries may admit several analytic sums in different directions. Stokes phenomena measure this remaining ambiguity.

\section{Sectors in the complex plane}
For $x\in\C$, a monomial such as $\e^{-Ax}$ is small only where
\[
  \Re(Ax)>0.
\]
Thus analytic realization is usually sectorial. A series may be summable in one angular sector, fail to be summable on a boundary ray, and acquire a different exponentially small contribution after crossing that ray.

\section{Three levels of assertion}
It is helpful to distinguish the following statements:
\begin{enumerate}
\item \emph{Formal:} substitution into an equation cancels every monomial.
\item \emph{Asymptotic:} finite truncations approximate an actual function with controlled remainders.
\item \emph{Summability:} a specified operation reconstructs a canonical analytic function from the formal series in a chosen direction.
\end{enumerate}
A formal calculation can be correct even when summability has not been proved. Conversely, numerical agreement does not replace a formal residual calculation.

\section{Why summation should respect algebra and calculus}
An ideal summation map $\mathscr S$ should satisfy, whenever both sides make sense,
\begin{align*}
  \mathscr S(A+B)&=\mathscr S(A)+\mathscr S(B),\\
  \mathscr S(AB)&=\mathscr S(A)\mathscr S(B),\\
  \mathscr S(A')&=(\mathscr S A)',\\
  \mathscr S(F\circ G)&=(\mathscr S F)\circ(\mathscr S G).
\end{align*}
No single elementary method covers every transseries, but Borel--Laplace summation achieves much of this program for Gevrey and resurgent classes.

\chapter{Borel transformation and Laplace summation}
\label{chap:borel}

\section{\texorpdfstring{Why factorial divergence suggests dividing by $n!$}{Why factorial divergence suggests dividing by n factorial}}
Consider a formal inverse-power series
\[
  \widetilde\phi(x)=\sum_{n=0}^{\infty}a_nx^{-n-1}.
\]
If $a_n$ grows like $n!$, the ordinary series diverges. The \defterm{formal Borel transform} removes this factorial growth:
\[
  \Borel\widetilde\phi(p)
  =\widehat\phi(p)
  :=\sum_{n=0}^{\infty}\frac{a_n}{n!}p^n.
  \label{eq:borel-transform}
\]
The new series often has a nonzero radius of convergence.

The inverse operation is suggested by
\[
  \int_0^\infty\e^{-xp}p^n\dd p
  =\frac{n!}{x^{n+1}},
  \qquad \Re x>0.
\]
Thus the directional Laplace transform
\[
  \Laplace_\theta\widehat\phi(x)
  =\int_0^{\e^{i\theta}\infty}
    \e^{-xp}\widehat\phi(p)\dd p
  \label{eq:laplace-transform}
\]
formally reconstructs the original powers.

\section{The Borel--Laplace recipe}
To sum $\widetilde\phi$ in direction $\theta$:
\begin{enumerate}
\item form the local Borel series \cref{eq:borel-transform};
\item analytically continue it along the ray $p\in\e^{i\theta}\R_{\ge0}$;
\item check that its growth is at most exponential along the ray;
\item apply the Laplace integral \cref{eq:laplace-transform}.
\end{enumerate}
When all steps work, define
\[
  \mathscr S_\theta\widetilde\phi
  :=\Laplace_\theta\Borel\widetilde\phi.
\]

\section{Euler's alternating factorial series}
Take
\[
  \widetilde F_-(x)=
  \sum_{n\ge0}(-1)^n n!x^{-n-1}.
\]
Its Borel transform is the convergent geometric series
\[
  \widehat F_-(p)=\sum_{n\ge0}(-p)^n=\frac1{1+p}.
\]
There is no singularity on the positive $p$-axis, so
\[
  \mathscr S_0\widetilde F_-(x)
  =\int_0^\infty\frac{\e^{-xp}}{1+p}\dd p
  =-\e^x\Ei(-x),
  \qquad x>0.
\]
The Borel sum is exactly the function introduced in \cref{chap:formal-analytic}'s motivating examples.

\section{Watson's lemma: why the Laplace integral has the right expansion}
Suppose $\widehat\phi$ is analytic near $p=0$ with
\[
  \widehat\phi(p)=\sum_{n=0}^{N-1}\frac{a_n}{n!}p^n+R_N(p),
  \qquad R_N(p)=\bigO(p^N).
\]
Insert this into the Laplace integral. The polynomial part gives
\[
  \sum_{n=0}^{N-1}a_nx^{-n-1}.
\]
On a short interval, the remainder contributes $\bigO(x^{-N-1})$ after rescaling $p=u/x$. On the remaining interval, the exponential kernel makes the contribution smaller still, under an exponential-growth bound. Therefore
\[
  \Laplace_0\widehat\phi(x)
  \asympseries\sum_{n\ge0}a_nx^{-n-1}.
\]
This elementary argument is the core of Watson's lemma.

\section{Gevrey order}
A formal series is \defterm{Gevrey-1} if, roughly,
\[
  |a_n|\le CA^n n!
\]
for some $C,A>0$. Then its Borel transform converges near the origin. More generally, growth like $\Gamma(1+sn)$ leads to Gevrey order $s$ and may require generalized Borel transforms. Most elementary resurgent examples are Gevrey-1 after a suitable change of variable.

\section{Products become convolution}
If
\[
  \widetilde f=\sum a_nx^{-n-1},
  \qquad
  \widetilde g=\sum b_nx^{-n-1},
\]
then, up to conventions concerning constant terms,
\[
  \Borel(\widetilde f\widetilde g)
  =\widehat f*\widehat g,
\]
where
\[
  (\widehat f*\widehat g)(p)
  =\int_0^p\widehat f(q)\widehat g(p-q)\dd q.
\]
The Laplace transform turns convolution back into multiplication. This is why Borel--Laplace summation can respect nonlinear equations rather than merely sum isolated series.

\section{\texorpdfstring{Derivatives and multiplication by $x$}{Derivatives and multiplication by x}}
With the convention \cref{eq:borel-transform}, differentiation and multiplication correspond to elementary Borel-plane operations. For example,
\[
  \Borel(\widetilde\phi')=-p\widehat\phi(p),
\]
provided no polynomial or constant boundary term is present. Thus differential equations become integral or convolution equations in the Borel plane, where analytic continuation can often be studied directly.

\chapter{Stokes phenomena: exponentially small jumps}
\label{chap:stokes}

\section{A Borel singularity on the integration ray}
Now consider the nonalternating factorial series
\[
  \widetilde F_+(x)=\sum_{n\ge0}n!x^{-n-1}.
\]
Its Borel transform is
\[
  \widehat F_+(p)=\frac1{1-p},
\]
which has a pole at $p=1$ on the positive integration ray. The ordinary Laplace integral is ambiguous. We define two \defterm{lateral sums}, one just above and one just below the ray:
\[
  \mathscr S_{0^\pm}\widetilde F_+
  =\int_0^{\infty\e^{\pm i0}}
   \frac{\e^{-xp}}{1-p}\dd p.
\]

\section{Computing the jump}
The difference of the two contours is a small loop around $p=1$. With the orientation convention above,
\[
  \mathscr S_{0^+}\widetilde F_+
  -\mathscr S_{0^-}\widetilde F_+
  =2\pi i\e^{-x}.
  \label{eq:euler-stokes-jump}
\]
The sign reverses if the lateral-orientation convention is reversed, but the scale and magnitude are fixed.

This formula is the simplest complete Stokes phenomenon:
\begin{itemize}
\item the two sums have the same power-series expansion;
\item they differ by an exponential smaller than every power;
\item the exponential scale is determined by the Borel singularity location $p=1$.
\end{itemize}

\section{Why the singularity location becomes an action}
A singularity at $p=A$ contributes, after Laplace transformation, a factor
\[
  \e^{-Ax}.
\]
Indeed, shifting $p=A+q$ gives
\[
  \e^{-xp}=\e^{-Ax}\e^{-xq}.
\]
The number $A$ is often called an \defterm{action}, especially in WKB theory and physics. A constellation of Borel singularities therefore predicts a constellation of exponential transmonomials.

\section{The transseries parameter absorbs the jump}
Suppose a formal solution has the form
\[
  \widetilde y=\widetilde\Phi_0
  +C\e^{-Ax}\widetilde\Phi_1+\cdots.
\]
If the lateral sum of $\widetilde\Phi_0$ jumps by
\[
  S\e^{-Ax}\mathscr S\widetilde\Phi_1,
\]
then the represented analytic solution remains unchanged when $C$ jumps oppositely:
\[
  C_+=C_- -S.
\]
The constant $S$ is a \defterm{Stokes constant}. The formal parameter is not a flaw; it is the coordinate that makes sectorial analytic continuation possible.

\section{Stokes lines and anti-Stokes lines}
Terminology varies between communities. A useful geometric distinction is:
\begin{itemize}
\item a ray where a Borel singularity obstructs directional summation is a Stokes direction in the Borel--Laplace sense;
\item a ray where two exponentials have equal magnitude is often called an anti-Stokes line, although some authors exchange the names.
\end{itemize}
Always check the convention in a source. The invariant facts are which exponential is dominant, where the Laplace contour meets singularities, and what jump occurs.

\section{A jump can be invisible to all perturbative orders}
If
\[
  F_+-F_-=2\pi i\e^{-x},
\]
then for every $N$,
\[
  F_+-F_-=\smallo(x^{-N}).
\]
Thus no finite or infinite inverse-power expansion can detect the jump. The transseries sector $\e^{-x}$ is exactly the missing resolution.

\section{Berry smoothing}
An idealized Stokes jump looks discontinuous when one compares asymptotic expansions on the two sides of a ray. The underlying analytic function is usually smooth. Near optimal truncation, the effective multiplier of the small exponential often changes rapidly but smoothly according to an error-function profile. This phenomenon, discovered and developed by Berry and others, explains how a discontinuous asymptotic bookkeeping rule emerges from a smooth function.

\chapter{Resurgence and alien calculus, gently}
\label{chap:resurgence}

\section{The central resurgent idea}
A divergent series may become an analytic germ after Borel transformation. Continue that germ through the Borel plane. If its singularities are isolated and their local expansions are themselves related to other sectors of the original problem, the series is called \defterm{resurgent} in an appropriate technical sense.

The slogan is:
\begin{quote}
The different exponential sectors are not independent; each sector reappears in the singularity structure of the others.
\end{quote}

\section{A toy large-order relation}
Let
\[
  \widehat\phi(p)=\frac1{1-p/A}.
\]
Expanding at $p=0$ gives
\[
  \widehat\phi(p)=\sum_{n\ge0}\frac{p^n}{A^n},
\]
so the original formal coefficients are
\[
  a_n=\frac{n!}{A^n}.
\]
The Borel singularity at $A$ and the factorial-over-power growth of $a_n$ are two views of the same fact. More complicated singularities produce refinements such as
\[
  a_n\sim
  \frac{S}{2\pi i}
  \frac{\Gamma(n+\beta)}{A^{n+\beta}}
  \left(b_0+\frac{A b_1}{n+\beta-1}+\cdots\right).
  \label{eq:large-order-template}
\]
The coefficients $b_j$ typically belong to the neighboring exponential sector.

\section{Reading another sector from late coefficients}
Suppose the perturbative coefficients $a_n$ are known to high order. One can:
\begin{enumerate}
\item estimate $A$ from ratios such as $n a_{n-1}/a_n$;
\item remove the leading factorial-over-power growth;
\item extrapolate the remaining sequence to estimate $S b_0$;
\item examine $1/n$ corrections to recover $b_1,b_2,\ldots$.
\end{enumerate}
This is one reason large-order numerical data can reveal nonperturbative physics.

\section{Alien derivatives as singularity detectors}
Ordinary differentiation measures change with respect to $x$. An \defterm{alien derivative} $\Delta_A$ is designed to measure the singular behavior at Borel location $A$. Informally,
\[
  \Delta_A(\text{sector at action }0)
  \propto
  \text{sector at action }A.
\]
Alien derivatives satisfy algebraic rules and can be assembled into the Stokes automorphism. Their exact normalization depends on conventions and on whether one uses pointed alien derivatives, logarithms of Stokes maps, or related variants.

\begin{advancedbox}
In a one-parameter problem, a bridge equation often has the schematic form
\[
  \dot\Delta_A\widetilde y(x,C)
  =S_A(C)\frac{\partial}{\partial C}
   \widetilde y(x,C).
\]
The left side acts in the Borel-singularity direction; the right side acts on the transseries parameter. The equation ``bridges'' alien calculus and ordinary parameter calculus. Its precise polynomial $S_A(C)$ is problem dependent.
\end{advancedbox}

\section{Multi-instanton and logarithmic sectors}
A simple nonlinear problem may have sectors
\[
  \e^{-kAx}\Phi_k(x),
  \qquad k=0,1,2,\ldots.
\]
Resonances can force additional powers of $\log x$:
\[
  \e^{-kAx}x^\beta(\log x)^j\Phi_{k,j}(x).
\]
Several independent actions $A_1,\ldots,A_r$ lead to a lattice
\[
  \exp\bigl(-(k_1A_1+\cdots+k_rA_r)x\bigr).
\]
This resembles a grid-based transseries, but analytic resurgence imposes extra relations between the grid sectors.

\section{Not every transseries is resurgent}
The formal LE field is enormous. Resurgent transseries form special analytic subclasses with controlled Borel continuation and singularities. Therefore:
\begin{itemize}
\item ``transseries'' does not automatically mean ``resurgent'';
\item an LE-transseries may be perfectly meaningful algebraically without a chosen Borel sum;
\item a resurgent expansion may use complex actions and logarithmic sectors not most naturally viewed inside the simplest ordered real LE field.
\end{itemize}
The two theories overlap powerfully but should not be identified.

\chapter{Optimal truncation, hyperasymptotics, and exponential improvement}
\label{chap:hyperasymptotics}

\section{Why more terms eventually hurt}
Suppose
\[
  a_n\sim C A^{-n}\Gamma(n+\beta).
\]
The magnitude of the $n$th term $a_nx^{-n}$ has ratio approximately
\[
  \frac{|a_{n+1}|x^{-n-1}}{|a_n|x^{-n}}
  \sim\frac{n+\beta}{|A|x}.
\]
Terms decrease until $n\approx|A|x$ and grow afterward. Truncating near the least term is \defterm{optimal truncation}.

\section{Estimating the least term}
Stirling's formula gives, at $n\approx Ax$ with $A>0$,
\[
  \Gamma(n+\beta)A^{-n}x^{-n}
  \asymp x^{\beta-1/2}\e^{-Ax}
\]
times a constant factor. Thus the optimally truncated error has exactly the scale predicted by the nearest Borel singularity.

\section{First-level hyperasymptotics}
Ordinary asymptotics expands a function around the Borel origin. Hyperasymptotics proceeds further:
\begin{enumerate}
\item truncate the first series near its least term;
\item express the remainder as a Laplace integral localized near the nearest Borel singularity;
\item expand the singularity contribution in its own asymptotic series;
\item truncate that new series optimally;
\item continue recursively to more distant singularities.
\end{enumerate}
Each level resolves a smaller exponential scale.

\section{Stirling's series as an example}
For \cref{eq:stirling}, Bernoulli growth places the nearest relevant Borel singularities at distances related to $2\pi$. Truncation near $n\approx\pi x$ leaves an error on the scale $\e^{-2\pi x}$. Exponentially improved versions of Stirling's formula express the remainder in terms of terminant functions whose switching captures the Stokes phenomenon.

\section{Superasymptotics versus hyperasymptotics}
Terminology is useful:
\begin{itemize}
\item \defterm{superasymptotics} means optimally truncating the first divergent series;
\item \defterm{hyperasymptotics} means re-expanding and optimally truncating the exponentially small remainder, possibly repeatedly.
\end{itemize}
A full transseries organizes all these levels in advance. Hyperasymptotics is the numerical strategy for extracting them efficiently; the classical sources include Dingle and Berry--Howls \cite{Dingle1973,BerryHowls1990}.

\section{Practical coefficient acceleration}
When coefficients are known numerically, Richardson extrapolation, Borel--Pad\'e approximation, conformal maps, and sequence transformations can estimate actions and Stokes constants. These are powerful tools, but they do not replace structural analysis: spurious Pad\'e poles and finite-data effects can imitate singularities.

\begin{checkpointbox}
The analytic story now fits into one chain:
\[
 \text{factorial divergence}
 \Longleftrightarrow
 \text{Borel singularity}
 \Longleftrightarrow
 \text{exponential scale}
 \Longleftrightarrow
 \text{Stokes jump}.
\]
Resurgence enriches this chain by identifying the local data at one singularity with the coefficients of another transseries sector.
\end{checkpointbox}

\chapter{Saddle points and transseries of integrals}
\label{chap:saddles}

\section{Laplace-type integrals}
Consider
\[
  I(x)=\int_\Gamma \e^{-xV(t)}a(t)\dd t,
  \qquad x\to+\infty.
  \label{eq:laplace-integral}
\]
The dominant contributions come from points where $\Re V$ is smallest along the contour: endpoints, poles, or critical points satisfying
\[
  V'(t_*)=0.
\]
Near a nondegenerate saddle,
\[
  V(t)=V(t_*)+\frac12V''(t_*)(t-t_*)^2+\cdots.
\]
After the scaling $t-t_*=u/\sqrt{x}$, one obtains a sector
\[
  I_{t_*}(x)
  \sim\e^{-xV(t_*)}x^{-1/2}
  \sum_{n\ge0}c_nx^{-n}.
  \label{eq:saddle-sector}
\]
Each saddle supplies its own exponential monomial.

\section{Several saddles produce a transseries}
If saddles $t_1,\ldots,t_r$ can contribute, the full asymptotic object has the schematic form
\[
  I(x)\sim
  \sum_{j=1}^r n_j
  \e^{-xV(t_j)}x^{-1/2}
  \sum_{n\ge0}c_{j,n}x^{-n}.
\]
The coefficients $n_j$ depend on how the original contour decomposes into steepest-descent cycles. As a parameter or direction changes, this decomposition can jump. That jump is the geometric origin of Stokes phenomena for integrals.

\section{Relative actions}
Choose one saddle as a reference. Then
\[
  \e^{-xV(t_j)}
  =\e^{-xV(t_0)}
   \e^{-x(V(t_j)-V(t_0))}.
\]
The differences
\[
  A_j=V(t_j)-V(t_0)
\]
are the actions that appear as Borel singularity locations for the perturbative expansion around $t_0$.

\section{Why the coefficients diverge}
The local expansion around one saddle has a finite region of analytic influence. High derivatives eventually feel neighboring saddles or singularities of $V$ and $a$. This produces factorial growth. Steepest-descent geometry and Borel-plane resurgence are therefore two descriptions of the same global analytic structure.

\section{Endpoint contributions}
If the dominant point is an endpoint rather than a stationary point, the scaling may be $t=u/x$ rather than $u/\sqrt{x}$. This gives powers $x^{-1},x^{-2},\ldots$. Degenerate saddles produce fractional powers, and coalescing saddles require uniform special-function models such as Airy or Pearcey functions.

\section{Infinite-dimensional analogies}
In quantum mechanics and quantum field theory, formal path integrals are expanded around classical solutions or saddles. The resulting perturbative series are commonly divergent, while instantons and other nonperturbative saddles produce factors like
\[
  \e^{-A/g},
\]
where $g$ is a small coupling. Resurgent transseries provide a disciplined language for combining these sectors. Infinite-dimensional applications require substantial analytic care, but the finite-dimensional saddle picture explains the basic grammar.

\part{The wider algebraic landscape}

\chapter{Germs at infinity and Hardy fields}
\label{chap:hardy}

\section{What is a germ?}
Two functions $f,g:(a,\infty)\to\R$ define the same \defterm{germ at infinity} if they agree for all sufficiently large $x$. A germ forgets behavior on bounded intervals and remembers only the eventual function. This is the natural analytic setting for asymptotics.

Addition, multiplication, and differentiation are defined using representatives. For example, changing a function on $[0,100]$ does not affect its germ at infinity.

\section{Hardy fields}
A \defterm{Hardy field} is a field of germs at infinity closed under differentiation. Classical examples are generated by rational functions, exponentials, and logarithms. Because a field cannot contain a nonzero germ with infinitely many zeros, every nonzero element has eventually constant sign. Hence Hardy fields are built for nonoscillatory asymptotics.

\begin{example}
The germs of
\[
  x,\quad \log x,\quad \e^x,
  \quad x^{\sqrt2}(\log x)^{-3},
  \quad \exp\left(\frac{x}{\log x}\right)
\]
can belong to suitable Hardy fields. The germ of $\sin x$ cannot belong to any Hardy field: it has infinitely many zeros, so it would be a nonzero field element without a multiplicative inverse germ.
\end{example}

\section{Orders of infinity become field comparisons}
For positive germs $f,g$, write
\[
  f\prec g\quad\Longleftrightarrow\quad f/g\to0.
\]
In a Hardy field, l'H\^opital-type principles often turn comparisons of functions into comparisons of logarithmic derivatives. This is the analytic ancestor of the valuation and asymptotic-couple structure in transseries.

\section{Hardy's logarithmic-exponential functions}
Starting from real constants and $x$, repeatedly apply field operations, exponentiation, and logarithm wherever defined. The resulting LE-functions form a Hardy field of germs. Examples include
\[
  \frac{\exp(\sqrt{x}+\log\log x)}
  {x^3+\exp(x^{1/4})}.
\]
Hardy studied their orders of infinity long before the modern formal transseries fields were constructed.

\section{From a germ to a formal expansion}
A germ may admit a transseries expansion obtained recursively by dominant comparison. For example,
\[
  \frac1{x+\log x}
  =\frac1x\left(1+\frac{\log x}{x}\right)^{-1}
\]
has expansion
\[
  x^{-1}-x^{-2}\log x+x^{-3}(\log x)^2-\cdots.
\]
The formal field provides a canonical arena in which such expansions can be added, multiplied, differentiated, and composed.

\section{Embedding results}
The field of Hardy's LE-germs embeds naturally into the field of LE-transseries, sending the germ of $x$ to the transseries $x$ and preserving order, field operations, exponentiation, and differentiation. This makes the formal objects faithful representatives for that tame analytic class \cite{DMM1997,DMM2001,ADH2018}.

Much broader Hardy fields can also be embedded into suitable extensions of transseries. Determining how far this relationship extends is an active structural theme: germs emphasize analytic existence, while transseries emphasize canonical formal normal forms.

\section{Why germs and transseries are complementary}
A germ has actual values but may lack a transparent normal form. A transseries has a transparent normal form but may need summation before it represents a function. The theory of H-fields abstracts the common differential-order behavior of both.

\chapter{\texorpdfstring{$H$}{H}-fields and asymptotic differential algebra}
\label{chap:hfields}

\section{The basic axioms}
An \defterm{$H$-field} is an ordered differential field $K$ with constant field
\[
  C=\set{c\in K:c'=0}
\]
that satisfies two asymptotic principles. In a standard formulation:
\begin{enumerate}
\item if $f>C$, meaning $f$ exceeds every constant, then $f'>0$;
\item every bounded element has the form $c+\varepsilon$ with $c\in C$ and $\varepsilon$ infinitesimal.
\end{enumerate}
The second condition says that the valuation ring is $C+\mathfrak o$.

Both Hardy fields and the standard field of transseries are $H$-fields. The axioms capture eventual monotonicity and the existence of a constant ``standard part'' without mentioning functions or formal sums.

\section{Asymptotic couples}
Let $v(f)$ denote the valuation class of a nonzero element. The logarithmic derivative
\[
  f^\dagger=\frac{f'}f
\]
depends asymptotically mainly on $v(f)$. This induces a map on the value group, often denoted
\[
  \psi(v(f))=v(f^\dagger).
\]
The pair consisting of the ordered value group and this derivative map is an \defterm{asymptotic couple}. It encodes the scale-changing behavior of differentiation.

For example,
\[
  (x^a)^\dagger=a/x,
  \qquad
  (\e^{x^2})^\dagger=2x.
\]
Differentiating a power lowers its magnitude, while differentiating a rapid exponential preserves its monomial and multiplies it by a large factor. The asymptotic couple organizes both regimes.

\section{Liouville closure}
A differential field is \defterm{Liouville closed} in the relevant ordered setting when it is real closed and has solutions to the basic first-order equations
\[
  y'=f,
  \qquad
  \frac{z'}z=f\quad(z\ne0).
\]
The second equation amounts to exponentiating an antiderivative. The standard transseries field is Liouville closed. This explains, in structural language, why antiderivatives and exponential integrals remain inside the field.

\section{Newtonianity}
For an algebraic polynomial, Newton's method cancels successive leading errors. Differential polynomials admit an analogous asymptotic Newton analysis. An $H$-field satisfying an appropriate differential Newton property is called \defterm{newtonian}. Informally, differential equations of Newton degree one have roots in the field.

The term should not be confused with numerical convergence of Newton iteration. It is a structural closure property expressed using valuations and differential polynomials.

\section{The model-theoretic picture}
Aschenbrenner, van den Dries, and van der Hoeven developed an axiomatic theory of transseries as $H$-fields. Conditions including Liouville closure, newtonianity, and an ``omega-free'' property isolate the first-order behavior of the standard transseries field. Their work yields strong completeness and model-theoretic transfer results \cite{ADH2017,ADH2018}.

For a beginning reader, the philosophical consequence is more important than the formal logic:
\begin{quote}
Many theorems about transseries depend only on a small set of order, valuation, and differential axioms, not on the particular syntax used to construct the series.
\end{quote}

\section{Differential intermediate values}
Ordinary real closed fields satisfy the intermediate value property for polynomials. Suitable transseries fields satisfy analogues for differential polynomials under sign-change hypotheses. Such results help explain why formal differential equations often have transseries solutions, while also clarifying the hypotheses needed for existence.

\section{Small derivation}
A derivation is called \defterm{small} when it sends infinitesimals or bounded elements into suitably controlled bounded elements. The exact valuation formulation depends on conventions. The familiar derivative on LE-transseries is small after a logarithmic change of independent variable in the settings where the model theory requires it. This technical adjustment aligns differentiation with the valuation topology.

\begin{advancedbox}
The model theory does not merely classify one field. It provides a common language for Hardy fields, transseries, and surreal differential fields, and it suggests universal embedding and extension questions. This is why topics that look like formal asymptotics can lead naturally to valuation theory and logic.
\end{advancedbox}

\chapter{Surreal numbers and transseries}
\label{chap:surreals}

\section{A very brief introduction}
Conway's surreal numbers form a proper class $\mathbf{No}$ containing the real numbers and all ordinal numbers, together with infinitely large and infinitesimal elements. Every surreal has a canonical normal form
\[
  a=\sum_{i<\alpha}r_i\omega^{a_i},
\]
where $\alpha$ is an ordinal, $r_i\in\R\setminus\{0\}$, and the exponents $a_i$ strictly decrease. This looks strikingly like a Hahn series with surreal exponents.

The simplest infinite ordinal $\omega$ plays a role analogous to the transseries variable $x$ evaluated at an infinite value.

\section{Exponential and logarithm}
Gonshor constructed an exponential map
\[
  \exp:\mathbf{No}\to\mathbf{No}_{>0}
\]
that extends the real exponential and has an inverse logarithm. Consequently, expressions such as
\[
  \exp(\omega),
  \qquad
  \log\omega,
  \qquad
  \exp\bigl(\sqrt\omega+\log\log\omega\bigr)
\]
are genuine surreal numbers, not merely formal symbols.

\section{\texorpdfstring{Embedding transseries by sending $x$ to $\omega$}{Embedding transseries by sending x to omega}}
Standard fields of transseries can be embedded into $\mathbf{No}$ by sending
\[
  x\longmapsto\omega
\]
and interpreting sums, exponentials, and logarithms in the surreal sense. A formal transseries then becomes a number of enormous or infinitesimal size. This makes precise the old intuition that a transseries is what an asymptotic expression would evaluate to at an ``infinitely large $x$.''

\section{A Hardy-type derivation on the surreals}
Berarducci and Mantova constructed a derivation
\[
  \partial:\mathbf{No}\to\mathbf{No}
\]
compatible with the field structure and exponential, with constant field $\R$, and satisfying Hardy-type order properties. With this derivation, the surreal numbers are Liouville closed; in particular, the derivation is surjective \cite{BM2018}.

The basic normalization is consistent with
\[
  \partial(\omega)=1,
\]
so the embedded transseries variable behaves like an independent variable.

\section{Omega-series}
Inside $\mathbf{No}$, the \defterm{omega-series} form the smallest subfield containing $\R$ and $\omega$ and closed under exponential, logarithm, and all allowed summable families. They contain the familiar LE- and EL-series but are much larger. Berarducci and Mantova showed that omega-series can be composed and differentiated and can be interpreted as functions on positive infinite surreals \cite{BM2019}.

\section{Why not compose every surreal?}
It is tempting to ask for a composition operation
\[
  a\circ b
\]
for every surreal $a$ and every positive infinite surreal $b$, compatible with the derivation and extending transseries composition. For the simplest surreal derivation, such a composition cannot exist on all of $\mathbf{No}$ with the desired compatibility. Omega-series form a large domain where composition works; the full surreal universe is too broad \cite{BM2019}.

\section{Numbers, functions, and formal expressions meet}
A transseries can be viewed in three ways:
\begin{enumerate}
\item as a formal Hahn-like sum;
\item as an asymptotic germ after analytic realization;
\item as a surreal number after evaluation at an infinite argument.
\end{enumerate}
These viewpoints are not identical, but the modern theory builds rigorous bridges among them.

\begin{warningbox}
A surreal number is not automatically a function. Composition on omega-series supplies a functional interpretation for a large subclass, but it is incorrect to assume that every element of $\mathbf{No}$ can be substituted into every other one.
\end{warningbox}

\chapter{LE-series, EL-series, omega-series, and hyperseries}
\label{chap:variants}

\section{Why several constructions exist}
The phrase ``close under exponentials, logarithms, and infinite sums'' leaves choices about the order and extent of closure. Different choices produce different fields. They agree on a large common core but need not be isomorphic.

\section{Logarithmic-exponential series}
LE-series are built by a controlled hierarchy of finite exponential height and finite logarithmic depth, together with well-based or grid-based summation. This is the classical field used throughout most of this tutorial.

Informally, each individual LE-transseries can be produced from powers $x^r$ by applying field operations, exponentiation, logarithm, and allowed infinite sums through finitely many nested construction stages.

\section{Exponential-logarithmic series}
EL-series arise from a different closure strategy based on prelogarithmic Hahn fields and repeated exponential extensions. Kuhlmann and collaborators developed these constructions in ordered exponential algebra.

Kuhlmann and Tressl proved that the LE-series field embeds as an exponential field into any of the relevant EL-series fields, but no non-Archimedean EL-series field embeds as an exponential field into the LE-series field. Thus EL-series are not merely a different notation for the same object \cite{KuhlmannTressl2011}.

\section{Omega-series}
Omega-series use the enormous summation resources of the surreal numbers. They contain LE- and EL-series and remain closed under composition and differentiation. Their supports and nesting complexity can extend far beyond any fixed finite-height construction.

\section{Why hyperseries are considered}
Even omega-series and familiar LE-series do not exhaust every desirable transfinite logarithmic scale. One may wish to represent iterates such as
\[
  \log_\alpha x
\]
for transfinite indices $\alpha$, or solve functional and differential equations whose natural supports pass through infinitely many logarithmic levels. \defterm{Hyperseries} and \defterm{logarithmic hyperseries} are frameworks designed for such transfinite iteration and summation \cite{HoevenHyperseries}.

\section{A containment picture}
The following diagram is schematic rather than a literal equality of every author's conventions:
\[
  \text{Laurent/Puiseux series}
  \subset
  \text{generalized power series}
  \subset
  \text{LE-transseries}
  \subset
  \text{omega-series}
  \subset
  \mathbf{No}.
\]
EL-series overlap the middle of this picture but are constructed differently and can contain elements not belonging to the LE field.

\section{Grid-based and well-based variants revisited}
At each level, one may impose a finitely generated grid or allow arbitrary reverse well-ordered support. Grid-based fields are more algorithmic; well-based fields have stronger closure and conceptual generality. Statements about composition, convergence, and fixed points can depend on which variant is chosen.

\section{A comparison table}
\begin{center}
\small
\begin{tabularx}{\textwidth}{>{\bfseries}p{0.18\textwidth}X X}
\toprule
Framework & Main strength & Main caution\\
\midrule
Poincar\'e series & Simple coefficient recurrences & Cannot see flat exponentials\\
Hahn series & Very general ordered supports & No exponential or derivation until added\\
LE-transseries & Rich differential, exponential, and compositional field & Finite construction height/depth per element\\
Resurgent transseries & Analytic summation and Stokes data & Special subclass; complex analytic hypotheses required\\
EL-series & Broad exponential closure of Hahn fields & Intrinsically different from LE-series\\
Omega-series & Huge surreal field with composition and derivation & Proper-class scale and deeper set-theoretic machinery\\
Hyperseries & Transfinite logarithmic/exponential iteration & Technically advanced and still developing\\
\bottomrule
\end{tabularx}
\end{center}

\chapter{What transseries do not automatically solve}
\label{chap:limits}

\section{Oscillation}
A totally ordered real transseries field cannot contain a nonzero element that changes sign forever. Functions such as
\[
  \sin x,
  \qquad
  \e^{-x}\sin(\e^x)
\]
require oscillatory extensions, complex phases, Fourier-type monomials, or other structures. One can develop such theories, but the simple dominance order must be modified.

\section{Global information}
An expansion at $x=+\infty$ does not determine values near $x=0$, zeros on a bounded interval, global topology, or a definite integral involving the whole domain. Transseries are local to an asymptotic regime unless supplemented by analytic continuation or global equations.

\section{Uniqueness of analytic realization}
A formal transseries need not have a unique sum. Directional Borel sums can differ by Stokes jumps, and some formal series are not Borel summable. Stronger methods such as multisummation or accelero-summation may be required, and some series fall outside known summability classes.

\section{Arbitrary rearrangement}
Infinite transseries sums cannot be rearranged freely. Summability is a structural condition. A family that is summable in one grouping can become meaningless if a proposed rearrangement destroys well-based support or introduces infinitely many contributions to one coefficient.

\section{Every differential equation}
The transseries field solves many algebraic and differential equations, but not every imaginable equation has a solution in the smallest LE field. New exponential heights, iterated logarithms, hyperseries, oscillatory extensions, or analytic objects may be needed.

\section{Convergence of formal composition}
Formal composition is exact algebraically on its domain, but analytic substitution also requires that the represented functions and sectors overlap. A formal identity can outlive the numerical region in which a particular truncation is useful.

\section{Automatic numerical accuracy}
A transseries with ten displayed terms is not necessarily more accurate than one with five. Divergent sectors should be optimally truncated. Exponentially small sectors can be important near Stokes lines or boundary layers even when they are negligible elsewhere.

\section{A diagnostic checklist}
Before using a transseries answer, ask:
\begin{enumerate}
\item What is the limiting variable and direction?
\item Which monomial ordering is intended?
\item Is the support legal in the chosen framework?
\item Is the statement formal, asymptotic, or summability-based?
\item What is the first omitted scale?
\item Are there Stokes directions or boundary layers?
\item Which constants or transseries parameters are fixed by the problem?
\end{enumerate}

\chapter{A guided reading map}
\label{chap:reading}

\section{For a first conceptual pass}
Begin with Edgar's \emph{Transseries for Beginners}, which develops grid-based and well-based monomials, examples, differentiation, composition, and convergence with unusually concrete explanations \cite{Edgar2009}. The present tutorial was designed to make that entry point still more elementary.

\section{For formal algebra and computation}
Van der Hoeven's monograph \emph{Transseries and Real Differential Algebra} gives a systematic computational and differential-algebraic treatment \cite{Hoeven2006}. The papers of van den Dries, Macintyre, and Marker construct the LE field and establish its model-theoretic and analytic foundations \cite{DMM1997,DMM2001}.

\section{\texorpdfstring{For $H$-fields and model theory}{For H-fields and model theory}}
The monograph by Aschenbrenner, van den Dries, and van der Hoeven is the central modern reference for asymptotic differential algebra and the model theory of transseries \cite{ADH2017}. Their expository article on numbers, germs, and transseries is a shorter orientation \cite{ADH2018}.

\section{For Borel summation and resurgence}
Costin's book develops Borel summability and nonlinear differential equations rigorously \cite{Costin2009}. Dorigoni's introduction and the Physics Reports primer by Aniceto, Ba\c{s}ar, and Schiappa provide broad, calculation-rich introductions to resurgence and physical applications \cite{Dorigoni2019,ABS2019}.

\section{For surreal numbers}
Conway and Gonshor are the classical references for surreal arithmetic and exponentiation \cite{Conway1976,Gonshor1986}. Berarducci and Mantova construct the Hardy-type surreal derivation and then develop omega-series, composition, and surreal analytic functions \cite{BM2018,BM2019}.

\section{For composition and iteration}
Edgar's papers on composition and fractional iteration are direct continuations of the beginner article and include careful support control in well-based and grid-based settings \cite{EdgarComposition,EdgarFractional}.

\section{For the historical roots}
Hahn's generalized series, Hardy's orders of infinity, Borel's summation theory, and \`Ecalle's analyzable functions are four major roots of the subject \cite{Hahn1907,Hardy1910,Borel1899,Ecalle1992}. Modern transseries bring these roots into a common language.

\begin{checkpointbox}
A sensible learning order is:
\[
\begin{aligned}
 \text{asymptotic scales}
 &\longrightarrow \text{Hahn series}
 \longrightarrow \text{LE operations}\\
 &\longrightarrow \text{ODE examples}
 \longrightarrow \text{Borel--Stokes analysis}\\
 &\longrightarrow \text{$H$-fields and surreals}.
\end{aligned}
\]
Trying to begin with alien calculus or model theory before mastering dominant monomials usually makes the subject look harder than it is.
\end{checkpointbox}

\part{Exercises, solutions, and reference material}

\chapter{Exercises}
\label{chap:exercises}

The exercises are grouped by topic. A star indicates a problem that is longer or more conceptual, not one that requires advanced prerequisites. Complete solutions begin in \cref{chap:solutions}.

\section{Asymptotic scales and support}

\begin{exercise}
Arrange the following monomials from largest to smallest as $x\to+\infty$:
\[
 x^{-10},\quad \e^{-\sqrt{x}},\quad
 (\log x)^{100}x^{-3},\quad \e^{-x},\quad
 x^2\e^{-\sqrt{x}},\quad \e^{-x/\log x}.
\]
Justify every neighboring comparison.
\end{exercise}

\begin{exercise}
Show directly from the definition that
\[
  \log x=\smallo(x^a)
\]
for every $a>0$, and that
\[
  x^b=\smallo(\e^{x^a})
\]
for every $a,b>0$.
\end{exercise}

\begin{exercise}
For
\[
  F(x)=\int_0^\infty\frac{\e^{-xp}}{1+p}\dd p,
\]
perform $N$ integrations by parts and prove that
\[
  F(x)\asympseries\sum_{n\ge0}(-1)^nn!x^{-n-1}.
\]
Give an explicit bound for the remainder when $x>0$.
\end{exercise}

\begin{exercise}
Decide which of the following supports are legal at $+\infty$ under the reverse-well-ordered convention:
\begin{align*}
 S_1&=\set{x^{-n}:n\in\N},\\
 S_2&=\set{x^n:n\in\N},\\
 S_3&=\set{\e^{-nx}:n\in\N},\\
 S_4&=\set{x^{-1/n}:n\ge1},\\
 S_5&=\set{x^{-n}\e^{-mx}:n,m\in\N}.
\end{align*}
For each illegal support, exhibit a nonempty subset with no largest element.
\end{exercise}

\begin{exercise}
Let
\[
  A=3x^2-5x+\e^{-x},
  \qquad
  B=3x^2-5x+x^{-7}.
\]
Determine $\magop(A)$, $\domop(A)$, $\magop(A-B)$, and whether $A\sim B$.
\end{exercise}

\begin{exercise}[Hahn-product finiteness]
Let $S,T$ be reverse well-ordered subsets of an ordered abelian group. Explain why a fixed group element $g$ can have only finitely many factorizations $g=st$ with $s\in S$ and $t\in T$. You may argue by contradiction using monotone subsequences.
\end{exercise}

\section{Algebra, exponentials, and logarithms}

\begin{exercise}
Expand
\[
  \frac{1}{x+1+\e^{-x}}
\]
through algebraic order $x^{-4}$ and through the first exponential sector down to $x^{-3}\e^{-x}$.
\end{exercise}

\begin{exercise}
Compute the displayed terms of
\[
  \log\left(2x^3\e^{\sqrt{x}}
  \left(1-\frac1x+2\e^{-x}\right)\right)
\]
through algebraic order $x^{-3}$ and through terms of order $x^{-2}\e^{-x}$.
\end{exercise}

\begin{exercise}
Expand
\[
  \left(x^2+x+\e^{-x}\right)^{3/2}
\]
through algebraic order $x^{-1}$ and include the leading exponential correction.
\end{exercise}

\begin{exercise}
Let
\[
  T=x+2\log x+x^{-1}.
\]
Expand $\e^{-T}$ through order $x^{-4}\e^{-x}$.
\end{exercise}

\begin{exercise}
Show that if $A\succ B>0$, then
\[
  A+B\sim A,
  \qquad
  \log(A+B)=\log A+\frac BA+\bigO((B/A)^2).
\]
State the formal support condition under which the error expansion is meaningful.
\end{exercise}

\section{Differentiation, integration, and composition}

\begin{exercise}
Compute the logarithmic derivative of
\[
  \m=x^a(\log x)^b
  \exp\left(-x+\sqrt{x}\right).
\]
Use it to find the dominant term of $\m'$.
\end{exercise}

\begin{exercise}
Find coefficients $b_0,b_1,b_2,b_3$ such that
\[
  \frac{\dd}{\dd x}\left[
  \e^{x^3}\left(
  b_0x^{-2}+b_1x^{-5}+b_2x^{-8}+b_3x^{-11}
  \right)\right]
  =\e^{x^3}+\bigO(\e^{x^3}x^{-12}).
\]
\end{exercise}

\begin{exercise}
Expand
\[
  \log\left(x^2+x+\e^{-x}\right)
\]
through algebraic order $x^{-4}$ and through the leading exponential term.
\end{exercise}

\begin{exercise}
Let $F(x)=x+\log x$. Derive the first three correction terms in the compositional inverse
\[
  F^{[-1]}(x)=x-\log x+\cdots.
\]
Verify by substitution through order $x^{-2}$.
\end{exercise}

\begin{exercise}
Prove formally that, when $U\prec G$ and the Taylor family is summable,
\[
  \e^{G+U}=\e^G\sum_{n\ge0}\frac{U^n}{n!}.
\]
Explain why the same formula is not used with $U=G=x$ at infinity.
\end{exercise}

\section{Special functions and equations}

\begin{exercise}
Starting from $W+\log W=\log x$, derive
\[
  W=L_1-L_2+\frac{L_2}{L_1}
  +\frac{L_2^2/2-L_2}{L_1^2}
  +\bigO\left(\frac{L_2^3}{L_1^3}\right).
\]
\end{exercise}

\begin{exercise}
For $y'+xy=1$, seek
\[
  y\sim\sum_{n\ge0}a_nx^{-2n-1}.
\]
Find the recurrence, solve it, and identify the exponentially small homogeneous solution omitted by the power sector.
\end{exercise}

\begin{exercise}
Use \cref{eq:riccati-recurrence} to compute $a_1,\ldots,a_5$ for \cref{eq:riccati-model}. Then derive the leading factor $\e^{-x}x^2$ of the first nonperturbative sector from the linearized equation.
\end{exercise}

\begin{exercise}
For the Airy Riccati equation $u'+u^2=x$, take
\[
  u=-x^{1/2}+c_1x^{-1}+c_2x^{-5/2}+c_3x^{-4}+\cdots.
\]
Find $c_1,c_2,c_3$. Integrate to recover the leading exponential and algebraic prefactor of the decaying Airy solution.
\end{exercise}

\begin{exercise}
Solve \cref{eq:boundary-two-point} directly and identify which terms are visible in the outer expansion for fixed $0<x<1$. What changes when $x=\varepsilon X$?
\end{exercise}

\begin{exercise}
Apply the operator expansion
\[
  \frac{D}{\e^D-1}=1-\frac D2+\frac{D^2}{12}
  -\frac{D^4}{720}+\cdots
\]
to $\log x$ and integrate to recover Stirling's series through $x^{-3}$.
\end{exercise}

\begin{exercise}[Abel coordinate]
For $T(x)=x+x^{-1}$, substitute
\[
 A(x)=\frac{x^2}{2}+c\log x+a x^{-2}+\bigO(x^{-4})
\]
into $A(T)-A=1$. Determine $c$ and $a$.
\end{exercise}

\begin{exercise}
Use \cref{eq:fractional-example} to write $T^{[1/2]}(x)$ through $x^{-7}$. Check formally that composing this half-iterate with itself gives $x+x^{-1}+\bigO(x^{-9})$.
\end{exercise}

\section{Borel summation and resurgence}

\begin{exercise}
Compute the Borel transform and Borel sum of
\[
  \sum_{n\ge0}(-1)^n n!x^{-n-1}.
\]
Verify the first three asymptotic coefficients by integration by parts.
\end{exercise}

\begin{exercise}
For
\[
  \widehat\phi(p)=\frac{1}{(1-p/A)^2},
\]
find the coefficients $a_n$ of the original formal series
\[
  \widetilde\phi(x)=\sum_{n\ge0}a_nx^{-n-1}.
\]
What exponential scale is predicted by the Borel singularity?
\end{exercise}

\begin{exercise}
Compute the difference between upper and lower lateral Laplace integrals of
\[
  \frac{1}{1-p/A},
  \qquad A>0,
\]
along the positive ray. State clearly your orientation convention.
\end{exercise}

\begin{exercise}
Show that the Borel transform turns the product
\[
  x^{-1}\cdot x^{-1}=x^{-2}
\]
into a convolution, using the convention
$\Borel(x^{-n-1})=p^n/n!$.
\end{exercise}

\begin{exercise}
Suppose
\[
  a_n\sim C A^{-n}\Gamma(n+\beta),
  \qquad A>0.
\]
Use the ratio of successive terms to estimate the optimal truncation index for
$\sum a_nx^{-n}$. Use Stirling's formula to identify the exponential scale of the least term.
\end{exercise}

\begin{exercise}[Large-order extraction]
Assume exact data
\[
  a_n=S\frac{\Gamma(n+\beta)}{2\pi i A^{n+\beta}}
  \left(b_0+\frac{A b_1}{n+\beta-1}\right).
\]
Construct a sequence that tends to $Sb_0/(2\pi i)$ as $n\to\infty$, and a second rescaled difference that tends to $Sb_1/(2\pi i)$.
\end{exercise}

\section{Structural questions}

\begin{exercise}
Explain why a Hardy field cannot contain the germ of $\sin x$. Which field axiom fails if one tries to adjoin it?
\end{exercise}

\begin{exercise}
Give an example of two distinct functions with the same complete inverse-power asymptotic expansion at $+\infty$. Then give a transseries scale that distinguishes them.
\end{exercise}

\begin{exercise}
Explain the difference among the following three claims about a formal expression: it solves an equation formally; it is asymptotic to a function; it is Borel summable to a function.
\end{exercise}

\begin{exercise}
Why does sending $x$ to the surreal number $\omega$ make $x^{-1}$ infinitesimal and $\e^x$ larger than every power $x^n$? Express the answer using order-preserving field and exponential operations.
\end{exercise}

\begin{exercise}
Make a two-column comparison of grid-based and well-based transseries. Include support, computational convenience, and closure/generalization.
\end{exercise}

\begin{exercise}[Designing an ansatz]
A nonlinear equation has a perturbative power series and its linearization has two independent small solutions
\[
  \e^{-Ax}x^\alpha,
  \qquad
  \e^{-Bx}x^\beta,
\]
where $A,B>0$ are rationally independent. Write a natural two-parameter transseries ansatz and describe the action lattice generated by nonlinear products.
\end{exercise}

\begin{exercise}[Residual certification]
Suppose a first-order equation $E[y]=0$ has linearization $L$ at a formal solution and an approximation $y_N$ satisfies
\[
  E[y_N]=r\m+\text{smaller terms}.
\]
Explain how to choose a correction $c\n$ from the dominant relation
$L[c\n]\sim-r\m$. What can go wrong if $L[\n]$ is anomalously small?
\end{exercise}

\begin{exercise}[Saddles]
For the Gaussian integral
\[
  I(x)=\int_{-\infty}^{\infty}\e^{-xt^2/2}\dd t,
\]
derive its saddle scaling and evaluate it exactly. Then explain why a nonquadratic analytic potential generally produces an infinite inverse-power prefactor series.
\end{exercise}

\chapter{Solutions to the exercises}
\label{chap:solutions}

\section{Asymptotic scales and support}

\subsection*{Solution 1}
The order is
\[
 (\log x)^{100}x^{-3}
 \succ x^{-10}
 \succ \e^{-\sqrt{x}}x^2
 \succ \e^{-\sqrt{x}}
 \succ \e^{-x/\log x}
 \succ \e^{-x}.
\]
The first comparison follows because $(\log x)^{100}x^7\to\infty$. Any inverse power dominates $x^2\e^{-\sqrt{x}}$ because exponentials beat powers. Multiplying by $x^2$ makes $x^2\e^{-\sqrt{x}}$ larger than $\e^{-\sqrt{x}}$. Since
\[
  \frac{x}{\log x}\gg\sqrt{x},
\]
we have $\e^{-\sqrt{x}}\gg\e^{-x/\log x}$. Finally $x/\log x<x$ for large $x$, so $\e^{-x/\log x}\gg\e^{-x}$.

\subsection*{Solution 2}
For the first limit, repeated l'H\^opital gives
\[
  \lim_{x\to\infty}\frac{\log x}{x^a}
  =\lim_{x\to\infty}\frac{x^{-1}}{ax^{a-1}}
  =\lim_{x\to\infty}\frac1{ax^a}=0.
\]
For the second, take logarithms of the positive ratio:
\[
  \log\frac{x^b}{\e^{x^a}}=b\log x-x^a\to-\infty
\]
by the first result. Hence the ratio tends to zero.

\subsection*{Solution 3}
Let
\[
  F_k(x)=\int_0^\infty\frac{\e^{-xp}}{(1+p)^{k+1}}\dd p.
\]
Integrating by parts with $u=(1+p)^{-k-1}$ and
$\dd v=\e^{-xp}\dd p$ gives
\[
  F_k=\frac1x-\frac{k+1}{x}F_{k+1}.
\]
Iteration yields
\[
 F(x)=\sum_{n=0}^{N-1}(-1)^nn!x^{-n-1}
 +(-1)^N N!x^{-N}F_N(x).
\]
Since $0<F_N(x)\le\int_0^\infty\e^{-xp}\dd p=x^{-1}$,
\[
 |R_N(x)|\le N!x^{-N-1}.
\]
This proves the asymptotic expansion.

\subsection*{Solution 4}
$S_1,S_3,$ and $S_5$ are legal. Every nonempty subset has a largest monomial because decreasing an exponent $n$ or $m$ increases the monomial, and the relevant set of integer pairs has a lexicographically minimal element. The support $S_2$ is illegal: the whole set has no largest element. The support $S_4$ is illegal because
\[
 x^{-1}\prec x^{-1/2}\prec x^{-1/3}\prec\cdots
\]
has no largest member.

\subsection*{Solution 5}
Both $A$ and $B$ have magnitude $x^2$ and dominant term $3x^2$. Their difference is
\[
  A-B=\e^{-x}-x^{-7}\sim-x^{-7},
\]
so $\magop(A-B)=x^{-7}$. Since $A-B\prec A$, we have $A\sim B$.

\subsection*{Solution 6}
Assume infinitely many factorizations $g=s_nt_n$. After passing to a subsequence, reverse well-ordering lets us take $s_1\prec s_2\prec\cdots$ or an analogous strict monotone pattern among distinct $s_n$. Then
\[
  t_n=g/s_n
\]
moves strictly in the opposite direction, producing an infinite increasing sequence in $T$ with no largest element. This contradicts reverse well-ordering. A precise proof uses the fact that every infinite subset of a reverse well-ordered set has a monotone subsequence whose forbidden direction cannot continue indefinitely.

\section{Algebra, exponentials, and logarithms}

\subsection*{Solution 7}
Factor $x$ and set $u=x^{-1}+x^{-1}\e^{-x}$:
\[
 \frac1{x+1+\e^{-x}}=x^{-1}(1+u)^{-1}.
\]
Expanding and retaining the requested scales gives
\begin{align*}
 \frac1{x+1+\e^{-x}}
 ={}&x^{-1}-x^{-2}+x^{-3}-x^{-4}\\
 &{}-x^{-2}\e^{-x}+2x^{-3}\e^{-x}
 +\text{smaller terms}.
\end{align*}

\subsection*{Solution 8}
Let $u=-x^{-1}+2\e^{-x}$. Then
\[
 \log T=\log2+3\log x+\sqrt{x}+\log(1+u).
\]
Using $u-u^2/2+u^3/3$ gives
\begin{align*}
 \log T={}&\sqrt{x}+3\log x+\log2
 -x^{-1}-\frac12x^{-2}-\frac13x^{-3}\\
 &{}+2\e^{-x}+2x^{-1}\e^{-x}
 +2x^{-2}\e^{-x}+\text{smaller terms}.
\end{align*}

\subsection*{Solution 9}
Factor $x^2$:
\[
 (x^2+x+\e^{-x})^{3/2}
 =x^3(1+x^{-1}+x^{-2}\e^{-x})^{3/2}.
\]
The binomial coefficients are $1,3/2,3/8,-1/16,3/128,\ldots$. Hence
\[
 x^3+\frac32x^2+\frac38x-\frac1{16}
 +\frac{3}{128x}
 +\frac32x\e^{-x}+\text{smaller terms}.
\]
Cross terms with the exponential are smaller than the leading $x\e^{-x}$ term.

\subsection*{Solution 10}
We have
\[
 \e^{-T}=\e^{-x}x^{-2}\e^{-x^{-1}}.
\]
Therefore
\[
 \e^{-T}=\e^{-x}\left(
 x^{-2}-x^{-3}+\frac12x^{-4}
 +\bigO(x^{-5})\right).
\]

\subsection*{Solution 11}
Write $A+B=A(1+B/A)$. Since $B/A\prec1$,
\[
 A+B\sim A
\]
and
\[
 \log(A+B)=\log A+\log(1+B/A)
 =\log A+B/A+\bigO((B/A)^2).
\]
The family of powers $(B/A)^n$ must be summable; this holds in the standard Hahn setting when $B/A$ is infinitesimal with well-based support.

\section{Differentiation, integration, and composition}

\subsection*{Solution 12}
Taking a logarithm,
\[
 \log\m=a\log x+b\log\log x-x+\sqrt{x}.
\]
Thus
\[
 \m^\dagger=-1+\frac1{2\sqrt{x}}
 +\frac ax+\frac{b}{x\log x}.
\]
The parenthesis is asymptotic to $-1$, so
\[
  \m'\sim-\m.
\]

\subsection*{Solution 13}
Let $B=b_0x^{-2}+b_1x^{-5}+b_2x^{-8}+b_3x^{-11}$. We need
\[
 B'+3x^2B=1+\bigO(x^{-12}).
\]
Matching powers gives
\[
 3b_0=1,
 \quad 3b_1-2b_0=0,
 \quad 3b_2-5b_1=0,
 \quad 3b_3-8b_2=0.
\]
Hence
\[
 b_0=\frac13,
 \quad b_1=\frac29,
 \quad b_2=\frac{10}{27},
 \quad b_3=\frac{80}{81}.
\]

\subsection*{Solution 14}
Factor $x^2$ and set $u=x^{-1}+x^{-2}\e^{-x}$. Then
\begin{align*}
 \log(x^2+x+\e^{-x})
 ={}&2\log x+x^{-1}-\frac12x^{-2}
 +\frac13x^{-3}-\frac14x^{-4}\\
 &{}+x^{-2}\e^{-x}+\text{smaller terms}.
\end{align*}

\subsection*{Solution 15}
Write $G=x-L+a/x+b/x^2$ with $L=\log x$ and allow $a,b$ to depend polynomially on $L$. Expanding
\[
 \log G=L+\log\left(1-\frac Lx+\frac a{x^2}+\frac b{x^3}\right)
\]
gives, through $x^{-2}$,
\[
 F(G)-x=\frac{a-L}{x}
 +\frac{b+a-L^2/2}{x^2}+\cdots.
\]
Thus $a=L$ and $b=L^2/2-L$. Therefore
\[
 F^{[-1]}(x)=x-L+\frac Lx
 +\frac{L^2/2-L}{x^2}+\cdots.
\]
Substitution makes the residual $\bigO(L^3/x^3)$.

\subsection*{Solution 16}
The additive decomposition regards $G$ as the large part and $U$ as infinitesimal relative to the chosen Taylor scale. By the definition of transseries exponentiation,
\[
 \e^{G+U}=\e^G\e^U
 =\e^G\sum_{n\ge0}U^n/n!.
\]
For $U=x$, the powers $1,x,x^2,\ldots$ have no largest term and are not a legal decreasing support at infinity; $\e^x$ must be treated as one monomial.

\section{Special functions and equations}

\subsection*{Solution 17}
Set $W=L_1-L_2+u$. The reduced equation is
\[
 u+\log\left(1-\frac{L_2-u}{L_1}\right)=0.
\]
Assume $u=P_1/L_1+P_2/L_1^2+\cdots$. At order $L_1^{-1}$,
$P_1=L_2$. At order $L_1^{-2}$,
\[
 P_2=L_2^2/2-L_2.
\]
This gives the stated expansion.

\subsection*{Solution 18}
Substitution into $y'+xy=1$ gives
\[
 a_0=1,
 \qquad a_{n+1}=(2n+1)a_n.
\]
Hence $a_n=(2n-1)!!$. The homogeneous equation $y'+xy=0$ has solution
\[
  C\e^{-x^2/2},
\]
which is smaller than every inverse power and therefore absent from the algebraic sector.

\subsection*{Solution 19}
The recurrence gives
\[
 a_1=1,
 \quad a_2=2,
 \quad a_3=8,
 \quad a_4=44,
 \quad a_5=296.
\]
The linearized perturbation obeys
\[
 \delta'+(1-2y_0)\delta=0.
\]
Thus
\[
 \delta=C\exp\left(-x+2\int y_0\dd x\right)
 =C\e^{-x}x^2(1+\smallo(1)).
\]

\subsection*{Solution 20}
Insert the ansatz into $u'+u^2=x$. Matching $x^{-1/2}$ gives
\[
 -\frac12-2c_1=0,
 \quad c_1=-\frac14.
\]
Matching $x^{-2}$ gives
\[
 -c_1+c_1^2-2c_2=0,
 \quad c_2=\frac5{32}.
\]
Matching $x^{-7/2}$ gives
\[
 -\frac52c_2+2c_1c_2-2c_3=0,
 \quad c_3=-\frac{15}{64}.
\]
Integrating begins with
\[
 \log y=-\frac23x^{3/2}-\frac14\log x+\cdots,
\]
so
\[
 y\asymp x^{-1/4}\e^{-2x^{3/2}/3}.
\]

\subsection*{Solution 21}
Solving once gives $y'=C\e^{-x/\varepsilon}$ and then
\[
 y=A+B\e^{-x/\varepsilon}.
\]
The boundary conditions yield
\[
 y=\frac{1-\e^{-x/\varepsilon}}{1-\e^{-1/\varepsilon}}.
\]
For fixed $0<x<1$, both exponentials are beyond all powers and the outer expansion is $y\sim1$. If $x=\varepsilon X$, then $\e^{-x/\varepsilon}=\e^{-X}$ is order one and resolves the layer at $x=0$.

\subsection*{Solution 22}
Applying the operator to $\log x$ gives
\[
 F'=\log x-\frac1{2x}-\frac1{12x^2}
 +\frac1{120x^4}+\cdots.
\]
Integrating,
\[
 F=\left(x-\frac12\right)\log x-x+C
 +\frac1{12x}-\frac1{360x^3}+\cdots.
\]
For $F=\log\Gamma$, $C=\frac12\log(2\pi)$.

\subsection*{Solution 23}
Use
\[
 T^2=x^2+2+x^{-2},
 \qquad
 \log T=\log x+\log(1+x^{-2}).
\]
The constant term in $A(T)-A$ is already $1$. Since
\[
 \log(T/x)=\log(1+x^{-2})=x^{-2}-\frac12x^{-4}+\cdots
\]
and
\[
 T^{-2}-x^{-2}=-2x^{-4}+\cdots,
\]
the coefficient of $x^{-2}$ is $1/2+c$, so $c=-1/2$. The coefficient of $x^{-4}$ is $-c/2-2a$, so $a=1/8$. Thus
Thus
\[
 A=\frac{x^2}{2}-\frac12\log x+\frac1{8x^2}+\cdots.
\]

\subsection*{Solution 24}
Set $s=1/2$ in \cref{eq:fractional-example}:
\[
 T^{[1/2]}(x)=x+\frac1{2x}
 +\frac1{8x^3}+\frac1{16x^5}
 +\frac{3}{128x^7}+\bigO(x^{-9}).
\]
Taylor-composing this expression with itself and collecting powers gives
\[
 x+x^{-1}+\bigO(x^{-9}).
\]
The group law of the Abel construction guarantees the cancellation; direct expansion verifies it coefficient by coefficient.

\section{Borel summation and resurgence}

\subsection*{Solution 25}
The Borel transform is
\[
 \sum_{n\ge0}(-p)^n=\frac1{1+p}.
\]
Therefore
\[
 \mathscr S_0\widetilde F
 =\int_0^\infty\frac{\e^{-xp}}{1+p}\dd p
 =-\e^x\Ei(-x).
\]
Three integrations by parts give
\[
 -\e^x\Ei(-x)
 =x^{-1}-x^{-2}+2x^{-3}+\bigO(x^{-4}),
\]
matching $0!, -1!, 2!$ with alternating signs.

\subsection*{Solution 26}
Since
\[
 (1-z)^{-2}=\sum_{n\ge0}(n+1)z^n,
\]
we have
\[
 \widehat\phi(p)=\sum_{n\ge0}(n+1)A^{-n}p^n.
\]
Thus
\[
 a_n=n!(n+1)A^{-n}=(n+1)!A^{-n}.
\]
The singularity at $p=A$ predicts the exponential scale $\e^{-Ax}$.

\subsection*{Solution 27}
The integrand is
\[
 \frac{\e^{-xp}}{1-p/A}.
\]
Its residue at $p=A$ is $-A\e^{-Ax}$. With upper contour from $0$ to $+\infty$ and lower contour subtracted in the same direction, the resulting closed contour is clockwise, so
\[
 \mathscr S_{0^+}-\mathscr S_{0^-}
 =-2\pi i(-A\e^{-Ax})
 =2\pi iA\e^{-Ax}.
\]
Changing the convention reverses the sign.

\subsection*{Solution 28}
We have
\[
 \Borel(x^{-1})=1.
\]
The convolution is
\[
 (1*1)(p)=\int_0^p1\dd q=p.
\]
But
\[
 \Borel(x^{-2})=p,
\]
so the rule holds in this simplest case.

\subsection*{Solution 29}
The successive-term ratio is approximately
\[
 \frac{n+\beta}{Ax}.
\]
It crosses one near $n\approx Ax-\beta$, so the optimal index is $N\approx Ax$. Stirling's formula then gives a least-term scale
\[
 x^{\beta-1/2}\e^{-Ax}
\]
up to a constant and convention-dependent power of $A$.

\subsection*{Solution 30}
Define
\[
 R_n:=\frac{A^{n+\beta}a_n}{\Gamma(n+\beta)}.
\]
Then
\[
 R_n=\frac{S}{2\pi i}\left(
 b_0+\frac{Ab_1}{n+\beta-1}\right)
 \longrightarrow \frac{Sb_0}{2\pi i}.
\]
Let $L=\lim R_n$. Then
\[
 \frac{n+\beta-1}{A}(R_n-L)
 \longrightarrow \frac{Sb_1}{2\pi i}.
\]
With finite data, one first extrapolates $R_n$ to estimate $L$ and then applies the second rescaling.

\section{Structural questions}

\subsection*{Solution 31}
The germ of $\sin x$ has infinitely many zeros but is not eventually zero. In a field, every nonzero element must have a multiplicative inverse. The reciprocal $1/\sin x$ is not a germ of a function defined on any entire tail interval because poles occur arbitrarily far out. Hence field closure fails.

\subsection*{Solution 32}
The functions $0$ and $\e^{-x}$ have the same complete inverse-power expansion, namely zero. The enlarged scale
\[
  1,x^{-1},x^{-2},\ldots,\e^{-x},x^{-1}\e^{-x},\ldots
\]
distinguishes them at the $\e^{-x}$ monomial.

\subsection*{Solution 33}
A formal solution means substitution produces the zero formal transseries. An asymptotic representation means an actual function has remainders controlled by every finite truncation on a stated domain. Borel summability is stronger and more specific: after Borel continuation and a directional Laplace transform, a prescribed analytic function is reconstructed from the formal coefficients.

\subsection*{Solution 34}
An order-preserving embedding sends $x>n$ for every integer $n$ to $\omega>n$. Taking reciprocals reverses positive inequalities, so
\[
  0<\omega^{-1}<1/n
\]
for every $n$. Because the exponential is order preserving and grows faster than powers in the embedded transseries field,
\[
  \exp(\omega)>\omega^n
\]
for every fixed $n$.

\subsection*{Solution 35}
Grid-based supports lie in a finitely generated multiplicative grid, which makes exponents reducible to integer multi-indices and supports effective computation. Well-based supports need only be reverse well ordered; they include more series and often have stronger conceptual closure, but can require subtler support arguments and may be less directly algorithmic.

\subsection*{Solution 36}
Use two parameters $C_1,C_2$:
\[
 y=\sum_{k,\ell\ge0}
 C_1^kC_2^\ell
 \e^{-(kA+\ell B)x}
 x^{k\alpha+\ell\beta}\Phi_{k,\ell}(x).
\]
The action lattice is
\[
  \set{kA+\ell B:k,\ell\in\N}.
\]
Rational independence prevents accidental equality between distinct pairs, while nonlinear products add the index pairs.

\subsection*{Solution 37}
Compute the dominant effect of $L$ on candidate monomials. If
\[
  L[\n]\sim\lambda\m,
\]
choose $c=-r/\lambda$. The new residual is smaller. If $L[\n]$ is anomalously small because of resonance or because $\n$ lies near the kernel, division by the dominant coefficient fails. One may need a logarithmic factor, a different monomial, or a new free homogeneous sector.

\subsection*{Solution 38}
Set $u=\sqrt{x}\,t$. Then
\[
 I(x)=x^{-1/2}\int_{-\infty}^{\infty}\e^{-u^2/2}\dd u
 =\sqrt{\frac{2\pi}{x}}.
\]
The potential is exactly quadratic, so the saddle expansion terminates. For a nonquadratic analytic $V$, Taylor terms $V^{(k)}(0)t^k$ become powers of $x^{-1/2}$ after rescaling; Gaussian moments generate an infinite asymptotic series, usually factorially divergent because neighboring saddles or singularities control late orders.

\appendix

\chapter{A proof kit for formal asymptotics}
\label{app:proofkit}

This appendix collects short arguments that recur throughout the book.  They are not a substitute for the full support theorems in the specialist literature, but they explain why the elementary manipulations work and exactly where the nontrivial hypotheses enter.

\section{Uniqueness of coefficients in an asymptotic scale}

Let $\phi_0,\phi_1,\ldots$ be nonzero functions on a tail interval, with
\[
 \phi_{n+1}=\smallo(\phi_n)
 \qquad(x\to+\infty).
\]
Suppose a function $f$ has two asymptotic expansions in this same scale,
\[
 f\sim\sum_{n\ge0}a_n\phi_n,
 \qquad
 f\sim\sum_{n\ge0}b_n\phi_n.
\]
Then $a_n=b_n$ for every $n$.

\begin{proof}
If the coefficients differ, let $k$ be the first index with $a_k\ne b_k$. Subtract the two expansions through index $k$:
\[
 0=(a_k-b_k)\phi_k+\smallo(\phi_k).
\]
Divide by $\phi_k$. The left side remains zero, while the right side tends to $a_k-b_k\ne0$, a contradiction.
\end{proof}

\begin{intuitionbox}
Asymptotic coefficients are read from left to right.  Once all larger scales have been removed, the first surviving scale cannot be canceled by smaller ones.
\end{intuitionbox}

\section{Why exponentials beat powers}

For every fixed $a>0$ and every real $N$,
\[
 \e^{-ax}=\smallo(x^{-N})
 \qquad(x\to+\infty).
\]
Indeed,
\[
 \log\!\left(x^N\e^{-ax}\right)=N\log x-ax\longrightarrow-\infty,
\]
because $\log x/x\to0$. Exponentiating gives $x^N\e^{-ax}\to0$. Replacing $x$ by $x^\alpha$ with $\alpha>0$ gives the more general comparison
\[
 \e^{-a x^\alpha}=\smallo(x^{-N}).
\]
This elementary fact is the mathematical reason that an exponential correction is invisible to every finite inverse-power expansion.

\section{Leading monomials and inverses}

Let
\[
 T=a\m(1+\varepsilon),
 \qquad a\ne0,
 \qquad \varepsilon\prec1,
\]
where $\m$ is a monomial and every monomial in $\varepsilon$ is smaller than $1$. Formally,
\[
 T^{-1}=a^{-1}\m^{-1}(1+\varepsilon)^{-1}
       =a^{-1}\m^{-1}\sum_{n\ge0}(-\varepsilon)^n.
\]
The algebraic identity is just the geometric series. The support issue is subtler: for any fixed output monomial, only finitely many products from the powers $\varepsilon^n$ may contribute.  A well-based or grid-based support condition guarantees this local finiteness.

\begin{proposition}[geometric inversion in a finitely generated grid]
Let $\boldsymbol\mu=(\mu_1,\ldots,\mu_r)$ with each $\mu_i\prec1$, and suppose every monomial of $\varepsilon$ belongs to
\[
 \boldsymbol\mu^{\N^r}
 =\set{\mu_1^{k_1}\cdots\mu_r^{k_r}:k_i\in\N}
\]
and has positive total degree. Then the coefficient of any given grid monomial in
$\sum_{n\ge0}(-\varepsilon)^n$ receives contributions from only finitely many $n$.
\end{proposition}

\begin{proof}
A monomial occurring in $\varepsilon^n$ has total multi-index degree at least $n$. A fixed target multi-index has finite total degree, so it cannot occur once $n$ exceeds that degree.
\end{proof}

\section{Formal exponentials and logarithms}

If $S\prec1$, define
\[
 \exp(S)=\sum_{n\ge0}\frac{S^n}{n!},
 \qquad
 \log(1+S)=\sum_{n\ge1}\frac{(-1)^{n+1}}{n}S^n.
\]
In a formal transseries field, these formulas do not assert numerical convergence. They assert that the families of supports are summable: each monomial receives a well-defined coefficient, and the union of all resulting supports remains legal.

For a general positive transseries, first factor its dominant monomial:
\[
 T=a\m(1+S),
 \qquad a>0,
 \qquad S\prec1.
\]
Then
\[
 \log T=\log a+\log\m+\log(1+S).
\]
Likewise, split a general exponent into its purely large, constant, and small parts,
\[
 U=U_{\succ1}+c+U_{\prec1},
\]
and write
\[
 \e^U=\e^c\e^{U_{\succ1}}\e^{U_{\prec1}}.
\]
The first two factors determine a new monomial; only the last is expanded as a Taylor series.

\section{A one-step Newton principle}

Suppose we seek $y$ with $F(y)=0$, and an approximation $y_0$ leaves residual
\[
 r=F(y_0).
\]
Let $L=F'[y_0]$ denote the linearized operator. Choose a correction $\delta$ satisfying, to dominant order,
\[
 L\delta\sim-r.
\]
Taylor expansion gives
\[
 F(y_0+\delta)=r+L\delta+Q(\delta),
\]
where $Q$ is at least quadratic. If the chosen $\delta$ cancels the dominant monomial of $r$ and if $Q(\delta)$ is smaller, the new residual has strictly smaller magnitude. Repetition produces a transseries coefficient by coefficient.

\begin{warningbox}
The step can fail at a resonance. If the dominant coefficient of $L$ vanishes on the candidate monomial, the correction may require a logarithm, a changed power, or a new free homogeneous sector. Resonance is information, not merely an algebraic inconvenience.
\end{warningbox}

\section{Why formal differentiation is usually easy}

For a monomial
\[
 \m=x^a(\log x)^b\exp(L),
\]
logarithmic differentiation gives
\[
 \frac{\m'}{\m}
 =\frac{a}{x}+\frac{b}{x\log x}+L'.
\]
Thus
\[
 \m'=\m\left(\frac{a}{x}+\frac{b}{x\log x}+L'\right).
\]
A transseries derivative is obtained termwise, provided the differentiated family remains summable. In the standard LE fields this closure is built into the construction.  For arbitrary generalized series, however, termwise differentiation is not automatic: a derivative could move infinitely many terms onto the same monomial or destroy well-basedness.

\section{Asymptotic validity of algebraic operations}

Suppose actual functions $f$ and $g$ have asymptotic expansions in compatible scales and nonzero leading terms. Then finite truncation estimates immediately justify addition and multiplication. For instance, if
\[
 f=F_N+o(\phi_N),
 \qquad
 g=G_M+o(\psi_M),
\]
then
\[
 fg=F_NG_M+F_No(\psi_M)+G_Mo(\phi_N)+o(\phi_N\psi_M).
\]
After the products are reordered by size, this yields the product asymptotic expansion. Inversion follows from factoring the leading term and applying a geometric expansion with a controlled remainder.

Composition and differentiation require extra hypotheses. Differentiating an $o$-remainder need not preserve its order, and composing can move the argument outside the uniformity region. This is why a formal identity and an analytic asymptotic theorem are logically distinct statements.

\begin{checkpointbox}
Most elementary transseries calculations rest on four moves:
\begin{enumerate}
 \item factor the dominant monomial;
 \item expand only a genuinely small remainder;
 \item cancel the largest residual term;
 \item verify that supports remain locally finite and well ordered.
\end{enumerate}
The deep theory makes these moves into closure theorems for large fields.
\end{checkpointbox}

\chapter{Glossary}
\label{app:glossary}

The entries below favor working intuition, followed by the formal point that matters in calculations.

\begin{description}[style=nextline,itemsep=0.75em]
\item[Action]
A positive quantity $A$ appearing in an exponential scale such as $\e^{-Ax}$. In resurgence, Borel singularities often occur at locations related to actions.

\item[Alien derivative]
An operator, usually written $\Delta_A$, designed to extract the resurgent information carried by a Borel singularity at $p=A$. It is not ordinary differentiation with respect to $x$.

\item[Anti-Stokes direction]
A direction along which two competing exponentials have equal magnitude, so dominance can exchange. Terminology varies in the literature; some authors interchange the labels ``Stokes'' and ``anti-Stokes.''

\item[Asymptotic couple]
An ordered abelian group together with a map recording the dominant value of logarithmic derivatives. It abstracts the interaction between valuation and differentiation in an $H$-field.

\item[Asymptotic expansion]
A sequence of approximations whose successive remainders are smaller in a specified asymptotic scale. It need not converge for any fixed value of the variable.

\item[Asymptotic scale]
A sequence $\phi_0\succ\phi_1\succ\phi_2\succ\cdots$ with $\phi_{n+1}=o(\phi_n)$ on the limit regime under study.

\item[Beyond all orders]
Smaller than every term in a chosen perturbative scale. For example, $\e^{-x}$ is beyond all orders of $x^{-1}$ as $x\to+\infty$.

\item[Borel transform]
A transform that divides factorial growth from coefficients. In the convention of this book,
$\sum_{n\ge0}a_nx^{-n-1}$ maps to $\sum_{n\ge0}a_np^n/n!$.

\item[Borel sum]
An analytic function obtained by analytically continuing a Borel transform and applying a directional Laplace integral. It assigns meaning to many divergent Gevrey series.

\item[Compositional inverse]
A series $G$ satisfying $F\circ G=G\circ F=x$. It is different from the multiplicative inverse $1/F$.

\item[Dominance]
The comparison $F\succ G$, meaning $G/F\to0$ in an analytic setting or that the leading monomial of $F$ is larger in a formal ordered field.

\item[EL-series]
Exponential-logarithmic series built by a construction order that starts from exponentials and then closes under logarithmic operations. Their relation to LE-series is subtle rather than merely notational.

\item[Exponential height]
A rough measure of how many nested exponential layers are needed to construct a monomial or transseries.

\item[Flat function]
A nonzero function with zero asymptotic power series in a given scale. The standard example at $+\infty$ is $\e^{-x}$ relative to inverse powers.

\item[Formal solution]
A transseries that makes an equation vanish under formal substitution. This alone does not assert that an actual analytic function realizes the transseries.

\item[Gevrey order]
A quantitative rate of coefficient growth. Gevrey-1 typically means $|a_n|\le CA^n n!$; such growth is suited to ordinary Borel transformation.

\item[Germ at $+\infty$]
An equivalence class of functions that agree on some sufficiently far tail interval. Germs forget behavior on bounded sets.

\item[Grid]
A finitely generated multiplicative lattice of monomials, often written $\boldsymbol\mu^{\mathbf{k}}$ for integer multi-indices $\mathbf{k}$ in a translated cone.

\item[Grid-based transseries]
A transseries whose support lies in a finite-rank grid. This setting is especially convenient for algorithms.

\item[Hahn series]
A generalized formal sum $\sum_{\g\in G}a_\g\g$ over an ordered abelian group, with support satisfying a well-order condition that makes algebra possible.

\item[Hardy field]
A field of germs at infinity closed under differentiation. Ordering is eventual ordering of functions.

\item[Height wins]
A useful informal rule: a genuinely positive exponential layer generally outruns all combinations from lower exponential height, although signs and cancellations must still be checked.

\item[$H$-field]
An ordered differential field satisfying axioms modeled on Hardy fields, including compatibility between order and differentiation for large elements.

\item[Hyperasymptotics]
The repeated asymptotic analysis of optimally truncated remainders, exposing successively smaller exponential scales.

\item[Infinitesimal]
An element smaller in magnitude than every positive real constant. In transseries at $+\infty$, $x^{-1}$ and $\e^{-x}$ are infinitesimal.

\item[Instanton sector]
In physics-inspired terminology, a sector proportional to $\e^{-kAx}$, often multiplied by powers and a perturbative series. The label need not imply literal instantons in the physical problem.

\item[Lateral Borel sum]
A Laplace integral taken just above or just below a singular direction in the Borel plane. The difference of the two lateral sums is a Stokes jump.

\item[LE-series]
Logarithmic-exponential series obtained from a base variable by iterated algebraic operations, logarithms, exponentials, and suitable generalized summation.

\item[Leading coefficient]
The coefficient multiplying the largest monomial of a nonzero transseries.

\item[Leading monomial or magnitude]
The largest monomial in the support of a nonzero transseries. It determines first-order size.

\item[Liouville closed]
An ordered differential field in which algebraic equations, antiderivatives, and logarithmic antiderivatives have the expected solutions. This is a differential-algebraic closure property, not ordinary topological completeness.

\item[Logarithmic depth]
A measure of how many iterated logarithms are needed below the main variable, as in $\log x$, $\log\log x$, and so on.

\item[Monomial]
A formal object representing one pure order of growth, such as $x^a(\log x)^b\e^{-x}$.

\item[Newton polygon]
A geometric device for balancing terms and discovering candidate leading exponents. Differential versions organize dominant balances for ODEs.

\item[Newtonian $H$-field]
Roughly, an $H$-field in which differential polynomials of Newton degree one have zeros. This property plays a role analogous to algebraic closedness for differential equations.

\item[Normal form]
A canonical decomposition such as $T=a\m(1+\varepsilon)$ with $a\ne0$, dominant monomial $\m$, and $\varepsilon\prec1$.

\item[Omega-series]
A large class of surreal transseries generated from the simplest infinite surreal number $\omega$ by transserial operations and summation.

\item[Optimal truncation]
Stopping a divergent asymptotic series near its least term. The resulting error is often exponentially small and close to the best accuracy available from that sector alone.

\item[Ratio set]
A finite set of small monomials generating a grid. Ratio sets provide a computational encoding of relative smallness.

\item[Resonance]
A degeneracy in coefficient matching, often caused by an exponent difference or linearized eigenvalue hitting an integer relation. It can force logarithms or new free parameters.

\item[Resurgence]
A property in which the Borel transform has controlled analytic continuation and its singularities encode relations among different transseries sectors.

\item[Saddle point]
A critical point of an exponent in a Laplace- or steepest-descent integral. Different saddles produce different exponential sectors.

\item[Sector]
Either a geometric angular region in the complex plane or, by context, one exponential block $\e^{-kAx}\Phi_k(x)$ of a transseries.

\item[Stokes automorphism]
An operator describing how a resummed transseries changes when a Stokes direction is crossed.

\item[Stokes constant]
A problem-dependent constant measuring the strength of a Stokes connection between sectors. Its numerical normalization depends on conventions.

\item[Stokes direction]
A Borel-Laplace direction obstructed by singularities, across which lateral sums differ. In WKB language, naming conventions can vary.

\item[Summable family]
A formal family for which every monomial receives only finitely many coefficient contributions and whose union of supports remains legal.

\item[Support]
The set of monomials with nonzero coefficient in a generalized series.

\item[Surreal number]
An element of Conway's proper-class-sized ordered field containing the real numbers, infinite and infinitesimal numbers, and a rich exponential structure.

\item[Transmonomial]
A monomial in a transseries group, including nested exponential and logarithmic scales rather than only ordinary powers.

\item[Transseries]
A generalized formal series built from ordered monomials and closed, in an appropriate framework, under operations such as exponentiation, logarithm, differentiation, integration, and composition.

\item[Valuation]
A map recording the leading order of a nonzero element. Depending on convention, a larger valuation may correspond to a smaller magnitude.

\item[Well-based support]
A support with no forbidden infinite ascent in monomial size; equivalently, under the opposite valuation order it is well ordered. This ensures a largest term and controls products and substitutions.

\item[WKB method]
A method for differential equations with a large or small parameter, using exponential phases and amplitude series. It naturally produces transseries and Stokes phenomena.
\end{description}

\chapter{Notation and formula sheet}
\label{app:formulas}

\section{Basic symbols}

\renewcommand{\arraystretch}{1.25}
\begin{longtable}{>{\raggedright\arraybackslash}p{0.23\textwidth}p{0.69\textwidth}}
\toprule
Symbol & Meaning \\
\midrule
\endfirsthead
\toprule
Symbol & Meaning \\
\midrule
\endhead
$F\prec G$ & $F/G\to0$, or the corresponding formal dominance relation.\\
$F\asymp G$ & Comparable leading size; in some analytic contexts, bounded above and below by constant multiples.\\
$F\sim G$ & Usually $F/G\to1$ for functions; elsewhere in the book, an explicit series after $\sim$ denotes an asymptotic expansion.\\
$\m,\n,\g$ & Generic transmonomials.\\
$\G$ & An ordered multiplicative group of monomials.\\
$\supp T$ & The set of monomials having nonzero coefficient in $T$.\\
$\magop T$ & The leading or largest monomial of $T$.\\
$\Tfield$ & A standard field of logarithmic-exponential transseries.\\
$\Borel$ & Formal Borel transform.\\
$\Laplace_\theta$ & Laplace transform along the ray of argument $\theta$.\\
$\mathscr S_\theta$ & Directional Borel-Laplace summation operator.\\
$\Stokes_\theta$ & Stokes automorphism across direction $\theta$.\\
$\Delta_A$ & Alien derivative associated with a Borel singularity/action $A$.\\
$C$ or $\sigma$ & Free transseries parameter multiplying a homogeneous/exponential sector.\\
$A$ & Exponential action in $\e^{-Ax}$.\\
$W$ & Lambert's inverse of $w\mapsto w\e^w$.\\
$\omega$ & The simplest positive infinite surreal number.\\
\bottomrule
\end{longtable}
\renewcommand{\arraystretch}{1}

\section{Canonical algebraic forms}

\begin{ideabox}
For a nonzero transseries,
\[
 T=a\m(1+\varepsilon),
 \qquad a\ne0,
 \qquad \varepsilon\prec1.
\]
Then
\[
 T^{-1}=a^{-1}\m^{-1}\sum_{n\ge0}(-\varepsilon)^n,
\]
and, for positive $T$,
\[
 \log T=\log a+\log\m+
 \sum_{n\ge1}\frac{(-1)^{n+1}}{n}\varepsilon^n.
\]
For real $\alpha$,
\[
 T^\alpha=a^\alpha\m^\alpha
 \sum_{n\ge0}\binom{\alpha}{n}\varepsilon^n.
\]
\end{ideabox}

If $U=U_{\succ1}+c+U_{\prec1}$, then
\[
 \e^U=\e^c\e^{U_{\succ1}}
 \sum_{n\ge0}\frac{U_{\prec1}^n}{n!}.
\]
For a monomial $\m=\exp(L)$,
\[
 \m'=L'\m.
\]
For $\m=x^a(\log x)^b\e^{L}$,
\[
 \m'=\m\left(\frac{a}{x}+\frac{b}{x\log x}+L'\right).
\]

\section{A standard one-parameter sector ansatz}

For one action $A>0$ and power shift $\beta$,
\[
 y(x;C)=\sum_{k\ge0}C^k\e^{-kAx}x^{k\beta}
 \Phi_k(x),
 \qquad
 \Phi_k(x)=\sum_{n\ge0}a_{k,n}x^{-n}.
\]
With two independent actions,
\[
 y(x;C_1,C_2)=
 \sum_{k,\ell\ge0}C_1^kC_2^\ell
 \e^{-(kA+\ell B)x}
 x^{k\alpha+\ell\beta}\Phi_{k,\ell}(x).
\]
Nonlinearity adds sector indices because exponential monomials multiply.

\section{Borel--Laplace conventions used here}

For
\[
 \widetilde f(x)=\sum_{n\ge0}a_nx^{-n-1},
\]
define
\[
 \widehat f(p)=\Borel\widetilde f(p)
 =\sum_{n\ge0}\frac{a_n}{n!}p^n.
\]
The directional Laplace transform is
\[
 \mathscr S_\theta\widetilde f(x)
 =\int_0^{\e^{i\theta}\infty}
   \e^{-xp}\widehat f(p)\dd p,
\]
when continuation and growth permit. The basic identity
\[
 \int_0^{\e^{i\theta}\infty}\e^{-xp}p^n\dd p
 =\frac{n!}{x^{n+1}}
\]
holds in a compatible sector and explains why the Laplace transform recovers the original formal powers.

Multiplication becomes convolution:
\[
 \Borel(\widetilde f\widetilde g)
 =\widehat f*\widehat g,
 \qquad
 (\widehat f*\widehat g)(p)
 =\int_0^p\widehat f(q)\widehat g(p-q)\dd q.
\]
A simple Borel singularity on the ray produces a lateral jump proportional to $\e^{-Ax}$.

\section{Frequently used model expansions}

\subsection*{Lambert $W$ at $+\infty$}
With $L_1=\log x$ and $L_2=\log\log x$,
\begin{align*}
 W(x)={}&L_1-L_2+\frac{L_2}{L_1}
 +\frac{L_2^2/2-L_2}{L_1^2}\\
 &+\frac{L_2^3/3-(3/2)L_2^2+L_2}{L_1^3}
 +\bigO\!\left(\frac{L_2^4}{L_1^4}\right).
\end{align*}

\subsection*{Airy decay}
For the recessive solution of $y''=xy$,
\[
 \operatorname{Ai}(x)
 \sim\frac{1}{2\sqrt\pi}x^{-1/4}
 \e^{-2x^{3/2}/3}
 \left(1-\frac{5}{48x^{3/2}}
 +\frac{385}{4608x^3}-\cdots\right).
\]

\subsection*{Stirling expansion}
As $x\to+\infty$,
\[
 \log\Gamma(x)
 \sim\left(x-\frac12\right)\log x-x
 +\frac12\log(2\pi)
 +\sum_{n\ge1}
 \frac{B_{2n}}{2n(2n-1)x^{2n-1}}.
\]

\subsection*{Euler's factorial series}
The divergent series
\[
 \widetilde F(x)=\sum_{n\ge0}(-1)^n n!\,x^{-n-1}
\]
has Borel transform $1/(1+p)$ and Borel sum
\[
 F(x)=\int_0^\infty\frac{\e^{-xp}}{1+p}\dd p
 =\e^x E_1(x)
\]
for $\Re x>0$, where $E_1(x)=\int_x^\infty\e^{-t}t^{-1}\dd t$.

\subsection*{Fractional iteration through an Abel coordinate}
If $V(f(x))=V(x)+1$, then the formal $s$-th iterate is
\[
 f^{[s]}(x)=V^{-1}(V(x)+s).
\]
This automatically gives
\[
 f^{[s]}\circ f^{[t]}=f^{[s+t]},
 \qquad f^{[0]}=\operatorname{id},
 \qquad f^{[1]}=f.
\]

\section{A diagnostic checklist}

When confronted with a new asymptotic problem, ask:
\begin{enumerate}
 \item What is the limit regime, and which quantities are large or small?
 \item What dominant balance selects the first monomial?
 \item Which ratios are genuinely small enough to expand?
 \item Does the equation possess a homogeneous solution invisible to the perturbative scale?
 \item Does nonlinearity force multiples or combinations of exponential actions?
 \item Are logarithms required by resonance?
 \item Is the coefficient growth factorial, and where are the predicted Borel singularities?
 \item On which complex sectors should an actual solution realize the formal object?
 \item What changes on crossing a Stokes direction?
 \item Which claims are formal identities, and which need analytic estimates?
\end{enumerate}

\backmatter

\chapter*{Bibliography}
\addcontentsline{toc}{chapter}{Bibliography}
\markboth{Bibliography}{Bibliography}

\begin{thebibliography}{99}

\bibitem{ABS2019}
I.~Aniceto, G.~Ba\c{s}ar, and R.~Schiappa,
``A primer on resurgent transseries and their asymptotics,''
\emph{Physics Reports} \textbf{809} (2019), 1--135.
\url{https://arxiv.org/abs/1802.10441}

\bibitem{ADH2017}
M.~Aschenbrenner, L.~van den Dries, and J.~van der Hoeven,
\emph{Asymptotic Differential Algebra and Model Theory of Transseries},
Annals of Mathematics Studies 195, Princeton University Press, 2017.

\bibitem{ADH2018}
M.~Aschenbrenner, L.~van den Dries, and J.~van der Hoeven,
``On numbers, germs, and transseries,''
in \emph{Proceedings of the International Congress of Mathematicians---Rio de Janeiro 2018}, Vol.~II, World Scientific, 2018, 19--42.
\url{https://arxiv.org/abs/1711.06936}

\bibitem{BerryHowls1990}
M.~V.~Berry and C.~J.~Howls,
``Hyperasymptotics,''
\emph{Proceedings of the Royal Society of London A} \textbf{430} (1990), 653--668.

\bibitem{BM2018}
A.~Berarducci and V.~Mantova,
``Surreal numbers, derivations and transseries,''
\emph{Journal of the European Mathematical Society} \textbf{20} (2018), 339--390.
\url{https://arxiv.org/abs/1503.00315}

\bibitem{BM2019}
A.~Berarducci and V.~Mantova,
``Transseries as germs of surreal functions,''
\emph{Transactions of the American Mathematical Society} \textbf{371} (2019), 3549--3592.
\url{https://arxiv.org/abs/1703.01995}

\bibitem{Borel1899}
\'E.~Borel,
``M\'emoire sur les s\'eries divergentes,''
\emph{Annales scientifiques de l'\'Ecole Normale Sup\'erieure}, series~3, \textbf{16} (1899), 9--131.

\bibitem{Conway1976}
J.~H.~Conway,
\emph{On Numbers and Games},
Academic Press, 1976; second edition, A K Peters, 2001.

\bibitem{Corless1996}
R.~M.~Corless, G.~H.~Gonnet, D.~E.~G.~Hare, D.~J.~Jeffrey, and D.~E.~Knuth,
``On the Lambert $W$ function,''
\emph{Advances in Computational Mathematics} \textbf{5} (1996), 329--359.

\bibitem{Costin2009}
O.~Costin,
\emph{Asymptotics and Borel Summability},
Chapman \& Hall/CRC Monographs and Surveys in Pure and Applied Mathematics 141, 2009.

\bibitem{DMM1997}
L.~van den Dries, A.~Macintyre, and D.~Marker,
``Logarithmic-exponential power series,''
\emph{Journal of the London Mathematical Society}, second series, \textbf{56} (1997), 417--434.

\bibitem{DMM2001}
L.~van den Dries, A.~Macintyre, and D.~Marker,
``Logarithmic-exponential series,''
\emph{Annals of Pure and Applied Logic} \textbf{111} (2001), 61--113.

\bibitem{DahnGoring1987}
B.~I.~Dahn and P.~G\"oring,
``Notes on exponential-logarithmic terms,''
\emph{Fundamenta Mathematicae} \textbf{127} (1987), 45--50.

\bibitem{Dingle1973}
R.~B.~Dingle,
\emph{Asymptotic Expansions: Their Derivation and Interpretation},
Academic Press, 1973.

\bibitem{Dorigoni2019}
D.~Dorigoni,
``An introduction to resurgence, trans-series and alien calculus,''
\emph{Annals of Physics} \textbf{409} (2019), article 167914.
\url{https://arxiv.org/abs/1411.3585}

\bibitem{Ecalle1992}
J.~\'Ecalle,
\emph{Introduction aux fonctions analysables et preuve constructive de la conjecture de Dulac},
Actualit\'es Math\'ematiques, Hermann, Paris, 1992.

\bibitem{Edgar2009}
G.~A.~Edgar,
``Transseries for beginners,''
\emph{Real Analysis Exchange} \textbf{35} (2009/2010), 253--310.
\url{https://arxiv.org/abs/0801.4877}

\bibitem{EdgarComposition}
G.~A.~Edgar,
``Transseries: composition, recursion, and convergence,''
preprint, 2009.
\url{https://arxiv.org/abs/0909.1259}

\bibitem{EdgarFractional}
G.~A.~Edgar,
``Fractional iteration of series and transseries,''
\emph{Transactions of the American Mathematical Society} \textbf{365} (2013), 5805--5832.
\url{https://arxiv.org/abs/1002.2378}

\bibitem{Gonshor1986}
H.~Gonshor,
\emph{An Introduction to the Theory of Surreal Numbers},
London Mathematical Society Lecture Note Series 110, Cambridge University Press, 1986.

\bibitem{Hahn1907}
H.~Hahn,
``\"Uber die nichtarchimedischen Gr\"o\ss ensysteme,''
\emph{Sitzungsberichte der Kaiserlichen Akademie der Wissenschaften, Wien, Mathematisch-Naturwissenschaftliche Klasse} \textbf{116} (1907), 601--655.

\bibitem{Hardy1910}
G.~H.~Hardy,
\emph{Orders of Infinity: The `Infinit\"arcalcul' of Paul du Bois-Reymond},
Cambridge Tracts in Mathematics and Mathematical Physics 12, Cambridge University Press, 1910.

\bibitem{Hoeven2006}
J.~van der Hoeven,
\emph{Transseries and Real Differential Algebra},
Lecture Notes in Mathematics 1888, Springer, 2006.

\bibitem{HoevenHyperseries}
L.~van den Dries, J.~van der Hoeven, and E.~Kaplan,
``Logarithmic hyperseries,''
\emph{Transactions of the American Mathematical Society} \textbf{372} (2019), 5199--5241.

\bibitem{Kuhlmann2000}
S.~Kuhlmann,
\emph{Ordered Exponential Fields},
Fields Institute Monographs 12, American Mathematical Society, 2000.

\bibitem{KuhlmannTressl2011}
S.~Kuhlmann and M.~Tressl,
``Comparison of exponential-logarithmic and logarithmic-exponential series,''
\emph{Mathematical Logic Quarterly} \textbf{58} (2012), 434--448.
\url{https://arxiv.org/abs/1112.4189}

\bibitem{Olver1997}
F.~W.~J.~Olver,
\emph{Asymptotics and Special Functions},
A K Peters, 1997; originally published 1974.

\bibitem{Schmeling2001}
M.~C.~Schmeling,
\emph{Corps de transs\'eries},
Ph.D. thesis, Universit\'e Paris~VII, 2001.

\bibitem{WikipediaTransseries}
Wikipedia contributors,
``Transseries,'' \emph{Wikipedia, The Free Encyclopedia},
accessed August 31, 2026.
\url{https://en.wikipedia.org/wiki/Transseries}

\end{thebibliography}

\end{document}
